\documentclass[11pt, twoside, openright]{book}

% Trim size: 6" x 9" (US Trade paperback)
\usepackage[
  paperwidth=6in, paperheight=9in, inner=0.875in, outer=0.75in, top=0.75in, bottom=0.75in
]{geometry}

% Typography
\usepackage{fontspec}
\usepackage{xeCJK}
\setmainfont{TeX Gyre Pagella}
\setCJKmainfont{Noto Serif CJK SC}
\emergencystretch=1em          % Allow extra stretch to avoid overflow with hyphenated/apostrophe words
\usepackage{setspace}
\setstretch{1.15}              % Standard book leading

% Graphics
\usepackage{graphicx}
\graphicspath{{../figures/}{../figures/book/}}

% Links (hidden for print - no colored text)
\usepackage{xcolor}
\usepackage[hidelinks, bookmarks=false]{hyperref}
\hypersetup{
  pdftitle={你称之为“我”的模拟},
  pdfauthor={Matthias Gruber},
  pdfsubject={Consciousness, Computation, and the Cosmos},
}

% Chapter styling
\usepackage{titlesec}
\titleformat{\chapter}[display]
  {\normalfont\huge\bfseries}
  {\chaptertitlename\ \thechapter}{20pt}{\Huge\raggedright}
\titleformat{name=\chapter,numberless}[display]
  {\normalfont\huge\bfseries}
  {}{0pt}{\Huge\raggedright}
\titlespacing*{\chapter}{0pt}{50pt}{40pt}

% Section styling
\titleformat{\section}
  {\normalfont\Large\bfseries}{}{0pt}{}
\titlespacing*{\section}{0pt}{20pt}{10pt}

% Headers and footers
\usepackage{fancyhdr}
\pagestyle{fancy}
\fancyhf{}
\fancyhead[LE]{\small\itshape 你称之为“我”的模拟}
\fancyhead[RO]{\small\itshape\leftmark}
\fancyfoot[C]{\thepage}
\renewcommand{\headrulewidth}{0.4pt}

% For plain pages (chapter starts, front matter)
\fancypagestyle{plain}{
  \fancyhf{}
  \fancyfoot[C]{\thepage}
  \renewcommand{\headrulewidth}{0pt}
}

% Quote formatting
\usepackage{csquotes}
\renewenvironment{quote}
  {\list{}{\leftmargin=1.5em\rightmargin=1.5em}\item\relax\itshape}
  {\endlist}

% Prevent widows and orphans
\widowpenalty=10000
\clubpenalty=10000

% Tables
\usepackage{booktabs}
\renewcommand{\lightrulewidth}{0.8pt}   % KDP minimum line weight is 0.75pt
\usepackage{longtable}
\usepackage{array}
\usepackage{tabularx}

% Figure placement
\usepackage{float}

% Landscape pages for wide tables
\usepackage{pdflscape}
\renewcommand{\textfraction}{0.1}
\renewcommand{\topfraction}{0.9}
\renewcommand{\bottomfraction}{0.9}

% Make blank pages (from \cleardoublepage) truly blank
\makeatletter
\def\cleardoublepage{\clearpage\if@twoside \ifodd\c@page\else
  \thispagestyle{empty}\hbox{}\newpage
  \if@twocolumn\hbox{}\newpage\fi\fi\fi}
\makeatother

\begin{document}

% ==== FRONT MATTER ====
\frontmatter
\pagestyle{empty}

% ---- Half-title page (recto) ----
\vspace*{3in}
\begin{center}
{\Huge\bfseries 你称之为“我”的模拟\par}
\end{center}
\cleardoublepage

% ---- Full title page (recto) ----
\vspace*{2in}
\begin{center}
{\Huge\bfseries 你称之为“我”的模拟\par}
\vspace{0.8cm}
{\Large 意识、计算与宇宙的架构\par}
\vspace{2cm}
{\large Matthias Gruber\par}
\end{center}
\clearpage

% ---- Copyright page (verso of title) ----
\thispagestyle{empty}
\vspace*{\fill}
{\small
\noindent \textcopyright\ 2026 Matthias Gruber。保留所有权利。\par
\vspace{0.5cm}
\noindent 未经作者事先书面许可，不得以任何形式或任何方式复制、分发或传播本出版物的
任何部分，但评论中的简短引用及版权法允许的某些非商业用途除外。\par
\vspace{0.5cm}

\vspace{0.5cm}
\noindent 第二版，2026年\par
\vspace{0.5cm}
\noindent www.matthiasgruber.com\par
}
\cleardoublepage

% ---- Dedication page (recto) ----
\thispagestyle{empty}
\vspace*{3in}
\begin{center}
\textit{献给每一位曾经好奇过的人——为什么任何事物，竟会让人感觉到些什么。}
\end{center}
\cleardoublepage

% ---- Table of contents ----
\pagestyle{plain}
\renewcommand{\contentsname}{目录}
\renewcommand{\figurename}{图}
\tableofcontents
\cleardoublepage








\chapter*{那个并不存在的圆点}
\addcontentsline{toc}{chapter}{那个并不存在的圆点}
\markboth{那个并不存在的圆点}{}

在我向你讲述任何东西之前------在讲理论之前，在介绍我是谁之前，在解释这一切为什么与你有关之前------我想请你先亲自向自己证明一件事。只需要十秒钟。

看看下面这两个标记：左边是一个\textbf{X}，右边是一个实心的\textbf{圆点}。


\begin{figure}[htbp]
  \centering
  \includegraphics[width=0.95\textwidth]{../figures/blind-spot-test.png}
\end{figure}


把书举到一臂远的地方。闭上\textbf{左}眼。用右眼直直地盯住那个X------只看X。别作弊，别偷瞄那个圆点；你能感觉到它就在你的余光里，这就够了。现在，让目光牢牢锁定X，慢慢把书移向你的脸。

在书离你的脸大约30厘米的时候，留意视野边缘的那个圆点。

它不见了。

不是变模糊，也不是变淡。是\textbf{消失了}------而这里才是值得停下来想一想的地方：它原来所在的位置上并没有一个洞。没有缺口，没有暗影，没有一块空白。就在纸上印着一个粗黑圆点的确切位置，页面却干干净净、天衣无缝、连绵一片地白着------就在你睁着的、功能完好的眼睛面前。你的眼睛正对着它。而你看不见它，因为你的大脑已经悄悄用白色把它涂掉了。

刚才发生的事情是这样的。每只眼睛都有一个盲点------视网膜上的一小片区域，正好位于视神经穿出、通向大脑的地方，那里完全没有感光细胞。你的视野里有一个洞，大得足以吞下九个满月，每只眼睛里都有，时时刻刻都在。而你一次也没有注意到它。不是因为它小------它大得惊人------而是因为你的大脑从不让你看见这个洞。它读取缺口周围的信息，然后把中间\textbf{编造}出来：编出纸页的白色，编出墙纸，编出任何“理应在那里”的东西，再把这幅完工的图景递给你，就好像那是未经加工的真相。

再把书移开。圆点回来了。多试几次。你正在亲眼看着自己的大脑在“如实报告世界”和“凭空发明世界”之间来回切换，而你感觉不到那个切换的瞬间。它从不预告自己。它从来没有过。

接下来这部分听起来像疯话，却恰恰是已经确立的神经科学：这种填补并不是大脑专门留给盲点的特殊把戏。它是你的大脑在整个视野范围内、在每一个清醒的瞬间都在做的事。盲点只不过是唯一笨拙到会露馅的地方------那一小块完全没有光线输入的区域，编造必须从零开始，于是失了手。在其他任何地方，原始信号几乎同样稀薄、同样嘈杂，由同一套猜测机制填补；你之所以毫无察觉，只是因为那些猜测通常都是对的。那个转瞬熄灭的圆点，是你的大脑被当场抓了个正着。你所看到的其余一切，是同一场表演------只是演得太好，好到无从察觉。

顺着这条线索走下去。你这一生，从来没有直接看见过现实。一次也没有。你手中的这一页不是，你周围的房间不是，你所爱过的每一个人的面孔也不是。它们无一例外，都是重建------是一个模型，在你的颅骨之内，用涓涓细流般的神经冲动和堆积如山的先验猜测搭建而成，然后如此平滑地呈现给你，以至于“这也许不是真实的东西”这个念头，直到此刻之前，多半从未在你脑海里闪现过。

你正生活在一个模拟之中。它由大约一公斤半的湿润组织在黑暗中生成，它是你此生唯一能够进入的世界，而在运行它的一辈子里，它一次也没有提起过自己的存在。

你此刻也许正感到的那阵眩晕------抓住它，别放手。那种感觉，正是这本书的全部主题。

我从内部体会过它。二十五岁那年，在因斯布鲁克的一座桥上，光天化日之下，我泪流满面，同时控制不住地大笑，我一生中最沉重的一块石头从身上落了下来。我不确定有没有人看见我。就算有，我也不会在乎。刚刚一下子全部落定的，是一个框架：它不仅解释了大脑如何构建那个天衣无缝的世界，还解释了这种构建为什么\textbf{会有任何感觉}------为什么“作为你”这件事本身有它的滋味，就在你读到这里的此刻。从那一刻起，在我看来，我余生的待办清单已经写完了。

2015年，我把它写了下来：一本300页的书，用德语写成，自费出版，通篇密布着技术性的论证装置。它卖出了零本。一本都没有。我说这个不是为了博取同情------我说它，是因为它与这个故事有关。那本无人翻阅的书里的理论，消解了或许是科学中最难的悬而未决的问题，做出了任何竞争理论都给不出的预测，还交出了一份建造真正有意识的机器的具体蓝图------不是一个模仿心智的聊天机器人，而是一台\textbf{拥有}心智的机器。然后，它在一块硬盘上躺了十年，与此同时，证据悄悄地一件件站到了它这一边。这本书是第二次尝试：更短，用英语写成，写给每一个曾经好奇“为什么任何东西会有任何感觉”的人。如果我是对的，意识的“难问题”其实并不难。它是一个范畴错误，在你看清它的那一瞬间就会消解------就像一页之前，那个圆点消失的方式一样。


\chapter*{关于作者}
\addcontentsline{toc}{chapter}{关于作者}
\markboth{关于作者}{}

我应该先告诉你，提出这些主张的人是谁。我不隶属于任何大学，没有博士学位------只有一个生物信息学硕士。我从未拿过科研经费，也从未在神经科学实验室里工作过。如果你在相信一个论证之前习惯先查验资历------这是个合理的本能------那么读到这里，你也许就要把书放下了。不妨拿它去垫垫那张摇晃的桌子，好歹别让那些树白白牺牲。

我拥有的，是一段自学局外人的经历：几十年的阅读，每次撞上一堵墙就转个方向，而且从未学会尊重学科之间的边界。这一点远比听起来重要。这套理论之所以存在，只是因为从来没有人告诉我，哪些学科是不许彼此接触的------数学与元胞自动机、机器学习、从临床神经病学到精神药理学的神经科学、演化生物学、心灵哲学。大多数意识理论只活在其中一两个世界里，而这一套试图把它们全部绑在一起。仔细想想，这恰恰就是大脑每天在做的事：它把来自完全不同来源的各路信息流，编织成一个单一的体验。一套不能像大脑那样跨越学科的意识理论，本身就该让你起疑。我的理论正是它所描述的那个递归循环的产物------知识推动表现，表现激发动机，动机又反哺知识------这个循环我会在书的后面细细展开。至于这套理论到底好不好，由你自己来判断。现在，进入理论本身。


\mainmatter
\pagestyle{fancy}
\chapter*{第1章：科学中最难的问题}
\addcontentsline{toc}{chapter}{第1章：科学中最难的问题}
\markboth{第1章：科学中最难的问题}{}

你正在读这句话。你正在经历一段体验。

那段体验（纸页上字母的视觉印象、诵读文字的内心声音、理解或困惑的感觉）既是你生命中最熟悉的东西，也是宇宙中最神秘的东西。我们对黑洞内部的了解，都要多于我们对“阅读为何会有某种感觉”的了解。

这并非夸张。（不过公平起见我得说明，数学里有些问题我认为还要更难，只是它们不会让大多数人夜不能寐。）物理学家有标准模型，生物学家有演化与遗传学，化学家有元素周期表。可是意识（此刻正在读这段话的你，“成为你”这件事确有“某种感觉”）却没有一套确立的理论，没有一个占主导的框架，没有一个公认的解释。

这并不是因为没人尝试。自20世纪90年代起，意识在经历数十年行为主义的放逐之后重新成为一个体面的科学话题，此后已有数千篇论文发表、数十种理论被提出、数亿美元被投入。结果呢？这个领域正处于科学哲学家托马斯·库恩（Thomas Kuhn）所说的“前范式状态”------众多相互竞争的观点，没有共识，还有一种日渐加深的感觉：某样根本性的东西也许一直缺席。

\section*{“难问题”究竟在问什么}

1995年，哲学家戴维·查默斯（David Chalmers）为这个谜题起了它经典的名字：意识的“难问题”。

它问的是这样一件事。想想看见红色的体验。当你看见红色时，神经科学家能告诉你大脑里发生的许多事：某个波长的光照到你视网膜上的视锥细胞，信号沿视神经传导，在视觉皮层被加工，各个脑区协同产生这一知觉。所有这些都已得到很好的理解，至少在轮廓上如此。

但这一切都没有解释\textbf{看见红色为什么会有某种感觉}。

原则上，你可以为大脑对红光的反应建立一个完整的神经模型------每一个神经元、每一个突触、每一条信号通路。你会得到一份完美的功能性描述。而你仍然没有解释红的那种感觉。那个“感觉起来是怎样的”。哲学家所说的\textbf{感受质}（qualia）。

查默斯把这一问题与意识的“容易问题”区分开来（后者其实一点都不容易，只是原则上可解）：大脑如何整合信息？它如何调配注意力？它如何报告自身的状态？这些是关于机制的问题。它们很难，但属于神经科学懂得如何着手的那种难。难问题则不同：它问的是，为什么这些机制竟然伴随着体验。为什么大脑不像一台计算机那样，只是“在黑暗中”加工信息？

\section*{当前的局势}

倒不是缺少竞争者。\textbf{整合信息理论}把意识归结为一个数字，$\Phi$------并为这份优雅付出了代价：它意味着一个足够大的逻辑门阵列也带着一点点意识；2023年，超过120位研究者联名发表公开信，称它不可证伪。\textbf{全局神经元工作空间理论}说，当信息被全脑广播时它就成为有意识的------这是对\textbf{何时}的清晰交代，却对\textbf{为何}广播应当伴随任何感觉保持着意味深长的沉默。\textbf{预测加工}把大脑刻画成一台预测引擎，把知觉刻画成一场“受控的幻觉”------在体验的\textbf{内容}上精彩非凡，而据其自身的构建者承认，它甚至不曾瞄准那个难的部分。高阶理论、注意图式、循环加工、电磁场理论------每一个都是真切的洞见，包裹着同一个缺失的核心。

随后，在2025年，裁判走上了赛场。COGITATE合作项目对两位领跑者------IIT对阵GNW------做了一场持续多年的对抗性检验，并把结果发表在《自然》（\textit{Nature}）上。裁定是：两者都没赢。与意识关联最强的信号出现在大脑后部，而那里两个阵营都没有插旗。历经三十年、数亿美元之后，这个领域可以说比起步时离共识更远了。这不是一门科学正在逼近答案的画面。这是一门科学缺了一块的画面。

\section*{阻碍进展的两条教条}

在我告诉你我认为缺了什么之前，我得先给两种偏见起个名字------几十年来，它们一直在悄悄地破坏这个领域。我在自己最初那本书里就给它们起了名字，因为我觉得没有名字的偏见更难对付。

第一个我称之为\textbf{nSAI教条}------“没有强人工智能”。这是一种广为流传的坚信：真正有智能的机器是不可能的；这份坚信并非扎根于证明，而是扎根于20世纪60年代早期人工智能研究的失败以及随之而来的反弹。凡是相信强AI可能的人，如果想在主流研究里被认真对待，就学会了闭口不谈。这不是理性的怀疑。这是旧日败绩留下的伤疤，硬化成了教义。

第二个更深，也更有害。我称之为\textbf{nSU教条}------“没有自我理解”。它是这样一种信念：人类的心智、人类的意识，原则上无法被这同一个心智所理解。人们搬出哥德尔不完备定理，或者搬出与“从宇宙内部观测宇宙学”之局限的模糊类比，又或者（最诚实地）他们只是觉得被彻底解释这一前景太可怕，不忍去想。如果意识不过是一台机器，那么灵魂会怎样？意义会怎样？身而为人的那份特别会怎样？

这两条教条彼此加固。如果你无法理解意识（nSU），那你当然也造不出一个来（nSAI）。而如果你造不出一个来（nSAI），那么意识也许真的超乎理解之外（nSU）。这是一个封闭的、体制性悲观的循环，它挡住了大批聪明的研究者，让他们连尝试这项工作都没有尝试。

我并不是说这些教条是出于恶意才被抱持的。许多研究者是真诚地相信它们的。但两条教条都从未被证明过。它们是信条，而它们给意识研究造成的损害，比任何一次失败的实验都要大。

\section*{缺了某样东西}

我认为，之所以没有哪种理论攻克了难问题，是因为大多数人一直在错误的地方寻找意识。他们盯着神经机器（那些神经元、突触、振荡、连接），然后问：“这些过程里哪一个是有意识的？”

我相信正确的问题是另一个：“在信息加工的哪一个层级上、动用哪一种架构，体验才会发生？”

这就是四模型理论（FMT）的出发点------一个在这本书还没真正开始之前你就已经向自己证明了的事实。你这一生，从未直接体验过现实。你从纸页上擦掉的那个点，并不是纸张玩的把戏；那是你的大脑在用编造的内容填补视网膜上的一个洞------就像它填补每一个洞那样，在每一次目光扫过之间，天衣无缝到你从未起过丝毫疑心。这并不是什么边缘论调：知觉是被建构出来的、而非直接给予的，这是主流神经科学，几乎该领域的每一位研究者都接受。四模型理论所补上的，是几乎无人愿意迈出的那一步------这唯一一个事实，一路追究到底，会把难问题化解掉。

\section*{一把剃刀}

我是在一条唯一可见的纪律之下建起这套理论的：\textbf{奥卡姆剃刀。}没有新物理，没有微管里的量子魔法，没有洒遍物质的原初意识------只有我们已经拥有的原料（神经网络、学习、模拟、自我指涉），以正确的架构排布起来。另有两条古老的原则在幕后悄悄发挥着作用，与其现在就让它们从你面前列队走过，我宁愿把每一条都恰好留到它真正管用、真正咬人的地方再登场：一条留到我们谈动物的时候，另一条留到我们谈僵尸与自由意志的时候。眼下，一把刀就够了。

现在，是时候来看看这四个模型了。


\chapter*{第2章：四个模型}
\addcontentsline{toc}{chapter}{第2章：四个模型}
\markboth{第2章：四个模型}{}

想象你正在看一个苹果。

苹果就摆在你面前的桌子上。红润、圆实、发亮，离你的手大约十五厘米远。你能看见它，你知道那是什么，你伸手就能把它抓过来。这似乎再简单不过------你在看一个苹果。

但实际发生的事情，要复杂得多得多。

苹果表面反射的光进入你的眼睛，落在视网膜上的光感受器细胞上。这些细胞把光转换成电信号。信号沿着视神经一路传到大脑后部的视觉皮层，在那里经过一层层越来越精细的特征探测器处理：边缘、方向、颜色、纹理、形状，最终是物体。在这一连串处理的某个环节里，对应“苹果”的神经活动被激活了。与此同时，你的运动系统正在准备可能的动作（伸手、抓握），你的记忆系统正在唤起联想（味道、口感、上一次你吃苹果的情景），你的空间系统正在追踪苹果相对于你身体的位置。

这一切发生在不到一秒的时间里。而这些都不是你所\textbf{体验}到的东西。你并不体验光子撞上视锥细胞，也不体验信号沿着轴突传递，或是特征探测器在放电。你体验到的是\textbf{一个苹果}。一个统一、稳定、三维的物体，安放在一个连贯的空间环境里，有着特定的模样、质感和意义。你所体验到的是一个\textbf{模型}------一场关于苹果的实时模拟，由你的大脑根据原始数据，以及它此前学到的关于苹果、物体、桌子和物理的一切生成出来。

正如我在第1章所论证的，这在神经科学上毫无争议。每一位神经科学家和知觉哲学家都同意：你体验到的是一个模型，而不是现实本身。你看到的那个苹果，是大脑对那里究竟有什么所做出的\textbf{最佳猜测}------它以感官数据为依据，却又与之不完全相同。（视错觉就是活生生的证据：当一个错觉崩塌时------当你忽然能同时看出它的两种样子时------你就当场逮住了正在运作的模拟。你从来没有直接看到过现实。你看到的一直是模型。错觉只是把这一点摆到了明处。）

而我的理论正是从这里开始：大脑不只是给苹果建模。它还给\textbf{正在看苹果的你}建模。正是这第二个模型（关于自我的模型），把信息加工变成了意识。

\section*{你大脑的四种表征}


\begin{figure}[htbp]
  \centering
  \includegraphics[width=0.95\textwidth]{../figures/figure2-real-virtual-split-bw-zh.png}
  \caption{真实／虚拟之分。基底（真实一侧）把知识储存在突触权重里------它是物理的、结构性的、无意识的。模拟（虚拟一侧）实时生成体验------它是短暂的、动态的、有意识的。}
\end{figure}


哪怕是只有三层的简单神经网络，也能学会对自己的输入建模------给它看足够多的例子，它就会为所遇到的模式建立起内部表征。你的大脑做的正是这件事，只不过规模要大得多、丰富得多。它不是建一个模型，而是建许许多多个，涵盖从视野到四肢位置、从一段嗓音到双脚踩在地面上的压力等等一切。这些模型既涵盖你之外的世界，\textbf{也}跨越你自身的身体，把所有可用的感官输入连结成连贯的表征。

神经科学认识这些模型已经有一个多世纪了。在运动皮层和体感皮层里，身体被实实在在地铺展成一幅扭曲的地图------手和嘴唇被夸张地放大，因为它们的神经末梢更多，躯干和腿则被压缩成细条。这些皮层地图被称为\textbf{小人图}（homunculi），最早由威尔德·彭菲尔德（Wilder Penfield）在20世纪30年代通过脑外科手术中的直接电刺激绘制出来。它们只是最鲜明的例子；大脑在其整个架构中维持着类似的地图和模型。（关于皮层组织的更多内容，见附录A。）


\begin{figure}[htbp]
  \centering
  \includegraphics[width=0.95\textwidth]{../figures/book/homunculi.en.png}
  \caption{彭菲尔德的皮层小人图。体感皮层分配给手、嘴唇和舌头的面积，远远超过整个躯干------这是一幅扭曲的身体地图，反映的是神经末梢的密度，而不是身体部位的大小。}
\end{figure}


我把这些叫作\textbf{隐式模型}：隐式世界模型（IWM）和隐式自我模型（ISM）。它们储存在大脑的结构之中------在突触连接的强度里，在神经回路的架构里，在一辈子积累下来的学习里。它们是大脑的硬盘。你从来不会直接体验到它们，就像你不会体验到手机里的硅片一样。可它们编码着大脑所知道的关于世界、关于你的一切。

现在，来看关键的洞见。这些隐式模型并不是干坐在那里。它们\textbf{生成}出某种东西。在工程学里，\textbf{数字孪生}是一个物理系统（一台喷气发动机、一张电网、一片工厂车间）的实时虚拟副本，它不断被传感器数据更新，好让工程师能够监测这个系统、与它互动，而无需直接触碰它。你的隐式模型做的正是这件事。它们生成出一场关于世界的实时虚拟模拟，以及一场关于你的实时虚拟模拟。这些就是\textbf{显式模型}：显式世界模型（EWM）和显式自我模型（ESM）。你所看、所听、所感、所想的一切，都发生在这些模拟之内，而不是发生在世界本身。

两组模型（隐式与显式），每一组都同时包含一个世界模型和一个自我模型。总共四个模型------有了它们，我们就有了一套语言，来谈论意识究竟在做什么。（给神经科学家和技术型读者的一点说明：数字“四”是一个原则性的最小值，而不是对大脑实际维持之物的字面清点。如果这让你介意，请在继续之前读一读附录E------它直接回应了这个问题。）

那么，现在来看这四个模型。

\textbf{隐式世界模型（IWM）}是你关于世界所知道的一切。不是你此刻正在想的东西------而是你\textbf{能够}去想的一切。物理定律（你知道松手的东西会落下）。你公寓的格局（你能在黑暗中在里面走动）。你母语的语法（你能判断一句话是否合乎语法，却不必知道规则）。你这辈子认识过的每个人的脸。巧克力的味道。雨声。

所有这些知识都储存在你大脑的突触连接里------储存在神经元之间那些连结的强度里。它是在你整整一生中，通过经验和学习一点点建立起来的。而你永远、永远都不会直接觉察到它。你无法对自己的神经连接做内省。你感受不到自己的突触。隐式世界模型就像一座你从不踏入的浩瀚图书馆------你只是读它送到你书桌上的那些书。

\textbf{隐式自我模型（ISM）}是你关于自己所知道的一切。你的身体图式------对四肢在哪里、有多大、如何运动的无意识表征。你的运动技能（骑车、打字、演奏乐器）。你的性格特质、社交本领、情绪模式、习惯。你的自传体记忆结构（把你的种种记忆组织成一段人生故事的那套框架）。

和世界模型一样，自我模型储存在突触权重里，从不直接进入意识。你体验到的不是你的身体图式；你体验到的是这套图式所生成的身体。你体验到的不是你的性格；你体验到的是你的性格所产出的种种念头和感受。隐式自我模型是幕后的班底------对演出不可或缺，却从不被观众看到。

\textbf{显式世界模型（EWM）}是你真正体验到的那个世界。就在此刻。你所在的房间，你听到的声音，这本书在你手里的重量（或者你正在其上阅读的屏幕的辉光）。这就是那场模拟------大脑的实时虚拟现实，由隐式世界模型加上当下的感官输入生成而来。它生动、细致，天衣无缝地令人信服。你会在它里面度过你的一生，从不曾迈出一步。

\textbf{显式自我模型（ESM）}就是\textbf{你}。那种身为一个主体的感觉。那个“我”的感觉------那个在看、在听、在想、在做决定的人。这，同样是一场模拟：一个由隐式自我模型加上当下的身体信号生成而来的实时模型。它是大脑创造出来、用以栖居在自己虚拟世界里的那个角色。

你就是大脑创造出来、用以栖居在自己虚拟世界里的那个角色。

\section*{真实的一侧与虚拟的一侧}


\begin{figure}[htbp]
  \centering
  \includegraphics[width=0.95\textwidth]{../figures/figure1-four-model-architecture-bw-zh.png}
  \caption{四模型架构。大脑维持两类模型（一类关于世界，一类关于自我），每一类都有两种模式：隐式（存储在大脑的结构之中）与显式（作为实时模拟主动运行）。意识栖居在显式模型之中。}
\end{figure}


这四个模型分成两侧，而这一划分是接下来一切内容的基础。

\textbf{真实的一侧}（两个隐式模型）是物理的、结构性的，且相对刚硬（由学习缓慢地塑造）。它是大脑存储的知识：突触权重、网络连接、受体的配置。你可以把它想成大脑\textbf{已经学会}的一切------凝结进脑组织本身的物理结构之中。它没有任何体验。一个突触放电，就像水流过管道一样，谈不上什么“体验”。真实的一侧，灯是关着的。

有一点值得强调：真实的一侧正是神经科学一直在研究的东西。当一位研究者把你送进 fMRI 扫描仪时，他看的就是真实的一侧------放电模式、连接方式、流向不同脑区的血流。当一位神经外科医生刺激某个皮层区域、观察随之发生的事情时，他探测的也是真实的一侧。神经科学绘制这片疆域已经一百多年了，并取得了非凡的进展。四模型理论并不否定这些工作中的任何一项。它只是说，所有这一切描述的，仅仅是画面的一半。

\textbf{虚拟的一侧}（两个显式模型）是被模拟出来的、转瞬即逝的、动态的。它在每一刻都由真实的一侧加上当下的输入重新生成。你可以把它想成大脑此刻\textbf{正在拿}它学到的东西\textbf{做什么}------是现场演出，而不是存档的剧本。而它就是体验的\textbf{全部}。你曾有过的每一次看见、每一声听闻、每一个念头、每一份感受、每一段记忆、每一场梦、每一次幻觉，都发生在虚拟的一侧之内。虚拟的一侧，灯是亮着的。

但难点在这里：虚拟的一侧从外部是看不见的。哪怕我们最先进的脑成像，也只能间接地捕捉它。一台 fMRI 显示的是哪些脑区处于活跃状态------那是真实的一侧在工作。要真正从脑数据里\textbf{读出}意识体验，你得破译大脑的编程语言------不仅要理解哪些神经元在放电，还要理解那放电的模式在模拟层面\textbf{意味着}什么。那需要类似于一个完整模拟的连接组：大脑的一份完整数字复本，在软件中运行，生成与生物大脑相同的那个虚拟世界。

关于这套理论给你什么、不给你什么，我想说实话。四模型理论告诉你这场模拟\textbf{是什么}、它栖居在\textbf{哪里}、以及它\textbf{为什么}让你感觉到某种东西。它并不把那枚解码指环递到你手里。从真实的一侧读出虚拟的一侧，是一项未来的研究计划------理论清晰地界定了它，但眼下还无法兑现。不过，这项计划的基础已经在铺设之中。人类连接组计划及相关工作，正以越来越精细的分辨率绘制大脑的接线图。我们还无法从结构数据中解码出虚拟的一侧，但结构数据正在到来。

如果你有科学的头脑，也许已经看出这将通向何方。如果体验只存在于虚拟的一侧，那么到真实的一侧去寻找体验------在神经元里、在突触里、在物理机器里寻找------就完全找错了地方。这就好比在 DVD 播放机的电路里搜寻一部电影的情节。

这就是关键所在。

\section*{你的大脑为何具备自我建模的能力}

那么我们已经确立：意识依赖于这四个模型，其中显式自我模型承担着最吃重的部分。但人类大脑一开始为什么会拥有这种能力，而更简单的动物却没有？答案就明摆在眼前：人类皮层的架构，说得直白些，对于基本的信息处理而言是超规格的。

人类的新皮层有六层。这是一个众所周知的解剖学事实。你在任何一本神经生物学教科书里都能看到它。但有意思的地方在这里：处理信息并不需要六层。三层就够了。

想一想一个标准的神经网络需要做什么。第一层：接收输入，过滤它，把它清理干净。第二层：提取模式，识别特征，完成繁重的计算活儿。第三层：整合结果，做出决策，产生输出。输入、处理、输出。这就是基本配方，而三层就能覆盖它。

但我们有六层。

那“多出来”的三层是干什么的？

它们是用来给头三层建模的。

一个三层网络处理世界。一个六层网络处理世界，\textbf{而且}观察自己在处理。多出来的这些层，为大脑提供了架构上的能力------不仅能构建一个关于外界的模型，还能构建一个关于自己正在为外界建模这件事的模型。自我模拟需要这种加倍------你需要一套层去做处理，另一套层去看着处理的发生。

这并不是在猜测各层单独“做”什么------我没有声称第4层做这个、第5层做那个。这是一项关于架构能力的观察。六层给你腾出了空间，既容得下隐式世界模型（IWM）（那已习得的、无意识的处理），也容得下显式世界模型（EWM）（实时模拟）。它们给你腾出了空间，既容得下隐式自我模型（ISM）（你的身体图式、运动程序、人格结构），也容得下显式自我模型（ESM）（那个体验到拥有身体、发起动作、身为一个人的“你”）。

现在来看看其他动物。爬行动物有三层皮层。哺乳动物有六层。而在哺乳动物当中，皮层最厚、褶皱最精巧繁复的那些（灵长类、鲸类、大象），恰恰就是展现出最丰富自我觉知迹象的那些。镜中自我识别、对未来的规划、社会性欺骗、哀伤。架构上的能力，与主观体验相互对应。

从三层到六层的这一跃迁，或许曾是一场遗传上的复制事故------演化的复制粘贴，恰好造出了意识日后将会利用的那套架构。爬行动物的祖先有三层皮层。在向哺乳动物过渡的某处，这个数字翻了一倍。那些指定皮层各层身份的转录因子（Tbr1、Satb2、Ctip2、Fezf2）都有旁系同源基因，暗示着基因复制事件。这究竟是一场单一的戏剧性事件，还是一步步的逐渐精化，至今仍有争议，但结果是清楚的：哺乳动物得到了双倍的层数，也随之得到了大多数爬行动物所缺乏的自我建模能力。

这就是从神经网络理论通往切身体验的桥梁。人类的皮层不只是一台庞大的模式识别器。它是一个超规格的、递归构造的网络，拥有足够多的层数，去为它自己的建模过程建模。而当一个网络为“自己如何为世界建模”这件事本身建模时，其结果------从内部看去------恰恰就是我们所称的意识。

我得说清楚：我并不是在声称六层皮层是\textbf{唯一}能够支撑意识的架构。它们是一种解决方案------哺乳动物演化出的那一种。但也可能有别的。章鱼，凭着它那彻底分布式的神经系统（八条半自主的腕足，每一条都有自己的神经处理中枢，约含4000万个神经元），代表着一种全然不同的架构路径，或许能达到相当的算力水平。鸟类提供了另一个引人注目的例子：鸦科鸟类和鹦鹉根本没有分层的皮层，它们的外套区组织成核团簇而非片层，可乌鸦却会制作工具、为未来做规划，并且有理由认为它们能在镜子里认出自己。如果关键在于自我建模的能力，而不在于具体的接线图，那么任何能够运行一份自身模拟的架构，原则上都可能有意识。章鱼、乌鸦，或者某种我们尚未造出来的东西。

不过其他那些架构可以先等一等。有一个更贴近的问题，正是它，在我第一次把这一路想到底时让我停了下来。如果\textbf{你}就是大脑写进它自己那场模拟里的那个角色------那么，是谁在看这场演出呢？


\chapter*{第3章：虚拟的那一面}
\addcontentsline{toc}{chapter}{第3章：虚拟的那一面}
\markboth{第3章：虚拟的那一面}{}

想象你正在玩一款电子游戏。一款好游戏------沉浸式的开放世界，画面惊艳，物理效果真实，故事引人入胜。你操控着一个角色，并通过这个角色，与一个细节丰富的虚拟世界互动。

现在想一想：这款游戏究竟存在于何处？严格说来，不在屏幕上------屏幕只是显示光的图案。严格说来，也不在显卡或CPU里------它们只是让电信号在硅电路中流动。游戏是作为一种\textbf{虚拟进程}而存在的------它是一种更高层级的现象，从硬件的活动中涌现出来，却不等同于任何一块具体的硬件。

游戏的虚拟世界拥有硬件所没有的属性。游戏里有高山、河流和城市，CPU里只有晶体管。游戏里有昼夜交替，GPU里只有时钟周期。你可以有意义地问“游戏里那座山有多高？”，但要是指着一个晶体管说“这个晶体管有3000米高”，那就荒谬了。游戏的属性存在于虚拟层级，它们是游戏真实的属性------尽管这款游戏“不过”是硬件里的一种活动图案。

这不是比喻。你的大脑就是这样运作的。

你的显式世界模型（EWM）------你所体验到的那个世界------是一个运行在神经硬件上的虚拟进程，就像游戏世界是一个运行在硅硬件上的虚拟进程一样。你体验到的世界拥有神经硬件所没有的属性（颜色、形状、距离、声音），而硬件拥有的是放电频率、突触强度和神经递质浓度。你所体验的世界的那些属性，是\textbf{模拟真实的属性}------尽管这场模拟“不过”是一种神经活动的图案。

你的显式自我模型（ESM）------那个体验着世界的“你”------同样是一个虚拟进程。它和这个类比里的游戏角色一样真实：在虚拟层级上确实存在，在虚拟层级上确实拥有属性，但在硬件层级上并不存在。

\section*{这个类比为何会失效（而且是在最要紧的地方）}

电子游戏的类比很有用，但它会在一个关键点上失效：游戏有一个\textbf{玩家}。有一个人在游戏之外------坐在沙发上的你------体验着这款游戏。游戏本身没有任何体验。它只是光的图案和代码。

大脑的这场模拟没有外部玩家。没有谁坐在你的头颅之外体验这场模拟。模拟包含着它自己的观察者------显式自我模型。模拟\textbf{就是}体验本身，而不是被别人体验的某样东西。

把自己代入游戏角色的位置。你\textbf{就是}主角。从游戏之外看，旁观者看到的是像素在屏幕上移动------没有任何东西可能感受到什么。但从模拟内部看呢？游戏世界就是一切。对角色来说，群山是真实的，阳光是温暖的，危险是可怕的。没有哪个外部观察者会猜到这堆代码能感受到什么，但那只是因为他们看错了层级。他们盯着的是硬件。而体验存在于软件层级。这就是我的主张，本书余下的部分将把证据一一摆出来。

这正是意识的特别之处，也正是意识的“难问题”显得如此棘手的原因。在电子游戏里，游戏（虚拟、无体验）和玩家（物理、有体验）之间有一道清晰的分界。而在大脑里，没有这道分界。模拟与体验者是同一样东西。显式自我模型并非从外部观看显式世界模型------它就\textbf{在}模拟内部，是同一个虚拟进程的一部分。

而这种自指闭合（模拟从内部观察自身），我认为，正是我们所说的意识。它不是附加到模拟之上的某样东西。当模拟包含了一个关于自身的模型时，它\textbf{就是}模拟本身。所以我才说，意识不是一样东西------它是一个过程。你无法通过把大脑拆开来找到它，正如你无法通过拆解CPU来找到一个正在运行的程序。

\section*{那些软件属性}

如果这些虚拟模型真的是运行在神经硬件上的类软件进程，那么它们理应以某些具体的、可检验的方式表现得像软件。而它们确实如此。虚拟那一面的四个属性会在本书中反复出现，所以让我现在就把它们摆出来。

\textbf{分叉（Forking）。} 单一基质可以同时运行多套虚拟配置。在软件里，你分叉一个进程，就会得到两个独立的实例运行在同一硬件上。在大脑里，这就是分离性身份障碍------多个自我模型，各有各的叙事和情绪特征，轮流掌控同一个神经基质。我们会在第9章看到这一点。

\textbf{克隆（Cloning）。} 把硬件在物理上分开，你就会得到软件的降级但完整的副本。切断胼胝体，每个半球都会运行它自己那一版模拟------不如原来那么强大，但功能上是完整的。这就是裂脑现象，同样在第9章。

\textbf{改道（Redirecting）。} 打断正常的输入流，模拟就会锁定当下占主导地位的那个信号。在鼠尾草（Salvia divinorum）的作用下，本体感觉输入压倒整个系统，显式自我模型便围绕身体感觉重新配置。在氯胺酮的作用下，外部输入退场，模拟便依靠内部噪声运行。虚拟模型并不会停下------喂给它们什么，它们就处理什么。第6章会详细讲这一点。

\textbf{重构（Reconfiguring）。} 修改基质的连接权重，你就会改变虚拟模型所产出的东西。认知行为疗法所做的正是如此------系统性地为基质重新布线，好让显式自我模型生成不同的叙事、不同的情绪反应、不同的行为。

四模型理论（FMT）对治疗给出了一个具体的预测：任何有效的治疗，都必须通过修改隐式模型（基质），从而让显式模型（模拟）相应地发生改变，才起作用。认知行为疗法做的正是这件事。它系统性地识别出ISM中那些适应不良的模式，通过有结构的练习为它们重新布线，从而改变ESM所产出的东西。这就是为什么认知行为疗法拥有一切心理治疗中最强的证据基础：它瞄准了正确的层级。

这就引出了一个令人不安的问题------对那些说不清自身机制的疗法而言尤其如此。如果一种治疗方法说不清它在基质里改变了什么，或者这种改变如何传导到模拟，那么它顶多是在通过一个自己并不理解的机制起作用，而最坏的情况是它根本没起作用。证据也印证了这一点：证据基础最薄弱的那些疗法，通常也正是改变理论最含糊的那些。如果你正在寻求治疗，不妨问你的治疗师一个简单的问题：“你具体想在我的大脑里改变什么，又是怎么改变的？”如果他们答不上来，就考虑另找一位答得上来的吧。

这些不是比喻。它们是结构性的预测。如果我的理论错了，虚拟模型\textbf{并不是}类软件的进程，那么这些相似之处就纯属巧合。但巧合通常不会横跨临床神经病学、精神药理学和心理治疗，四次全部对上号。接下来的各章会展示每一个属性在实际中如何运作。

有一个简单的实验，你现在就能做------好吧，需要一位朋友、一只橡胶手、一块硬纸板挡屏和两支画笔------它能展示显式自我模型有多容易被骗过。这就是橡胶手错觉，由Matthew Botvinick和Jonathan Cohen设计，是整个神经科学里最能揭示问题的派对小把戏之一。

装置很简单。你坐在桌前，一只手臂藏在一块硬纸板挡屏后面。一只逼真的橡胶手被放在你面前、看得见的地方，大致就在你那只藏起来的手本该在的位置。有人用两支画笔同时抚过橡胶手和你藏着的真手------同样的位置，同样的速度。经过一两分钟这样同步的抚摸，某种诡异的事情发生了：你开始\textbf{感觉到}画笔的抚摸落在了橡胶手上。不是落在挡屏后面你真正的手上。是落在你眼前那只假手上。

你的显式自我模型已经把橡胶手纳入了它的身体图式。它重新分配了归属------判定那只橡胶手是“你”的一部分。自我模型不是天生写死的。它是习得的。它会依据当下最可靠的证据持续更新，而当视觉证据（看见橡胶手被抚摸）与触觉证据（感觉到你的真手被抚摸）始终吻合时，ESM便得出了理性的结论：那只手是我的。这时如果有人威胁那只橡胶手（拿锤子朝它砸下去），你会缩一下，会感到一阵焦虑猛地窜起，你的皮肤电反应会飙升。对于那部分定义了“你”的大脑而言，那只手\textbf{就是}你自己的。

这不是故障。这恰恰是自我模型在按设计运作------不断依据多模态的感觉相关性来更新它的身体边界。正是同一个机制，让截肢者在使用一段时间后能把假肢“感觉”成自己的。也正是同一个机制，在身体失认症中崩溃了------患者否认自己真实肢体的归属；在异手综合征中也是如此------那只手会自行移动。

\section*{拼缀而成的全息图}

虚拟一侧还有第五个属性值得单独辟出一节来讲，因为它解释了一件让神经科学家困惑了将近一个世纪的事：为什么脑损伤会让功能\textbf{逐渐}退化，而不是删除某些特定的记忆。

在20世纪20年代和30年代，心理学家卡尔·拉什利（Karl Lashley）训练大鼠走迷宫，然后通过手术切除它们皮层的一块块组织，想找出记忆究竟存在哪里。他从来没找到。无论他切掉哪一块，大鼠依然记得那座迷宫。真正起作用的是他切除了\textbf{多少}皮层，而不是切掉了\textbf{哪些部位}。切掉一点，大鼠表现略差；切掉一大片，它们就差得多。但记忆从来不曾干干净净地\textbf{消失}，像从硬盘上删掉一个文件那样被利落地切除。拉什利穷尽整个职业生涯寻找“印痕”（engram）------记忆的物理痕迹------最终得出了那个著名的结论：它似乎并不存在。

他找错了对象。记忆并不是像文件存在硬盘某个特定扇区那样，被存\textbf{在}某一块特定的皮层里。它是\textbf{分布}在整个网络之上的，散布在数百万神经元之间的连接权重里。神经网络就是这样运作的：信息并不待在任何单一的节点上，而是被编码在所有节点之间的连接模式里。你没法指着某一个突触说“迷宫就存在这儿”，就像你没法指着某一个像素说“电影就存在这儿”一样。

这本质上是一种全息属性。如果你拿一张实体全息图把它对半切开，你得到的不是图像的两半，而是\textbf{完整}图像的两份拷贝，只是分辨率各降低了一些。切成四份，你就得到四幅完整的图像，只是更加模糊。全息图里的信息分布在整块底片上，所以每一块都含有完整的画面------只是细节更少。

神经网络也是这么做的。训练一个网络去识别人脸，然后随机摧毁它10\%的连接。它不会忘掉10\%的面孔，而是对\textbf{所有}面孔都变得略差一些。摧毁50\%，它对一切都会差得多，但仍然能认出点什么。信息弥散在整个网络之中，这正是拉什利找不到印痕的原因：它无处不在，又无处可寻。

不过------有意思的地方就在这里------大脑并不是\textbf{一张}全息图。它是我所说的\textbf{拼布全息图}。在单个功能区内部（比方说你的初级视觉皮层，大致相当于布罗德曼17区），各个皮层柱彼此相似，信息以全息方式存储。摧毁其中几根皮层柱，你几乎察觉不到。这块区域在局部是全息的（一个部分包含整体，只是分辨率更低）。

而在全局层面，不同区域做着不同的事。你的视觉皮层和运动皮层不能互换。把整个视觉皮层切除，你就失去视觉------没有任何模糊的备份。所以大脑在每个功能区内部是局部全息的，其皮层柱架构是分形自相似的，但在全局上却\textbf{不是}全息的。它是一床拼布：几十块全息瓷砖缝合在一起，拼成一个整体------而这个整体，明确地说，是非全息的。

这种拼布结构解释了临床神经病学中一再出现的一个模式。小面积的中风和小的病灶往往只造成出人意料的轻微缺损，因为在任何一个给定的皮层区内部，全息原理都在保护着你。剩下的组织会以更低的分辨率重建缺失的信息。但大面积中风一旦抹掉整个功能区，就会造成灾难性的、特定的丧失（失明、瘫痪、失语），因为你把整块瓷砖从拼布里扯了下来，没有别的瓷砖能顶替它。

它还解释了为什么记忆不会因为神经元死亡就“凭空消失”。每一天都有神经元死去，突触被修剪。如果记忆是像文件存在硬盘上那样存储的，你会预期偶尔弄丢一份------某天早上醒来，忘掉了自己的婚礼、童年养的那条狗，或者咖啡的味道。这种事从不发生。相反，记忆是逐渐褪色的，在多年里慢慢失去细节和鲜活。这恰恰是全息存储系统所预测的：衰减是温和的、成比例的、全局性的，绝不会突然、离散或局部地发生。

拼布全息图，正是我上面描述的那些软件属性（尤其是克隆）之所以真正管用的物理原因。把大脑一分为二，每一半都保留着一份退化但完整的模拟拷贝，因为在每个半球内部，全息原理都确保了每一块都含有完整的画面。模拟不会崩溃，它只是以更低的分辨率运行罢了。


\chapter*{第4章：为什么会有感觉（以及为什么这个问题问错了）}
\addcontentsline{toc}{chapter}{第4章：为什么会有感觉（以及为什么这个问题问错了）}
\markboth{第4章：为什么会有感觉（以及为什么这个问题问错了）}{}

现在，我们可以直面意识的“难问题”了。

问题是：\textbf{为什么物理加工会有感觉？}

答案是：\textbf{它并没有。}

物理加工（神经元放电、突触传递、隐式模型存储与计算）没有任何体验。一丝一毫也没有。身处真实的这一侧，并不存在“是何感受”这回事。真实的一侧，恰恰就是那种“在黑暗中”运转的加工------难问题假定意识需要解释的，正是它。

有感觉的，是那个\textbf{模拟}。显式世界模型（EWM）与显式自我模型（ESM）------虚拟的一侧------才是体验安身之处。而在模拟内部，体验并不是附加在过程之上的神秘添加物。当模拟包含一个自我模型时，体验就是模拟本身\textbf{所是}的东西。ESM“感知”EWM，这就是我们所说的感受质（qualia）。感受质，是虚拟自我登记虚拟世界的方式。

不妨这样想。如果你问“为什么晶体管的开关切换会感觉像在运行一款电子游戏？”答案会是：“它不会。晶体管的开关切换没有任何感觉。游戏是一个运行在晶体管之上的虚拟过程，却拥有晶体管所不具备的属性------风景、角色、物理规则和光影。这些属性是虚拟过程的真实属性，而不是晶体管的属性。”

同样：神经元放电并不会感觉像看见红色。神经元放电生成并维持着一个模拟，而在这个模拟内部，自我模型把某一类世界模型的内容感知为我们所说的“红”。红，是模拟的真实属性，而不是神经元的属性。

难问题假定，我们需要解释物理加工如何产生体验。但物理加工并不产生体验------它产生的是一个\textbf{模拟}。而这个模拟，因为包含一个自指回路（ESM在EWM之内为自身建模），从根本上说就\textbf{是}体验。

\section*{醒来的那一刻，你是谁？}

先从一件你几乎肯定体会过的事说起。

你从某个很深的地方浮上来------一次晕厥、一次被打晕、一场麻醉，或者只是一觉漆黑无梦的沉睡，醒来时置身于一个你认不出的房间。而在一次心跳的工夫里，“里面”空无一人。觉知是有的------天花板、一道光带、“存在着”这一赤裸的事实------但还没有一个\textbf{你}挂靠在上面。你不知道自己在哪里。在漫长的一秒钟里，你不知道自己是\textbf{谁}。然后一切涌了回来：你的名字、你的过往、你人生的轮廓、你正躺在其中的这具身体。接缝合拢得如此之快，你几乎错过了它。但你没有错过。在那一次心跳里，存在着没有自我来认领的体验------觉知在运转，而自我模型还在启动之中，在房间里四处搜寻可供锚定之物，却一无所获。

抓住那道缝隙。整章要讲的，就是它。

因为它告诉了你一件事：那个重新“上线”的你，并不是像重新打开一份保存好的文件那样，从存储器里完好无损地取出来的。它是被\textbf{重建}出来的------就地组装，用底下的基质当场拼合而成。每天夜里，你的ESM都会瓦解，深睡眠会抹去正在运行的模拟。每天早晨，它重新启动，用隐式自我模型（ISM）------那个缓慢的、被存储下来的你，睡眠留下的版本------把“你”重新搭建起来。通常，这场重建瞬间完成、天衣无缝，你从来抓不到那道接缝。而在那个陌生的房间里，这一次，你抓到了。

所以，当我在本章余下的篇幅里论证体验\textbf{就是}一场运行在基质之上的模拟------那个被感受到的“你”是一个虚拟过程，而不是那些神经元------你手里其实已经握着凭据了。你曾亲身站在模拟加载完成前的那半秒钟里。现在，让我们把那半秒钟拆开来，问一问：这个加载出来的东西，为什么居然会有任何感觉。

\section*{循环性问题}

大多数读者会问的第一个问题是：“你不是只把问题挪了个地方吗？为什么\textbf{这个}模拟有体验，而天气模拟却没有？”

答案是自指。天气模拟为天气建模，却不为\textbf{它自己}建模。天气模拟有一个“外部”------计算机、程序员、解读输出结果的科学家。这个模拟可以被完整描述而无需提及任何体验，因为它内部没有自我模型。

大脑的模拟为它自己建模。ESM是这个模拟为\textbf{其自身过程}建立的模型。这就造出了一个闭合回路：模型与被建模之物是同一个系统。不存在一个可以完整描述这个模拟的“外部”，因为描述者本身就是描述的一部分。

这不是魔法，而是自指的结构性后果。当一个过程为自身建模时，模型与被建模者之间的区分便坍塌了。自我建模的过程，与身为一个自我的体验，并不是需要一座桥来连接的两样东西------它们是同一件事，只是用不同的词汇加以描述。

难问题要求在物理加工与体验之间架一座桥。四模型理论（FMT）说：没有桥，因为它们从未分开过。体验就是这场自我模拟本身------从回路内部看到的样子。

归根结底，这是一个同一性断言（在科学中，这类断言标志着一个落脚点，而不是一个缺口）。“水是H$_2$O”就是一个同一性断言。你无法有意义地追问“可\textbf{为什么}水是H$_2$O？”------同一性本身\textbf{就是}解释。再往更深处要答案，就是在索要另一种宇宙。同样：体验，就是处于临界性之中的四模型自我模拟\textbf{所是}的东西。如果有人问“可\textbf{为什么}这场自我模拟会有感觉？”答案是：因为这个过程\textbf{就是}这样一种东西。这个同一性是可证伪的------如果理论的预测落空，这个同一性就是错的。但它无法被“进一步解释”，正如水的分子同一性无法被进一步解释一样。它就是停驻之处。

\section*{为什么模拟不能“摸黑”运行}

这里还有一个更深的问题，而回答它，会揭示出意识为什么会\textbf{有感觉}的某种本质。姑且承认大脑运行着一场自我模拟。姑且承认四模型架构、临界性、自指闭合。难道这一切不能照常发生，却没有任何身临其中的\textbf{感受}吗？难道模拟不能一边评估、建模、预测------一边什么都感觉不到吗？

这是穿上技术外衣的僵尸直觉。答案是否定的，而理解其中缘由，会揭示这套架构最重要的特征。

基质把虚拟模拟部署为自己的评估机制。这是信息流的主方向：隐式系统把情境呈交给模拟，让模拟去评估后果、登记结果。但要让这种评估行得通，被模拟的状态必须带有\textbf{效价}------它们必须对模拟而言“有所谓”。一个只是数字的疼痛信号，无法在模拟层面驱动回避行为。只有一个\textbf{在乎}结果的模拟，才能评估结果。

还记得第2章讲过的数字孪生吗------那种针对喷气发动机的工程模拟。典型的数字孪生并不只是被动地映照发动机。它会\textbf{添加}一个可视化层：警告、用不同颜色标示的指示灯、警报------这些东西在物理发动机上并不存在。发动机出现金属疲劳，孪生体亮起闪烁的红色警告；发动机温度升高，孪生体的仪表就从绿转黄、再转红。这个添加出来的层，正是全部意义所在。没有它，孪生体不过是一张电子表格------一堆呆在内存里的数字，技术上准确，功能上无用。正是可视化，让这场模拟成为一件\textbf{评估工具}。

你的大脑做的是同一件事，而且做得更彻底。有意识的模拟并不只是映照基质的加工。它\textbf{添加}了现象性的效价。疼痛、愉悦、紧迫、好奇、恐惧、欣喜------这些就是大脑版本的警示灯和仪表盘指示。它们在基质层面并不存在（神经元不会感到疼痛，正如金属不会感到疲劳）。它们存在于模拟层面，\textbf{由}模拟亲手添加，好让系统能够一眼评估复杂局面。基质需要模拟来评估那些新奇、含混的情境------反射应付不了的那一类。而要让这种评估行得通，被模拟的自我就必须登记享乐效价：威胁、机会、后果。这种登记------这种“\textbf{有所谓}”------就是现象性。拿走感受质，你就拿走了评估------就像把驾驶舱仪表盘的显示屏整个扯下来，还指望飞行员靠读原始传感器电压来飞行。

“可是，强化学习系统也有驱动行为的奖励信号啊，”你也许会反驳，“它有感受吗？”没有------因为它缺少处于临界状态的四模型架构。强化学习的奖励信号，只是第一类或第二类系统里的一个标量值。而现象效价，是 ESM 在一场以第四类动力学运行的完整自我模拟之中对后果的登记------这是一个性质上截然不同的过程。二者的差别不在程度，而在架构。

模拟不能摸黑运行，因为黑暗会让它的目的落空。现象性并不是意识的附加功能，而是模拟完成自身工作的机制。

\section*{这不是什么：错觉主义}

这不是错觉主义。而这个区别重要到值得把话说透。

哲学界有一个颇受尊敬的立场，叫作错觉主义，与 Daniel Dennett 和 Keith Frankish 的名字联系在一起。它主张：感受质（qualia）是错觉。按照这种观点，“看见红色是什么感觉”这回事根本不存在。经验的显现本身就是一种虚构------大脑讲的一个故事，背后没有任何经验实在。最强意义上的意识并不存在，只是看起来存在而已。

想想这究竟在主张什么。假如你此刻正感受到什么------对这个论证的好奇、怀疑、手中这本书的重量------错觉主义会说：那种感受是错觉。你并没有真正体验任何东西。当你说“我感觉到了什么”，按照这套理论，你错了。你对自己经验的亲口证词是错的。你实际上等于在撒谎------只不过并不存在一个可以撒谎的“你”。如果你觉得这显然荒谬，我同意。

四模型理论说的恰恰相反。

感受质是真实的。它们在模拟之内是真实的。它们是虚拟自我感知虚拟世界的方式。当你的显式自我模型（ESM）登记下你的显式世界模型（EWM）对一只红苹果的表征时，这个登记（这次“看见红”）就是虚拟过程的一项真实属性。它存在于模拟层面，就像电子游戏里一颗子弹击中角色时，真的会\textbf{伤到}那个角色。这不是比喻------在游戏之内，伤害是真实的：血量下降，角色踉跄，世界随之反应。从外面看，那只是内存里一个递减的数字；从游戏里面看，那就是疼痛。这就是层面之差。你的感受质，就住在这里。

这套理论采用一种双层本体论。基质层面（神经元、突触、隐式模型）没有经验，它是灯灭的。模拟层面（显式模型、虚拟世界与虚拟自我）拥有真实的经验，它是灯亮的。两个层面都是物理的，谁也不是错觉。它们是同一个物理系统的不同层面，各自带着不同的属性。

这套理论并不说你的疼痛是错觉。它说你的疼痛是真实的------只不过真实在模拟里，而不在神经元里。而既然你的一生都在模拟之内度过，这就是对你而言唯一要紧的那种真实。

这是关键的区别。错过它，你就会把这套理论与取消主义混为一谈，与错觉主义混为一谈，与一切靠“解释掉”来解释意识的框架混为一谈。四模型理论并不把意识解释掉。它解释的是意识住在哪里------而答案恰恰是你一直站着的地方。

\section*{“在模拟之内是真实的”意味着什么}

这里有一处值得展开的哲学微妙之处。当我说感受质“在模拟之内是真实的”，你可能听出两种意思。要么它们是\textbf{真正现象性的}------那样一来，我不过是把谜团从神经元搬进了模拟，意识的“难问题”只是换了个地址继续活着；要么它们\textbf{在功能上真实、却并非真正现象性}------那样一来，这不过是丹尼特换了种绕圈子的说法。

这是一个假两难。它只有在你坚持存在一个“上帝视角”时才成立------一个能裁决某物是否“真正”具有现象性的外部视角，能够检验模拟是真的在感受，还是只是表现得像在感受。但自指闭合恰恰消除了这个外部视角。ESM 是它自己的观察者。不存在任何外部制高点可以用来发问：“可它\textbf{真的}有感受吗？”发问这件事本身，就是这个过程的一部分。

“真正现象性”对“仅仅功能性”的对立，预设了现象性是一种属性，一个过程要么具有它、要么没有，可以由一位独立观察者来核验。而对一个处于临界状态、完全自指的系统来说，这样的观察者并不存在。问题就此消解------不是因为它无法回答，而是因为它无从提出。它要求一种视角，而这种视角恰恰被自指闭合判定为不可能。

这是过程物理主义之内所能走出的最强一步，也是 Thomas Metzinger 借“现象透明性”这一概念所指向的立场------只是四模型理论对透明性\textbf{为何}产生说得更明确。制造这种透明性的，正是隐式与显式之间的边界：你无法看穿它，于是你无法跨到自己的现象性之外，去追问它是否“真正如此”。这道边界不是缺陷。它正是“真正现象性还是仅仅功能性”这个问题对你这样的系统不适用的原因。

我不是第一个走到这里的人------如果我是，我反倒该对自己起疑。

认知科学家 Joscha Bach 把同一个核心想法说得比几乎任何人都干脆。“你的经验是一场模拟，”他写道，“一场投射给虚拟观察者的虚拟现实。”而那个观察者呢？它是“一个模型，模拟的是假如你存在、被上传进一只猴子体内、一心牵挂着它的种种事务，会是什么样子”------一个为了驾驶这只猴子穿行于物理世界而造出的模型。一场模拟，模拟里有一个被模拟出来的“某人”来接收这场模拟，还有一具由这整套安排掌舵的身体。这一传统更早的支脉，我们在本书里已经见过------Metzinger 的自我模型，Dennett 的叙事自我。

而本书要再往下走一层------走进架构本身。究竟是哪些模型：世界与自我、隐式与显式构成的 2×2------那四个模型。一个系统究竟从何时起才算有“内部”的那条清晰界线：自指闭合，那个折返回自身的环路。以及让这一切活起来的条件：临界性，那个让系统既不僵死冻结、也不沦为白噪声的混沌边缘。

Bach 告诉你意识\textbf{是什么}。我要告诉你的，是它由什么\textbf{构成}、又在何时\textbf{点亮}。目的地相同，只是我把图纸也带来了。

\section*{神秘感为何挥之不去}

即便难问题已被消解，仍有一个缠着人不放的疑问。如果答案如此干净利落，为什么意识\textbf{感觉起来}依然那么神秘？为什么即使有人把解法告诉了你，难问题看上去还是那么难？David Chalmers 把这称作“意识的元问题”------即解释我们为什么会\textbf{认为}存在一个难问题。

四模型理论对此有一个干净的回答，而且它直接从架构里推出来。

奇怪之处在这里：那个有意识的“你”（虚拟自我）看不见生成它的机器。你无法内省自己的突触权重，就像梦中的角色无法检查做梦者的大脑。创造你经验的那个系统，就其本性而言，对你的经验是隐形的。不是因为有谁在藏它，而是因为它运作的层面，根本不在你的经验所涵盖的范围之内。

不妨这样想。你是电子游戏里的一个角色------一款非常出色的游戏，你在游戏世界内部拥有完整的自我觉知。你能看见渲染出来的群山，听见渲染出来的风声，感到脚下渲染出来的地面。可你几乎从不看见图形引擎，几乎从没瞥见过源代码。渲染过程发生在一个游戏世界通常并不包含的层面上。我说“几乎”，是因为偶尔会有痕迹渗漏进来。在你的大脑里也是如此------致幻剂会打开这道边界，心流状态会让它变薄，甚至在寻常日子里你也能捕捉到几瞥：大脑替你填补的盲点、揉眼睛时冒出的光幻视、闭上眼睑后浮现的几何图案。这些不是故障。它们是基质的运算从模拟内部短暂现形的时刻。我们会在第6章细细探究。但在大多数时候，渲染过程对它所渲染的世界是隐而不见的。

这正是 ESM 的处境。当有意识的自我试图理解自身经验的根基时，它撞上的是一种原则性的不透明------不是当前知识的缺口，而是架构本身的结构特征。生成模拟的隐式模型不是模拟的一部分。它们不可能是，就像 GPU 不可能成为游戏里的一座山。

结果可想而知。ESM 无法观察自己的基质，于是得出结论：意识的机制必定是非物理的，或从根本上无法解释，或不知怎么就超出了科学所及的范围。这就是二元论的源头。这就是“解释鸿沟”。这就是那种顽固的直觉------总觉得每一种对意识的物理解释都“漏掉了”什么。因为从模拟内部看，确实\textbf{有}东西被漏掉了：基质。恰恰是那个生成经验的东西，对它所生成的经验而言是隐形的。

这份神秘是真实的，但它是架构的产物，而不是任何非物理之物的证据。而它之所以\textbf{感觉}神秘，也自有其缘由。你是运行在生物硬件上的一个虚拟过程，而在大多数时候，你与基质之间的边界是不透明的。但并非永远如此。有时------在意识的改变状态里，在极度专注的瞬间，在眼角的余光中------你会瞥见底下的机器。看不清楚，也看不完整，却足以让你察觉：在你经验的表面之下，有某种庞大的东西正在运转。那种隐隐发怵的感觉，那种意识似乎深过你所能触及之处的感觉------那就是身为一场模拟的滋味：它几乎、却又不完全能看穿自己的那道幕布。

这是理论的一个\textbf{预测}，而不是一个未了的悬念。如果你是一个模拟，与自己的基底之间隔着一道几乎不透明的边界，那么你\textbf{理应预期}意识感觉起来恰恰就像现在这样奇异、这样不可还原。“难问题”的直觉冲击力，并非来自意识真的无法解释，而是来自我们所处的架构位置------我们身在模拟之内，透过裂缝向外窥视。

\section*{把自己缝合起来的自我}

现在，带上这套架构，回到之前那段空白。

那半秒钟的“无人在家”并不是故障，而是重建过程在半途被撞了个正着。身份并不是基底的固定属性------它是一种\textbf{重建}，每天清晨从储存下来的自我模型中重新组装而成。而它赖以重建的基底，从来都不完全是你入睡时的那一个。那些你不记得的梦已经重写了突触权重；一夜之间，记忆巩固又把你的记忆重新排布了一遍。醒来的你，已不完全是昨晚上床的那个人。这种差异通常小到你永远不会察觉------但它就在那里，每一天都在。

有两样东西让“你”在这种漂移中保持连续：一是隐式自我模型（ISM），它变化缓慢；二是睡眠，它在变化发生时让模拟下线，从而把变化藏了起来。把漂移推得足够远------一夜之间改写ISM、替换记忆、重塑人格------旧的“你”也不会就此消失，而是会被\textbf{吸收}。清晨的显式自我模型（ESM）会从残存下来的一切中重建出一条连续的叙事，把旧与新缝进同一个故事里，天衣无缝，从不察觉针脚何在。ESM不做干脆的断裂，它总是缝合。只有彻底的抹除才会扯断这根线；只要还有东西留存，明天的“你”就会悄悄把今天的“你”写进自己的历史，并把这一切称作同一个人生。

所以，这就是你：一个虚拟的自我，每夜重新组装，从一个自己爬不出去的循环内部感受着自己的世界。这样一来，只剩下一个问题还立在那里。我说过，模拟运行在一个维持于某种特定状态的基底之上------可那个状态究竟\textbf{是}什么？是什么样的物理刀锋，让几磅布满线路的细胞每天早晨启动出一个“某人”，又在每个夜晚归于黑暗？正是这一环让整个系统得以运转，它就是接下来要讲的。


\chapter*{第5章：在混沌边缘}
\addcontentsline{toc}{chapter}{第5章：在混沌边缘}
\markboth{第5章：在混沌边缘}{}

到目前为止，我告诉了你这套架构长什么样（四个模型、两条轴、一场运行在基底上的模拟）。我告诉了你体验栖身在哪里（在虚拟的那一侧，在显式模型之中）。我也告诉了你身份是什么（一次重建，每天清晨从存储的隐式模型里重新拼装而成）。

但我还没告诉你，是什么让这整台机器\textbf{运转}起来。为什么模拟有时开着、有时关着？是什么物理性质把一个有意识的大脑和一个无意识的大脑区分开？为什么深睡会抹去模拟，而架构却完好无损？

还有最后一块拼图，正是它真正让我确信这套理论行得通。

四模型架构对意识是必要的，但还不充分。你还需要正确的\textbf{动力学}。具体来说，基底（那个运行模拟的物理系统）必须工作在数学家和物理学家称之为\textbf{混沌边缘}的地带。

2002年，博学者斯蒂芬·沃尔弗拉姆（Stephen Wolfram）出版了《A New Kind of Science》，在书中他依据动力学把计算系统分成四类。我认为沃尔弗拉姆的分类需要第五类------他把分形系统和真正混沌的系统混为一谈，但两者在结构上截然不同。完整的论证在附录C，留给想看数学细节的读者。在这里，关键的一点是：

计算系统落在一条从完美有序到完美无序的谱系上。一端是静态和周期性的系统，太简单，算不出任何有趣的东西。另一端是混沌系统，太紊乱，稳定的模式无从形成。而在两者之间，在\textbf{混沌边缘}，坐落着能够进行通用计算的系统：复杂到足以产生丰富、多变、不可预测的行为，却又有序到让这些行为得以持续并相互作用。康威的生命游戏就是那个经典范例------正是我小时候在一台286上编程玩过的同一种元胞自动机。平坦网格上的三条极其简单的规则，却能生出滑翔机、振荡器、自我复制的结构，以及（可以证明的）通用计算。你可以在它里面造出一台计算机。你还可以在那台计算机里再造一台计算机。原则上，你可以在一片二维的像素网格里运行一整个三维的虚拟世界。从近乎虚无中，长出一切。

我想告诉你这个例子从何而来，因为它是这本书里一切的种子。生命游戏是我这辈子写过的第一个真正的程序------小时候坐在那儿，看着三条玩具般的规则在一片非死即活的像素网格里组装出一台能运转的计算机，在当时是我学到过的最重要的东西。一个更高维的世界从一套更低维的规则里爬出来，这本该是不可能的，而它明明就是可能的------这个事实嵌进了我心里，再没离开过。那时我还不知道，自己会用一生去追随这一个直觉，也不知道我自己双眼后方那片折叠的皮层，最终竟会施展同样的把戏，只是往上多了一维。这正是我们接下来要去的地方。

意识就栖身在这里。只有混沌边缘的动力学同时具备你所需要的两种性质：\textbf{通用计算}（复杂到足以真正运行一场自我模拟）和\textbf{全局整合}（系统中相距遥远的部分相互影响，局部的变化在全局传播，信息被绑定成一个统一的整体）。这就是为什么有意识的体验感觉起来是\textbf{统一}的------你不会把这边看到的红和那边听到的声音当成两条分开的溪流。临界动力学把一切绑定成一场体验。绑定并不是大脑在其他计算之外\textbf{额外}做的一件事，它是这套动力学状态的一个后果。

一个深睡中的大脑，被缓慢的δ波贯穿，运行在周期性动力学里：重复往返，哪儿也去不了。模型仍在基底之中，但模拟没有运行。一个陷入全面性发作的大脑，被一场同步共激活的洪流推离了那道临界边缘，坍缩成单一而粗糙的状态------僵硬的过度同步、失控的混沌，或两者交替出现------模拟再也无法维系。只有在清醒状态下（悬停在混沌边缘），系统才能持续撑起有意识的体验。

大脑作为一台被数十亿年演化打磨过的通用计算机，把\textbf{所有}计算状态都当作各司其职的工具来用：用稳定的吸引子做长期记忆，用周期性振荡做定时与门控（α、θ、γ节律），用分形处理做尺度不变的识别与纹理分析（主要在视觉皮层的V2-V4区），用混沌边缘的动力学来运转皮层自动机本身（意识的引擎）。只有混沌边缘这一状态才产生意识。但意识要正常运作，离不开其他那些状态。

当我在2015年的书里发表这个论证时，我完全不知道实证神经科学正独立地走向同一个结论。事实上，元胞自动机这一框架是整套理论里我最没把握的部分。当时我找不到任何实证支持，甚至差点没把它写进书里------它感觉像是迈得太远的一步，是一个会让理论其余部分很容易被人一笔勾销的主张。我留下它，是因为那逻辑似乎无可回避，而不是因为我手里有证据。

但这里有一个至关重要的微妙之处。仅有临界性还不够。一堆恰好在其休止角处崩塌的沙堆，正端坐在混沌边缘上。它并没有意识。这套理论要求跨过\textbf{两道}阈值：物理的那道（基底必须处于临界状态）和功能的那道（基底必须实现四模型架构）。有临界性而无架构，你得到的是复杂的动力学，却没有意识。有架构而无临界性，你得到的是一个休眠的系统------模型存在于基底之中，但模拟没有运行。两道阈值都必须跨越。合在一起，它们才充分。

我真正想象皮层的方式是一片海洋------但这是一片海水被彼此接通的海洋：不是每一滴水都挨着旁边水滴的那种平铺水面，而是一张网络，每个点都由习得的突触权重连向它自己特定的一批邻居。丢一颗石子进去------一个刺激、一个念头------一圈涟漪就沿着这些连线扩散开。而在这片网络之海上，正如在池塘上一样，模式可以成形：彼此加强的涟漪、驻波、那些自我维持的小小形状，它们游走、碰撞、留存。那些形状就是滑翔机。一个念头，或一个自我，就栖身在那里------作为一个骑在海面上的稳定模式，而不是作为任何一滴水。

\section*{两个旋钮，而非一个}

在\textbf{临界性}这个词内部藏着一个微妙之处，我花了很久才看清楚，而一旦你看清它，好几个谜题会同时消解。当人们说一个大脑“处于临界”时，他们通常是把两件不同的事塞进了一个词里。在网络之海上，它们是两个分开的旋钮。

\textbf{第一个旋钮问的是：一圈涟漪能传多远？}把它拧得太低，海面就成了一面平镜------丢一颗石子，只得到一圈原地就死掉的涟漪。什么都不扩散，什么都不整合。把它拧得太高，一颗石子就把整片海洋抽打成一场风暴------单单一次扰动就在一瞬间淹没一切。最佳点在两者之间：一圈涟漪往外传，大致催生出一圈新的，既不熄灭也不爆炸。就像一堆沙子恰好停在那个角度，多加一粒就能触发一场任意规模的滑落。这个旋钮关乎\textbf{触及范围}------单单一次事件能把多少片海洋拉进同一个相连的过程里。就叫它\textbf{广度}旋钮吧。

\textbf{第二个旋钮问的是：这片海洋上究竟能活下来哪些形状？}这就是那个生命游戏旋钮，整套理论都骑在它上面。一面冻结的镜子没有任何形状；什么都不动。一片温和而规律的涌浪只有一种形状，同一道波一遍遍重复------好看，却愚钝；它什么都算不出来。一锅翻滚的沸水让每个点都在随机闪烁------那是混沌，同样什么都算不出来。只有在有序与沸腾之间的那道边缘上，你才得到\textbf{滑翔机}：能游走、碰撞、留存、并做出些名堂来的驻定模式。就在那里，一次计算------以及一个自我模型------才能作为水面上的一个稳定形状而存在。就叫它\textbf{复杂度}旋钮吧。

我所知道的、描绘第二个旋钮最干净利落的画面，是一只机翼。飞机的机翼恰好在边缘处产生最大的升力------在气流分离、机翼失速陷入湍流的前一刻，它所能保持的最陡角度。再过那道边缘一丝，平滑的气流就碎裂，升力随之崩溃。大脑正是飞行在那道失速前的边缘上：计算力达到最大，就差一度动力学便要坠入混沌。再往前推一步，那份“升力”------承托滑翔机、进行计算、保持有意识的能力------就会像失速时的升力那样崩塌。

大多数时候，这两个旋钮一同升降，所以一个词通常就够用了。但它们可以分开，而一旦分开，它们就解释了那个单旋钮说法解释不了的一些情形。一锅沸水把第一个旋钮拧到了顶------一次扰动到处传开------但第二个旋钮被摁死在零：没有滑翔机，只有翻腾。传得很广，却没有形状，也没有意识，恰如你所料。

全面性发作恰好是一个例子，它最终敲定了第一个旋钮究竟在测量什么。在一场发作中，大脑里极大一部分同时放电------单论数量，比任何正常清醒时刻都更“活跃”。如果意识只是“许多神经元活跃着”，那一场发作就该是最有意识的状态。事实恰恰相反。那股活动的洪流从混沌边缘的状态里坍缩成单一而粗糙的形式------整片皮层被卷进一种粗糙的齐步节律，或被抛入失控的混沌------那种分化的、承载滑翔机的计算随之消失。所以第一个旋钮测量的不可能只是单纯的共激活；它测量的必定是大脑有多大一部分卷入了真正的混沌边缘动力学。在一场发作里，尽管放电铺天盖地，那个数字却崩了。大脑活跃了很多，和大脑意识清醒了很多，不是一回事。把这两个旋钮记牢；大脑所能进入的那些最古怪的状态，不过是它们的不同旋钮设定罢了。

\section*{皮层自动机}

是时候把某个可能还让人觉得抽象的东西落到实处了。我一直在说皮层需要工作在混沌边缘、处于第四类动力学里。可这个第四类系统究竟\textbf{是}什么？它不是什么盘旋在大脑之上的神秘力量。它就是神经放电的模式本身。

想想皮层在运转时实际长什么样。数十亿个神经元，每一个要么放电要么不放，每一个都通过习得的连接权重影响着它的邻居。每个神经元都是元胞自动机里的一个元胞，这不是比喻，而是字面意义上的。这台自动机的规则就是突触权重、阈值、局部连线。每个“元胞”的输出是一个放电率。而最终的结果------那幅以10到40赫兹在皮层表面上起舞的宏大电活动图景------是一台运行在数千维空间里的沃尔弗拉姆第四类元胞自动机。

我把它叫作\textbf{皮层自动机}。

这跟我小时候在一台286上编程玩过的（康威的生命游戏）是同一个想法------只不过这里不是一片带三条规则的平坦网格，而是一片折叠的皮层，带着数十亿条因地而异的规则；也不是在二维里移动，而是它的那些模式穿行于一个浩瀚到令想象力彻底失灵的维度空间。就像一只长着无数条手臂的章鱼，皮层自动机可以在任何时刻触及皮层的任何部位，随需激活它所需要的任何存储模型------这儿一段记忆，那儿一个运动计划，别处一个语言碎片。它像抓起一个个乐高小人那样抓取这些模型，用它们从一个令人满意的状态导航到下一个。

而这里有一个至关重要的区分：\textbf{皮层自动机不是意识}。它是引擎，不是体验。数十亿神经元放电所构成的那幅看似混沌的图景，实际上是一台异常精巧的装置，它进行计算、思考，并操纵一具身体走完一生。但意识只是这台装置的一个\textbf{效应}------一个在条件恰当时、从自动机与皮层的相互作用中涌现出来的效应。当自动机以恰当的频率、按一段连贯的时间顺序，同步地扫过合适的皮层区域时，一场有意识的体验就从这一连串画面帧中浮现出来。自动机装载着我们世界模型和自我模型的那些运行实例；意识，就是这些模型在模拟中活跃运行时所发生的事。

顺便说一句，你可以直接观察到皮层自动机------不需要fMRI。

试试这个：找一个完全黑暗的房间，闭上眼睛，等任何残留的余像消退（如果你之前一直盯着明亮的东西看，这大约需要 30 到 60 秒）。起初你什么都看不见，或者几乎什么都看不见。但接下来，如果你耐心等待、仔细留意，就会开始看到黑暗中闪烁的彩色光点。

大多数人把这些当作“视网膜噪声”不予理会------眼睛里的感光细胞对压力或自发的化学事件作出反应而产生的随机放电。而且，如果你轻轻按压眼睑，确实能以这种方式触发局部的视觉感受。但你在完全黑暗中看到的那些彩色光点\textbf{并非}来自视网膜。它们太有条理了，不可能是视网膜的产物。你所看到的，是 V1（你的初级视觉皮层）的静息活动，由残留感觉信号和皮层自动机本身自上而下的投射共同驱动。自动机正在运行它的基线动力学，而你正在实时地看着它发生。

如果你继续看下去------不是集中注意力，而是放松，让你的注意力变得柔和------就会发生某种非凡的事。主动的专注其实会压抑这些图案；恰恰是在你不再努力去看的时候，你才开始看见。自动机开始征用视觉系统中更多的部分，去解读并放大那里仅有的一点点信号。闪烁的光点稳定下来，凝成形状。几何图案浮现出来：网格、螺旋、格栅。接着是脸孔，扭曲而不断变换。接着是人形。然后，只要你有足够的耐心（我说的是\textbf{好几个小时}，不是几分钟），就会出现完整的场景------精致、彩色、带有叙事的幻觉，在种类上与你每晚做的梦并无二致。

这背后是与临睡幻觉相同的机制------就在你即将入睡时，在你脑海中闪过的那些生动图像。这是皮层自动机在几乎没有外部约束的情况下运行，通过激活存储的图案并把它们投射进模拟之中，生成它自己的内容。你所经历的那个渐进过程（从微弱的噪声到连贯的幻觉），是一扇直接窥看自动机运作方式的窗口：它从 V1 这个最早期的视觉处理阶段开始，随后逐步征用 V2、V3 以及更高级的区域，努力从任何可得的信号中理出意义。当没有真实信号可用时，它就自己\textbf{生成}一个。这就是通透性泄漏在起作用。没有外部信号来主导模拟，基质自身的处理噪声就变得可见。你并不是在幻视\textbf{虚无}------你看到的是这台图形引擎的空转图案，相当于神经版的、没调好台的电视雪花。只不过这片雪花是有结构的，因为那套处理机器本身就有结构。

你还可以用这种方式诱发一种暂时的联觉。年少时，我就用它来“看见音乐”。如果你闭上眼睛听音乐，同时让那些视觉图案自己浮现出来------放松、被动、不刻意用力去看，这些图案就会逐渐与你所听到的节奏和频率同步。皮层自动机在被剥夺了外部视觉输入之后，开始把它的视觉动力学耦合到任何其他可得的强信号上------在这个例子里就是听觉输入。你所看到的，简直就是你大脑活动的直接可视化：自动机的 V1 层图案，被听觉皮层而非视网膜输入所驱动。真正的联觉者------那些感官被永久交叉接线的人，一听到声音就总会看到颜色------很可能拥有同一种耦合的更持久版本，多半是因为感觉区域之间存在更强或更多的连接，无论是在丘脑还是在皮层本身。机制是一样的：一种感觉模态泄漏进了另一种的处理流水线。皮层自动机并不太在乎它的输入来自哪里。它处理它所接收到的一切。

我并不建议你把这当成一项常规爱好来尝试。这种体验可能令人不安，尤其是在你心理上没有准备的情况下。而且还有一点小概率，持续的感觉剥夺可能会让某些有潜在精神疾病易感性的人失去稳定。但如果你曾经好奇，当你的意识基质空转时------当外部世界安静下来，系统只是在……运行时------它究竟是什么样子，那么这就是不借助脑部扫描仪所能获得的最直接的一瞥。

那个从几乎虚无到一个完整虚构视觉世界的渐进过程，由你的自我模型在一个虚拟宇宙中亲历，正是皮层自动机运作时的一幅直接写照。

当自动机出错时，你同样能看到。当自动机的一部分陷入第一类或第二类动力学（周期性、锁定、在计算上毫无用处），或者被推过第四类进入第五类混沌，癫痫发作就发生了。中风，则是皮层的一部分彻底退出运作。晕厥，则是维持清醒所需的最低频率再也达不到。自动机在某种程度上是脆弱的。但生成它的那套结构------新皮层，连同它习得的权重和演化而来的架构------却是稳健的，这正是我们能从这些扰乱中恢复得如此出色的原因。

\section*{汇合}

2003 年------在我甚至还没有这套理论的两年前------约翰·贝格斯（John Beggs）和迪特马尔·普伦茨（Dietmar Plenz）在皮层组织中发现了“神经元雪崩”：一种遵循自组织临界性数学特征的神经活动图案，而自组织临界性正是处于混沌边缘的系统的标志。

2014 年，罗宾·卡哈特-哈里斯（Robin Carhart-Harris）提出了熵脑假说：意识水平与大脑活动的熵（无序度）相关联，其最佳点位于一个中间水平------熵太少意味着无意识，太多则意味着不连贯的体验。

2016 年，恩佐·塔利亚祖基（Enzo Tagliazucchi）及其同事表明，LSD 会把大脑推向临界性，这与迷幻药使用者所报告的那种被增强（但有时混乱）的意识相吻合。到 2022 年，已经有一篇综述论文能够谈及“作为意识框架的自组织临界性”------证据正在累积。

而在 2025 至 2026 年间，经验证据的堤坝决口了。基思·亨根（Keith Hengen）和伍德罗·舒（Woodrow Shew）在《Neuron》（2025）上发表了一项针对 140 个数据集的元分析------这是有史以来对大脑动力学中临界性所做的规模最大的系统性分析------证实了大脑在多种测量模态下都运行在一个临界点附近。随后，因巴尔·阿尔戈姆（Inbal Algom）和奥伦·施里基（Oren Shriki）在《Neuroscience \& Biobehavioral Reviews》（2026）中提出了 ConCrit 框架（意识与临界性），主张临界的大脑动力学为所有主要的意识理论提供了一个统一的机制基础。他们的结论是：意识追随临界性。当大脑处于或接近临界点时，意识就在场。当它被推到临界性以下时（由麻醉、由睡眠、由脑损伤所致），意识就缺席。当它被推离临界点、朝另一个方向偏移时------坍缩为锁步般的过度同步，或爆发为失控的混沌，正如癫痫发作、也可能是某些药物状态所致------意识就变得不连贯。

两条路径。一条是理论的，从沃尔弗拉姆的计算框架出发，推理一个自我模拟需要些什么。一条是经验的，从神经记录出发，分析大脑活动在每一种可触及的意识状态下的统计特性。两者的起源相隔不过两年，最终汇聚到同一个结论上。

这正是那种会让你认真对待一套理论的汇合。

\section*{全息图遇上自动机的三种方式}

在写这一章时，我意识到某件让我僵住不动的事。

全息原理和第四类自动机不断出现在同一些讨论里------在物理学中，在神经科学中，在计算理论中。但似乎没有人问过那个显而易见的问题：\textbf{它们之间可能存在哪些关系？}

恰好有三种。

\textbf{关系 1：一个全息基质产生第四类动力学。} 这大概就是大脑所做的事。神经网络在局部上是全息的------第3章里那块拼布式的全息图、拉什利的老鼠等等------而这个基质在临界性上运作，产生了意识所要求的第四类动力学。已被充分确立、详尽记录，而且（请原谅我）是无聊的那一种。

\textbf{关系 2：一个第四类自动机，把全息图案作为一种涌现行为产生出来。} 这个自动机在它的规则里并不是全息的，但它的动力学自发地生成了全息结构------高维信息被编码进低维图案之中，从计算本身之中产生。如果一个第四类自动机自然而然地产生全息式的输出，那就意味着非局域的信息分布从纯粹局部的规则中涌现出来------而这，耐人寻味的是，正是量子纠缠看上去的样子。

这里我该提一提杰拉德·特·胡夫特（Gerard 't Hooft），因为这个关联实在太惊人，不容略过------即便它仍是推测性的。特·胡夫特是一位诺贝尔物理学奖得主，他提出量子力学本身就是普朗克尺度上的一台元胞自动机：我们的宇宙从根本上是决定论的，而量子效应是一种更深层、离散的动力学的涌现现象。如果他是对的，那么我一直在描述的这一原理就不只是靠类比适用于意识而已。它就是字面意义上宇宙运作的方式，一直贯彻到最底层。简单的局部规则产生了一个全息宇宙，而在这个宇宙之内，简单的神经规则产生了一种全息意识。同一条计算原理在两个尺度上运作：宇宙学的和神经学的。我发现这种分形般的一致性极具说服力，但老实说，特·胡夫特的诠释在物理学界仍属少数派观点，而从结构上的优雅推向物理实在的论证，也曾遭到过正当的批评。尽管如此------如果最终证明有一条单一的计算原理同时构成了宇宙以及那些为宇宙建模的心灵的基础，那将是人类有史以来发现的最美的事实。

\textbf{关系 3：一个第四类自动机，其规则结构本身就是全息的。} 这就是那个让我放下笔的关系。如果这样一种东西真的存在------一台元胞自动机，其规则本身就把高维信息编码进一个低维结构里，就像全息图把三个维度编码进两个维度里那样------那么你就拥有了一个系统，它天生就在做全息原理所说的宇宙所做的事。不是一个仅仅\textbf{运行在}全息基质\textbf{之上}（或产生一张全息图）的系统。而是一个\textbf{本身就是}一种全息编码的系统。也可能就是宇宙------不过我应当指出，这仍是推测性的，而从数学之美推出物理实在的那个论证，也曾受到过正当的批评。我会在第14章回到这一点，届时我将解释为什么我认为关系 3 或许是数学中最重要的未解问题，然后在第 15 到 17 章里全力去追寻答案。


\chapter*{第6章：致幻剂揭示了什么}
\addcontentsline{toc}{chapter}{第6章：致幻剂揭示了什么}
\markboth{第6章：致幻剂揭示了什么}{}

开始之前必须先说一句：本章的任何内容都不应被理解为建议你去尝试致幻剂。它们威力强大、难以预测，能够毁掉一个人的一生------字面意义上的、永久性的毁掉。它们能在有易感体质的人身上诱发精神分裂症，能引起精神病性发作、持续性焦虑障碍，以及永远不会消失的 HPPD（致幻剂持续性知觉障碍）。我在这里讨论它们，是因为它们揭示了关于意识架构的某种重要东西。这份科学价值并不能使它们变得安全。

如果你想理解意识，就去研究它出错时会发生什么。在我看来，致幻剂是我们所拥有的、通向意识架构最富启发性的一扇窗------比睡眠者的脑部扫描更能揭示问题，比损伤研究在理论上更有信息量，而且比裂脑手术要容易获得得多。

原因在于：致幻剂不只是\textbf{改变}意识。它们以\textbf{系统化、可预测的方式}去改变意识，而这种方式恰恰揭示出底层的架构------只要你知道该留意什么。

\section*{通透性梯度}

回想一下隐式模型与显式模型之间的那道边界。在正常的清醒生活中，这道边界是有选择地通透的：相关的信息能穿过去，无关的信息则留在图书馆里。你觉知到你需要的东西，而对其余的一切浑然不觉。

致幻剂把这道边界炸开了。

在致幻剂（LSD、裸盖菇素、DMT、麦司卡林）的作用下，隐式---显式边界的通透性在整体上升高。那些平时完全在真实一侧被处理、对意识不可见的信息，开始向模拟这一侧渗漏。

而关键之处在于：它是\textbf{按顺序}渗漏进来的。

在低剂量下，或在体验的早期，最简单的处理阶段最先变得可见。这些是最靠近原始感官输入的阶段：V1 级处理。你会看到色彩变得更鲜明、静止的表面出现呼吸般的起伏、余光里有细微的移动。这些是视觉皮层早期的特征探测器，平时不可见，如今进入了模拟。

随着剂量增加或体验加深，更复杂的处理阶段变得可见。V2/V3 级处理：几何图案、分形、镶嵌纹样，以及 Heinrich Klüver 在 1920 年代所编目的那些著名的“形式恒常”。这些是视觉系统的中间表征------它平时用来建构你视觉体验的那些积木，如今独立地显现了出来。

再往上，更高级的视觉区域也变得可以触及。面孔浮现出来。人影。场景。面孔加工区、物体识别区、场景建构区------平时全都在意识阈限之下运作------如今把它们的中间产物直接播送进模拟。

在最高剂量下，整个处理层级都暴露无遗，结果便是彻底成形的幻象体验：复杂的、叙事性的、梦境般的场景，由隐式处理最深的那些层级建构而成。

这种有序的递进------从简单到复杂、从 V1 到更高区域、随剂量而变------正是四模型理论（FMT）所预测的。它是通透性梯度的直接结果：较低层级的处理阶段因为更靠近边界，会在通透性升高时先于较高层级变得可以触及。

下面这张视觉处理层级表展示了每个区域平时做什么，以及当通透性屏障降下时会有什么变得可见：


{\small
\noindent
\begin{tabularx}{\linewidth}{X X X}
\toprule
\textbf{区域} & \textbf{正常功能} & \textbf{致幻特征} \\
\midrule
V1 & 边缘、空间频率、朝向 & 光幻视、Klüver 形式恒常、呼吸般起伏的表面 \\[4pt]
V2 & 轮廓整合、纹理、边界归属 & 镶嵌纹样、重复的几何图案 \\[4pt]
V3 & 整体形状、动态形状处理 & 流动、变形的几何形态 \\[4pt]
V4 & 颜色、曲率、复杂纹理 & 彩色分形、万花筒般的图案 \\[4pt]
V5/MT & 运动处理 & 图案的旋转与移动 \\[4pt]
梭状回/IT & 面孔、物体、词形 & 面孔、人影、实体 \\[4pt]
前部 IT & 语义类别、场景建构 & 完整的叙事性幻觉 \\[4pt]
\bottomrule
\end{tabularx}
}


每一行都代表一个更深的处理阶段。在正常情况下，你只体验到最终的输出------完成后的知觉。在致幻剂作用下，你则会随着通透性升高，按顺序体验到那些\textbf{中间}阶段。（这张表更完整的版本，附有感受野大小和更多细节，见附录A。）

我知道这听起来很诱人。你正读到视觉处理的层层结构如何变得可见，心里有一部分好奇那究竟是什么样子。我懂------我当年也好奇。两条路我都试过。我那时年轻、愚蠢，也走运。上一章描述的冥想路径（一间黑屋、放松的注意、耐心）能把你带到同一个地方。没那么快，头一回也没那么戏剧化。但同样震撼、同样真实，而且不冒永久损伤心智的风险。你所需要的，只是一间黑屋里一张温暖的床。

还有另一条路：清醒梦。如果你能学会在梦中认出自己正在做梦------而这是一项可以训练的技能------你就能进入那个不受约束、全速运转的完整模拟。没有感官输入，没有外部现实来校正模型。只有那个虚拟世界，以及有意识地身处其中的你。对一些人来说，这比持久的冥想更容易做到。相关技巧有翔实的记载（实用指南见附录D），而这种体验至少和任何药物所产生的一样富有启示------却不冒风险。清醒梦我们将在第7章再回来讨论。

\section*{五个嵌套的系统}


\begin{figure}[htbp]
  \centering
  \includegraphics[width=0.95\textwidth]{../figures/figure-five-layer-stack-bw-zh.png}
  \caption{大脑组织的五个层级。每一层级都随附于（“运行在”）它下面的那一层之上。意识只存在于最顶端的虚拟层级，在那里显式模型生成现象体验。}
\end{figure}


把你的大脑想象成拥有五个截然不同的组织层级，像俄罗斯套娃一样一层套一层：

\textbf{物理层。} 最底层，是原始的物质：原子、分子、大脑自身的物理基质。这就是化学------构成组织的碳、氢、氮、氧。它是遵循热力学定律的惰性物质。这里没有任何有意识的东西居住。

\textbf{电化学层。} 往上一层：神经信号传递。动作电位沿着轴突飞奔，神经递质涌入突触，离子穿过通道流动。这就是人们一想到“大脑在做事”时所设想的那种电活动与化学活动。这一层是神经元放电的地方。仍然没有体验，但现在你有了信息传递。

\textbf{蛋白质组层。} 再往上：蛋白质结构与分子机器。突触权重就存储在这里------神经元之间连接的物理强度。细胞膜上的受体、调节可塑性的酶、决定哪些突触增强、哪些减弱的分子支架。这就是学习的“硬件”。当你练习一项技能并越来越娴熟时，你改变的正是蛋白质组这一层。仍然是无意识的，但现在你有了记忆。

\textbf{拓扑层。} 更高一层：网络架构。连接的模式，哪些神经元连向哪些神经元，连得多密，以什么样的构型。布罗德曼分区住在这里，皮层柱住在这里，“视觉皮层与运动皮层对话”这种大尺度结构就存在于这里。它是接线图。改变这一层，你就改变了这个系统能进行何种加工。你的隐式模型（IWM 和 ISM）就存储在这里。仍然是无意识的。但现在你有了知识。

\textbf{虚拟层。} 在最顶端：被模拟出来的世界。皮层自动机------那在网络上舞动的电活动动态模式，整合信息，生成预测，实时运行着那些模型。你的意识体验就住在这里。显式模型（EWM 和 ESM）存在于这里，也只存在于这里。这是唯一一个“感觉起来像点什么”的层。

每一层都随附于它下面那一层，却又有自己的动力学。没有物理物质就没有电化学信号传递，没有化学就没有蛋白质结构，没有突触就没有网络拓扑，没有一个网络来运行它就没有模拟。但每一层都拥有更低层所没有的性质。一个突触并不“关于”任何东西------它只是一个连接。而一个由突触构成的网络\textbf{却是}关于某个东西的：它表征一张脸、一个词、一段记忆。至于运行在那个网络之上的模拟呢？那正是“关于”变成“体验”的地方。

这个五层等级结构解决了一个几乎让每个初次听到这套理论的人都栽跟头的问题：“如果意识是虚拟的，那它到底运行在什么之上？”答案是：它运行在拓扑层（网络）之上，而拓扑层实现于蛋白质组层（突触权重）之中，蛋白质组层又架设在电化学层（神经元放电）之上，电化学层则存在于物理层（物质）之中。意识并不因为是虚拟的就少了几分真实------它只是\textbf{在一个不同的层次上}真实，而不是像神经元那样真实。电子游戏里的那座山，在游戏这个层次上是真实的，尽管在硬件层次上它“只不过”是一堆晶体管。同样的道理。

这正是五层等级结构大显身手、发挥其解释力的地方。致幻剂作用于这一堆栈的中部，其效应则向上层层波及。像 LSD 和裸盖菇素这样的经典致幻剂结合到5-羟色胺2A受体上，作用于\textbf{电化学}层------它们改变神经元彼此通信的方式。这一扰动传播到\textbf{蛋白质组}层，在那里受体敏感性会在数小时内发生偏移。它重塑\textbf{拓扑}层，网络的连接模式在这里发生改变------在 fMRI 上可见为全局整合度的提升。它还改造\textbf{虚拟}层，意识模拟在这里被平常不可见的内容所淹没。经典致幻剂唯一不触碰的层，就是\textbf{物理}层------它们不摧毁神经元，不改变原始物质。它们改变物质\textbf{之上}的一切，按由下而上的顺序。这是一个至关重要的区分。经典致幻剂（LSD、裸盖菇素、DMT、麦司卡林）没有神经毒性。它们改变神经元传递信号的方式，却不摧毁神经元。许多别的药物就没这么仁慈了。甲基苯丙胺和酒精会实实在在地摧毁神经元；可卡因造成的损害则更为间接，主要是通过收缩血管、使神经组织缺氧。MDMA 在高剂量或反复使用时会损伤5-羟色胺轴突。就连毒蝇鹅膏------那种标志性的红白相间蘑菇，很多人把它跟致幻蘑菇搞混------也是一种谵妄剂，其作用机制截然不同，也更危险。如果这一章你只记住一件事：并非所有能改变意识的药物都一样，而“改变信号”与“摧毁硬件”之间的区别，字面上就是暂时的意识状态改变与永久性脑损伤之间的区别。那种随剂量而变的视觉进程正好可以对应到这里：低剂量对电化学层的扰动恰好足以影响 V1 加工；更高的剂量把扰动向上推过更多层次，招募进越来越复杂的加工阶段进入意识体验。

\section*{可被改道的自我}

但最戏剧性的证据，来自自我身上所发生的事。

你的显式自我模型（ESM）是一个需要输入的虚拟过程。在正常状况下，它接收着一股稳定的自我指涉信号流：你对自己身体所在位置的感觉（本体感觉）、你对自己器官感受的感觉（内感受）、内心独白那道叙事之流，以及那份持续不断、身体层面的自我觉知背景------你从不会注意到它，除非它被打断。

在高剂量致幻剂之下，这股输入被打断了。自我模型并不会死去------它会\textbf{改道}。一旦被剥夺了正常的自我指涉输入，它就会抓住当下占主导的任何输入。

对此展示得最为戏剧化的，是鼠尾草属的墨西哥鼠尾草（Salvia divinorum），一种作用于κ-阿片受体的分离型致幻剂（与 LSD 或裸盖菇素的5-羟色胺能机制完全不同）。使用鼠尾草的人无一例外地报告了\textbf{变成}事物的体验：

\begin{itemize}
  \item “我变成了那张沙发。”
  \item “我曾是那面墙。”
  \item “我变成了书里的一页。”
  \item “我是电视上的某个角色。”
  \item “我变成了一个分形，不是看着一个分形，而是\textbf{作为}一个分形存在。”
\end{itemize}

这些不是比喻。使用者报告的是完整的、在体验层面令人信服的身份转移。在体验持续的那段时间里，他们\textbf{就是}那个物体或那个实体。有人形容那感觉像是死了------不是正在死去，而是\textbf{已经死了}------因为当你是一把椅子时，你曾经所是的那个人就干脆不复存在了。

内容追随着感官环境。看电视的人变成一个电视角色。躺在沙发上的人变成那张沙发。盯着某个图案看的人变成那个图案。

这正是显式自我模型在做理论所预测的事：一旦正常的自我输入被打断，它就改道去追随当下占主导的任何输入。身份的内容并非随机------它由感官环境所决定。只要控制住环境，应该就能控制身份的体验。

我得在这里让理论暂停一下。据我们所知，墨西哥鼠尾草是地球上最强的天然致幻物质。我刚刚描述的那种彻底的本体感觉接管，意味着身体觉知与空间定向的完全丧失。在鼠尾草作用下的人，有的从十楼的窗户走了出去。有的走进了车流。有的死了。这不是派对上的助兴药，也不是周五晚上拿来一试的猎奇玩意儿。它是现存对显式自我模型最极端的药理学破坏，而这种破坏可能致命------不是因为这药有毒，而是因为你不再知道自己的身体在哪里，还可能彻底相信自己长了翅膀、能够飞翔。

许多尝试过鼠尾草的人报告说，那种体验感觉像是在死去------不是打个比方，而是一种真切的、令人瘫痪的确信：自己已经不复存在。这是显式自我模型彻底崩塌到如此地步，以至于模拟根本再也无法生成一个“你”。我们会在第8章看到这种情形的临床对应物，届时会谈到科塔尔妄想------那些在神经学层面确信自己已经死了的患者。鼠尾草则以药理学的方式把你带到那里，只需几秒钟，毫无预兆。想一想，这是不是你想去体验的东西。

我自己也体验过那种主观时间的延展。在鼠尾草的作用下，半秒钟的真实时间------由在旁看着我的人所确认------延展成了感觉像是十五分钟甚至更久的东西。我的知觉世界重建成了一段段繁复的序列，其中包括了长着翅膀四处飞翔的感觉（那种飞的感觉，我后来意识到，来自我向后倒在床上时从身旁掠过的气流）。我的整个现实崩塌又再生，全都发生在眨一次眼的工夫里。我把这些讲给一位替我计时的旁观者听，他说我“消失”了还不到一秒钟。这跟我早几年、在1998或1999年一次濒死事件中已经体验过的（那个机制我会在第13章描述）是同一种主观时间延展，只不过这一次是由药物诱发的，而且还要极端得多。

我并不是最戏剧化的案例。一份翔实记录的报告讲的是一个男人，在一次墨西哥鼠尾草发作期间，他经历了感觉像是整整八年的另一种人生------上学、交朋友、建立一份全新的存在------而这次发作按钟表时间只持续了大约四十五秒。经同行评审的研究证实了在受控条件下会出现极端的时间扭曲，其中一位参与者把时间描述为“像手风琴一样折叠起来”（Addy et al., 2015）。基底把如此之多的内容如此之快地推送进模拟，以至于主观时间与钟表时间彻底脱钩。

这从未在受控环境中经过实验检验。但它是可以被检验的，而那将是对这套理论最具特色的机制的一次戏剧性证实。

大脑为什么能把这么多内容这么快地推送进模拟，这跟一个将死之人为什么会在几秒钟里看见自己的一生闪过，是同一个问题------在第13章我会把两者放在一起再来谈，届时会发现它们共享着同一个机制。

值得追问的是，为什么这个角落------两个旋钮同时被推到高位，大脑的大部分被卷进单一进程之中，\textbf{并且}这一进程又以最大的丰富度运行着------会如此罕见。通常有两样东西禁止它出现，而这两样彼此互补。

第一样是简单的经济学。同时做这两件事在代谢上是毁灭性的；大脑负担不起超过片刻的时间，能量预算把这扇门死死关着。这就是为什么需要某种极端之事才能把它撞开------一个代谢刹车正在失灵的将死之脑，或是一种扰乱了那套配给预算之调节机制的药物。

第二样更为奇特，也正是这些状态之所以感觉像是把整个世界抛在身后的原因。你大脑的\textbf{内部}世界，其维度远远高于那些把它连向现实的细细管道------两只眼睛、两只耳朵、一具身体的触觉，对峙着一个丰富得令人瞠目的内部模拟。通常那条窄窄的感官管道就够用了，因为它在不断地\textbf{校正}那个内部模拟：一股持续不断的“这才是外面实际有的东西”，把那奔逃的想象拴系在世界上。但把两个旋钮都推高，内部动力学就会膨胀，直到那条细细的输入管道再也无法把它校正回来。模拟跑赢了自己的现实核查，锁进它自己的吸引子，飘离了世界。而妙就妙在这里：干这件事的正是那个\textbf{广度}旋钮------把更多的大脑向内卷进来，恰恰就是淹没那根拴绳的东西。往这个角落里推，就是切断了那条本该让你脚踏实地的绳索。

这就留下了一个至今仍让我感到惊愕的反转。那个角落同时是\textbf{意识的巅峰，与现实的最少接触}------最多的加工，最少的世界。这恰恰就是为什么那些双旋钮全开的状态（濒死时的人生走马灯、高剂量鼠尾草的坠落、最深的致幻沉浸、最鲜活的梦境）会是那些感觉最脱离现实的状态。它们并不是\textbf{尽管}强烈却仍脱离现实。它们脱离现实，恰恰是\textbf{因为}它们如此强烈。

如果你想看看这条原理能延伸到多远，不妨来做一个思想实验。设想有个人被长期维持在极高（但尚未完全解离）剂量的沙维诺林A（salvinorin A）之下------它是墨西哥鼠尾草（Salvia divinorum）中的活性成分，只作用于单单一种受体（κ-阿片受体）。这个人的显式自我模型（ESM）将永远无法稳定下来。它会无休止地随着当下恰好占据主导的输入而循环：这一刻他相信自己是一把椅子，下一刻是一张桌子，再下一刻是一只恐龙，然后是空气，然后是一张纸。他仍会\textbf{体验}到事物（视觉和听觉仍在运作），但他再也不会知道自己是谁、是什么。撤去药物，随着时间推移，正常的自我模型会从完好无损的隐式自我模型（ISM）中重新组装起来。

这一点很重要，因为它表明：意识并不需要一个\textbf{正确}的自我模型，它只需要\textbf{一个}自我模型。无论如何，这套架构照样运转。显式自我模型（ESM）不会因为收到荒谬的输入就停摆------它会用手头一切可用的信号，尽力构建出最好的一个自我。这与我们在科塔尔妄想（ESM在内感受信号缺失的情况下工作：“我一定是死了”）、安东综合征（眼睛失灵时，EWM从记忆里生成视觉）以及转换障碍（ESM为一个底层其实并不存在的瘫痪建模）中看到的是同一条原理。自我模型是个强迫性的建造者。它从不停止建造。它从不宣告数据不足。它只管建造，并且相信。

\section*{病感失认：反向的情形}

这里有一种美妙的对称。如果说迷幻药是隐式---显式边界变得\textbf{过于}可透过时发生的事，那么病感失认就是这道边界变得\textbf{过于}不可透过------至少在局部如此------时发生的事。

病感失认最常见于右半球中风之后，患者对自己的功能缺损确确实实浑然不觉。一位左臂瘫痪的患者会坚称手臂好好的，会为自己用不了它的种种失败找借口开脱，一旦被摆到瘫痪的证据面前，就会困惑或恼怒。他们并不是心理意义上的否认------手臂已经瘫痪这一信息，根本就从未抵达他们有意识的模拟。

在四模型理论（FMT）里，这是隐式---显式通透性的一次局部下降。隐式自我模型（ISM）\textbf{拥有}关于瘫痪的信息------底层记录下了损伤。但这道边界针对那个特定领域被堵住了，于是显式自我模型（ESM）从不把这项缺损纳入其中。患者的模拟里不包含一条瘫痪的手臂，所以患者也就体验不到它。

这套机制其实更为具体，一旦你看清它，它就以一种略带惊悚的方式显得优雅。当你的运动系统发出一道指令（比如“拍手”）时，它同时做了两件事。它把指令发给肌肉，又把\textbf{预测的反馈}发给意识：根据以往经验，拍手\textbf{应该}是什么手感、什么声音。这份预测的反馈会\textbf{先于}真实的感觉反馈抵达，因为真实的反馈得走更慢的神经通路。在正常情况下，这份预测很快就会被真实的感觉数据修正或确认。你预测那声拍手，接着你感觉到、听到那声拍手。吻合。继续。

在病感失认里，来自瘫痪肢体的真实反馈永远不会抵达。而那个本该发出“等等------什么都没发生”警报的机制，已经损坏了。于是预测的反馈就这样未经修正地通过。患者的运动系统命令双手拍在一起，把一次双手拍手的预测发给意识，意识便体验到了恰恰是这样的东西------一次双手完成的、完全正常的拍手。患者会满怀真诚地告诉你，他刚才用双手拍了手。他听到了。他感觉到了。他体验到了。在他的模拟里，这事发生了。只是它没在现实里发生。

意识\textbf{就是}这样运作的，一直如此，我们每个人都如此。唯一的区别在于，在健康的人身上，预测的反馈会在几毫秒内被修正。而在病感失认里，修正机制坏了，患者的模拟便干脆只靠预测独自运行。

迷幻药与病感失认是同一套机制朝相反方向的运行。一个让通透性在全局范围内升高，另一个让它在局部范围内下降。而这种对称催生了一个跨领域的预测：迷幻药应当能缓解病感失认。全局的通透性升高应当会压倒局部的堵塞，让那项关于缺损的信息得以抵达意识。

从来没有人检验过这一点，因为从来没有人拥有一套把这两种现象连起来的理论。若没有四模型理论（FMT），这条联系是看不见的。


\chapter*{第7章：当灯熄灭时会发生什么}
\addcontentsline{toc}{chapter}{第7章：当灯熄灭时会发生什么}
\markboth{第7章：当灯熄灭时会发生什么}{}

每天晚上，你都会失去意识。每天早晨，你又把它找回来。而这两者之间的过渡（穿越各个睡眠阶段的旅程），正是临界性原理每夜上演的一场演示。

\section*{深睡：阈值之下}

在深度非快速眼动睡眠中，大脑的动力学转入亚临界状态。它的标志是慢波：大规模、同步的振荡，其中庞大的神经元群一起放电，又一起归于沉默。这是第二类动力学------周期性的、重复的，太过有序，无法承载意识。

由马尔切洛·马西米尼（Marcello Massimini）及其同事提出的扰动复杂度指数（PCI）直接印证了这一点。PCI衡量的是大脑对一次磁脉冲的反应有多复杂：在清醒的意识中，这种反应复杂而有分化（PCI高）；在深睡中，它简单而刻板（PCI低）。深睡中的大脑无法维持一个有意识的模拟所需要的那种丰富、全局整合的动力学。

灯熄灭了。显式世界模型（EWM）和显式自我模型（ESM）都已崩塌。没有模拟，也没有体验。

\section*{梦：降级运行模式}

但在快速眼动（REM）睡眠中，灯又亮了起来。大脑的动力学重新朝临界性靠拢------没有完全回去，但已足够接近。模拟重新启动，你又体验到了一个世界。

这是一场降级运行的模拟。正常的外部输入被切断了（你的眼睛闭着，你的肌肉处于麻痹状态）。显式世界模型（EWM）依靠内部数据运行------从隐式世界模型（IWM）储存的知识中提取素材，而不是从当下的感官输入中提取。这就是为什么梦里会出现熟悉的地点和人物，却带着不可能的物理规律和支离破碎的叙事：模拟正在用有限的输入尽其所能。

显式自我模型（ESM）同样在降级运行模式下工作。你把梦体验为发生在“你”身上的事，但你的元认知监督被削弱了------你会毫不怀疑地接受不可能的事件，你很少察觉到自己在做梦，你的批判力被调低了。

那种降级运行的模拟从内部感受起来是什么样，我很清楚。大约七岁到十二岁之间，我做过一个反复出现的梦------在那些年里，同一个梦回来了好几次。它是一片风景，如果那也能叫风景的话：一个无限延展的分形结构，山谷向下坠落，永无止境，无论你看向哪一个深度，它都是自相似的。伴随它而来的是一种我在清醒时从未有过的感觉------一种深刻的脱离身体感，仿佛我的身体被剥去，而我被安放进了那几何结构本身之中。那个梦诡异、陌生，却又奇异地令人着迷。哪怕它吓到我，我也想让它再回来。

真正令人惊讶的是，我至今仍然能\textbf{感觉到}它。几十年过去，每当我想起那个梦，同样那种脱离身体的感受就会回来------不是对一种感觉的记忆，而是那感觉本身。那份隐式的编码依然在那里，藏在突触权重之中，被三十年的清醒经验触碰不到。基底把它储存得如此之深，以至于仅仅去想它，就足以让那场模拟部分地重新激活。

那个梦当时在做什么？用这套理论的语言来说：显式世界模型被切断了外部输入，就从基底自身的计算动力学中生成内容。而那些动力学又是什么呢？第四类------它把第三类（分形结构）作为一个子过程包含在内。我沉睡的大脑，除了自身的架构之外别无凭借，就这样运行着它的模拟，于是产出了这架构的视觉签名：分形。那种脱离身体感，就是显式自我模型在它最降级的形式下运行------模拟知道自己是\textbf{某个人}，在某个地方，但身体的表征几乎已经完全脱落。

多年以后，当我学会了清醒梦，我却从未试图把那个梦重新召唤回来。也许我该试试。

梦游是一场更加戏剧化的演示。在梦游中，运动系统部分地重新激活，而显式自我模型仍然离线，或近乎离线。基底在运行运动程序------走路、辨向，甚至完成复杂的动作------但模拟并没有真正参与其中。梦游者穿行于物理世界，靠隐式世界模型的空间知识引导，却几乎没有、甚至完全没有有意识的体验。

这一点我有亲身经历。青少年时期，我有过一阵子梦游。有一天早上，我醒来发现自己坐在书桌前，面前摊着潦草的笔记，是用左手写的------而我清醒时从不用左手写字。我有一段零碎的记忆：沿着墙一圈圈地走，想找到门，却找不到。但坐到书桌前、试图写字的那一段------完全是一片漆黑。基底在辨向，运动程序在执行，可模拟------那个“我”------并不在场。

这就是这套理论的微缩版。一具身体在世界中移动，处理空间信息，执行学过的运动程序，而这一切之中并没有一个有意识的自我。隐式模型在掌控全局。显式模型处于离线状态。结果就是一个会走路、会行动、甚至会写字的人，可里面空无一人。

\section*{清醒梦：那个开关}

然后还有清醒梦------在这种状态里，你还身处梦中，却意识到了自己正在做梦。在四模型理论（FMT）中，这是显式自我模型在梦境状态内更充分地“切换开启”。这是自我建模能力的一次阶跃式增强。

学会做清醒梦，就是学会抓住模拟在犯错的那一刻。我用的方法叫作电灯开关测试：一整天里，你养成拨动电灯开关的习惯，并问自己灯的变化对不对。在清醒的生活里，灯总是会正确地变化。在梦里，电灯开关不管用------依靠内部数据运行的显式世界模型，懒得去模拟电路的物理过程。当你在梦里拨动开关，灯却没有变化，或者房间反而变暗了，或者开关的手感不对劲------那种不匹配，就是你的显式自我模型察觉到了模拟里的一处不一致。而在你察觉到它的那一刻，你就知道自己在做梦了。

我只花了几天时间。我很幸运------大多数人需要几周甚至几个月的练习，这个习惯才会转移进他们的梦里。但当它奏效时，那个过渡正是这套理论所预测的阶跃式阈值。前一刻，我还是梦境叙事里一个被动的角色。下一刻，我就完全在场了------知道自己在做梦，知道周围的世界是被生成出来的，知道自己可以改变它。ESM切换开启了。基底没有改变。动力学也没有改变。但自我模型与模拟之间的耦合越过了一道阈值，一切都感觉不同了。

这套理论预测，这个过渡对应着一次临界性阈值的跨越。不是大脑复杂度的逐渐上升，而是一次骤然的阶跃。如果你在清醒感刚开始的那一刻前后、一个严格对齐的时间窗口内测量脑电（EEG）复杂度（利用那个成熟的范式------由清醒梦者事先约定好的眼动信号），你应当会看到一处不连续。

\section*{麻醉：两种类型}

麻醉为临界性原理提供了最干净的一次检验，因为不同的麻醉剂尽管被归在同一个标签之下，却会产生截然不同的体验。

\textbf{丙泊酚}把大脑推向亚临界。丘脑---皮层连接被破坏，皮层复杂度崩溃，PCI趋近于零。灯彻底熄灭。患者报告在丙泊酚麻醉期间没有任何体验。这正是这套理论所预测的：把动力学推到临界性之下，模拟就无法维持。

\textbf{氯胺酮}做的事情则完全不同。它\textbf{并不}把大脑推向亚临界。脑电研究表明，氯胺酮\textbf{提高}了神经熵------它把大脑推向临界性，或推过临界性，进入一个更混沌的状态。结果呢？就是“K洞”（K-hole）------那种生动、往往怪诞的体验：解离、扭曲的现实、灵魂出窍般的体验，以及身份的剧烈变换。

在四模型理论（FMT）中，K洞是意识在\textbf{错误}输入上运行。显式世界模型和显式自我模型仍然活跃（大脑仍处于临界性，或在其之上），但外部的感官处理被破坏了。模拟依靠内部信号和被扭曲的信号运行，于是产生了K洞特有的现象学。

这一区分（丙泊酚通过跌入亚临界而取消意识，氯胺酮则通过在输入被破坏的同时越过临界进入超临界而改变意识）是一项真正的解释优势。大多数理论都难以解释，为什么两种“麻醉剂”会产生如此天差地别的体验。而临界性框架则让这一区分变得顺理成章。

我从没做过化学麻醉。但我曾被一击打晕过------猛烈、突然、毫无预兆。回想起来，最让我震惊的，是那种彻底没有过渡的感觉。睡眠有各个阶段。致幻剂有一条起效的曲线。哪怕是起效极快的鼠尾草（Salvia），也会给你零点几秒的“有什么正在发生”。被打晕却什么都不给你。前一刻，模拟还在运行。下一刻，你就在地上醒来，完全不知道过了多久。没有变暗，没有褪去，没有隧道。模拟并不是逐渐劣化，而是直接终止。断路器跳闸了。

这正是你能预料到的、皮层动力学突遭剧烈扰动时会发生的事------一瞬间就被推到远低于临界性的状态。没有时间让系统沿着睡眠各阶段优雅地缓缓关闭。系统不会经由第四类过渡进第二类，而是彻底从第四类里崩出去，模拟就这么停了。

\section*{意识地图}


{\small
\noindent
\begin{tabularx}{\linewidth}{X X X X}
\toprule
\textbf{状态} & \textbf{临界性} & \textbf{模型} & \textbf{意识} \\
\midrule
正常清醒 & 处于临界 & 四个模型全部活跃 & 完整 \\[4pt]
REM 睡眠 & 近临界 & EWM/ESM 依赖内部输入 & 劣化（做梦） \\[4pt]
深度 NREM & 亚临界 & EWM/ESM 已崩塌 & 缺席 \\[4pt]
丙泊酚 & 被强制压到亚临界 & EWM/ESM 被抑制 & 缺席 \\[4pt]
氯胺酮 & 越过临界（$\uparrow$熵） & EWM/ESM 依赖错误输入 & 在场，但断联 \\[4pt]
致幻剂 & 处于/越过临界 & 全部活跃，$\uparrow$通透性 & 在场，但被改变 \\[4pt]
清醒梦 & 近临界，跨过阈值 & EWM 活跃，ESM 完全投入 & 增强的自我觉知 \\[4pt]
\bottomrule
\end{tabularx}
}


这张表总结了本章讲过的一切，也提供了一处你可以随时回来查阅的参照。你这辈子体验过的每一种意识状态，都能在这张地图上找到位置，而决定位置的是两个因素：你的基质是否处于临界性，以及四个模型中有哪些正在运行。睡眠、麻醉、致幻剂、梦境、K 洞------它们并不是各自独立的谜团，而是同一张地图上的不同坐标。

但这张地图上的每一种状态，都是一颗健康的大脑，到了早晨会自己找到回家的路。更棘手的问题是：当机器里有一部分零件回不来了------当损伤是永久的，架构必须在自身缺了一块的情况下继续运行------这张地图会是什么样子。要回答这个问题，你去的不是睡眠实验室，而是神经科病房。


\chapter*{第8章：临床之镜}
\addcontentsline{toc}{chapter}{第8章：临床之镜}
\markboth{第8章：临床之镜}{}

同一套四模型架构既能解释睡眠与麻醉，也能解释临床神经病学中一些最戏剧化、最令人费解的病症。它们不只是有趣的病例------它们是架构中某些特定部件失效时发生的情形。而每一次失效，都从不同角度照亮了这套架构，就像一根烧断的保险丝会告诉你它原本在保护哪条电路。

如果这套理论确有价值，那么对特定模型的损伤就应当产生特定的、可预测的缺陷。不是含糊其辞地说“意识受损了”，而是精确的预测：敲掉这个部件，你就会得到\textbf{那种}综合征；让另一个部件在缺少正常输入的情况下继续运转，你就会得到\textbf{这种}另外的综合征。临床文献里充满了这样的病症：在标准的意识模型下它们深不可解，可一旦你手握真实/虚拟之分、又有四个相互作用的模型可用，它们就自然而然地各归其位。

\textbf{盲视与安东综合征：完美的镜像}

如果本章你只记住一件事，那就记住这一对。其他每一种意识理论，连解释其中一种病症都举步维艰。而四模型理论对两者都做出了预测。

先从盲视说起。一位患者的初级视觉皮层受损------那是大脑中生成有意识视觉体验的部分。按任何标准的临床检查，这位患者都是盲的。问他看见了什么，他会告诉你：什么也没有。他是认真的，既不是谦虚，也不是糊涂。就他的有意识体验而言，视觉世界根本就不存在。

可接着，令人惊讶的事情发生了。研究者在走廊里摆上障碍物，请这位患者穿行而过。他反对------他什么都看不见，怎么可能走得过去？他们坚持。他叹了口气，站起身，走了起来。

他毫无差错地穿过了障碍路线。绕开椅子，从一道上次并不在那里的横杆下弯腰钻过，在两个障碍物之间的缝隙里灵活穿行------而与此同时，他仍然真诚而由衷地坚称自己一丝一毫都看不见。这有录像------我建议你去找来看看，因为光靠文字读到它，无法还原它的震撼。一个临床意义上的盲人像视力完好一样在障碍路线中穿梭的画面，是整个神经科学中最惊人的演示之一。在旁观看的研究者，个个像见了鬼一样。

怎么做到的？因为基底仍在处理视觉信息。隐式世界模型（IWM）通过绕开受损皮层的皮层下通路接收视觉输入（一条从视网膜经上丘到丘脑枕的快速通道），这条通路的演化远早于皮层的出现。它构建空间地图，引导运动行为，让身体避免与物体相撞。可这一切都没有抵达显式世界模型（EWM）。有意识的模拟里不含视觉。这位患者真真切切地体验着失明------又真真切切地凭视觉在导航。基底在没有模拟的情况下运转。

现在把它反过来。安东综合征（对皮层性失明的病感失认）恰恰是它的完全反面。这些患者是真真切切、彻底地失明。他们的视觉皮层或视觉通路已被摧毁，没有任何视觉信息抵达大脑。可他们却绝对地、不可动摇地坚信自己看得见。

他们撞上墙，还怪家具摆错了地方。他们满怀信心地描述房间里根本不存在的物体------“桌上有个蓝色花瓶”，可桌子是空的。请他们说出你举起的是什么，他们会给你一个答案，平静而毫不迟疑，而且这答案会是错的。若拿他们失明的证据当面质问，他们会先是困惑，继而恼火，再是愤怒。灯光不好。他们需要配副新眼镜。他们只是刚才没留神。他们没有撒谎，也不是心理学意义上的否认。他们真真切切地、在体验层面上\textbf{看见}------而他们所看见的，与真实世界毫无对应。

在四模型理论里，这正是显式世界模型（EWM）从隐式世界模型（IWM）储存的知识中生成一段视觉模拟------尽管此刻并没有任何视觉输入抵达。这段模拟靠的是旧数据，是预期，是大脑对世界应当是什么样子的最佳猜测。患者“看见”了一个并不在那里的世界。模拟在没有当前输入的情况下运转。

把它们并排放在一起。盲视：基底处理了视觉，但模拟没有把它显示出来。安东综合征：模拟显示了视觉，但基底根本没在接收。有基底而无模拟。有模拟而无输入。如果你把意识看作单一而统一的东西，这两种病症都深不可解；而对一套区分真实处理与虚拟体验的理论来说，它们却是自然而然、甚至可以预测的结果。就算你刻意去设计，也几乎设计不出比这更好的一对检验案例。

\textbf{隐匿的觉知：困在里面}

2006年，阿德里安·欧文（Adrian Owen）和同事发表了一项研究，改变了我们对植物状态的看法。他们把一位被诊断为植物状态（无反应、看似无意识）的患者放进 fMRI 扫描仪，请她想象自己在打网球。她的大脑亮起的模式，与一个健康、有意识的人想象同一件事时的模式完全一致。

她在里面。有意识、有觉知、在思考------却完全无法移动、说话，或向任何人示意自己的存在。

四模型理论在这里做出了一个干净利落的区分。一个真正的植物状态患者，基底处于亚临界状态。动力学已跌落到阈值之下。模拟没有在运行。里面空无一人------不是因为这个人“离开了”，而是因为生成模拟的那套计算架构已经离线。

可一个隐匿地保有意识的患者，则是全然不同的另一回事。基底处于临界状态------动力学丰富到足以维持一段模拟。显式世界模型（EWM）和显式自我模型（ESM）都在运行。这个人在体验、在思考、在感受。但输出通路已被摧毁。这段模拟没有办法把自己表达出来。这个人有意识，却被锁在里面，困在一具不肯回应的身体之中。

扰动复杂度指数（PCI，与区分睡眠阶段的是同一项指标）应当能把这些案例区分开来。事实也的确如此。一些被诊断为植物状态的患者，测出的 PCI 值端端正正落在有意识的范围里。他们根本不是植物状态。他们是囚徒。其中的医学与伦理意义极其重大，而四模型理论恰好告诉你这一区分为何存在、又该如何去检测它。

\textbf{科塔尔妄想：“我死了”}

再往下，还有一些患者坚信自己已经死了。

科塔尔妄想是精神病学中最离奇的病症之一。患者坚称自己已经死去。他们相信自己的器官已经溶解，血液已经流干，自己不复存在。有人相信自己正在腐烂。有人相信自己不朽------因为如果你早已死了，就不可能再死一次。他们不是在打比方。他们怀着彻底的、不可动摇的信念，字字当真。

到这里，你应该能认出其中的机制了。它和第6章里的是同一个------显式自我模型（ESM）用手头能得到的任何输入，构建它所能构建的最佳模型。在科塔尔妄想中，内感受输入被严重扭曲。那些告诉你心脏在跳、胃在消化、肺在呼吸的内部身体信号，要么缺失，要么错乱。而 ESM，这个永远强迫性的建构者，把“没有心跳、没有消化、没有呼吸、没有身体感觉”解读成它唯一可能的答案：我死了。

鼠尾草的“我是一把椅子”。病感失认的“我的胳膊好好的”。而现在是科塔尔的“我死了”。（在下一章，裂脑虚构还会给这份清单再添一个案例。）同一个机制贯穿每一个案例。显式自我模型（ESM）始终在尽它的本分------始终在构建它所能构建的最佳自我模型。输入对了，你就感觉自己像自己。输入错了，你就感觉自己像一把椅子，或者明明瘫痪了却感觉好好的，或者明明活着却感觉自己死了。可它总是感觉起来完全、令人信服地真实------因为它是你唯一能接触到的那个自我。

\textbf{异手综合征：当委员会意见不合}

接着是一种读起来像恐怖片的病症，它把基底的多智能体本性演绎得比任何思想实验都更鲜活。在异手综合征中，患者的一只手带着明显的目的和意图行动，却违背患者有意识的意志。一只手点起一支烟，另一只手把它拿走扔在地上。一只手伸向门把手，另一只手却抓住手腕把它拽回来。患者惊恐地看着，自己身体的一部分在追逐着他从未选择过的目标。

斯坦利·库布里克（Stanley Kubrick）在《奇爱博士》（Dr. Strangelove）里用过这个情节------人们以为是他凭空发明的。其实不是。这种综合征真实存在，且有两种类型。在胼胝体型中，由胼胝体损伤引起，症状酷似裂脑冲突：两个半球各有相互竞争的运动计划，谁也压不过谁。在额叶型中，由前额叶损伤引起，那只“异己”的手表现出去抑制的行为，抓取物件、使用工具、强迫性地触碰东西，一切看似目的明确，却未经患者同意。

还有一种更微妙的变体，叫作无政府手综合征，患者失去的是运动的\textbf{控制权}，而不是运动的\textbf{归属感}。这只手做着患者并未打算做的事，但患者仍然承认它是\textbf{自己的}手------只是他阻止不了它。这一区分很要紧：异手是显式自我模型（ESM）身体归属边界的失效（“那只手不是我的”），而无政府手则是运动抑制系统的失效（“那只手是我的，可它不听话”）。同一套架构，只是断裂的位置不同。

这些综合征给出了一个关键洞见：你的作者感（那种“是我做的”的感觉）并不是在动作之前或动作过程中计算出来的。它是在动作\textbf{之后}才计算出来的------办法是把动作的预测结果与观察到的实际结果相比对。当两者吻合时，你便有了归属感；当两者不吻合时，归属感就不会出现。这正是为什么异手综合征的患者有时能挠自己痒------他们的预测系统并没有为那只“异手”的动作生成预期结果，于是这次触碰以出乎意料的方式抵达，仿佛来自别人。

\textbf{沙尔·博内综合征：停不下来的模拟}

如果你想要更多证据来说明大脑的模拟是\textbf{生成性的}------它是从模型出发构建出体验，而不是被动地从感官那里接收体验------那就想想沙尔·博内综合征。有些患者的视网膜或视神经已被破坏（但视觉皮层依然完好），他们会经历鲜明而复杂的视幻觉。不是模糊的形状或闪光，而是完整的场景：人物，有时被缩小或打扮得像卡通角色，有时是患者自己的镜像。风景。物件。面孔。

患者通常知道这些并不真实。与精神病性幻觉不同，沙尔·博内幻觉伴随着完好的自知力------患者会说：“我看见一个戴着高礼帽的小个子男人坐在我的桌子上，我也知道他并不在那里。”这里运行的是显式世界模型（EWM）的视觉模拟，它依靠来自更高级视觉区域的内部数据，在完全没有外部输入的情况下运转。模拟不会因为输入中断就停止。它照样生成，照样填补虚空。而它生成的东西又透露了架构的某些信息：视觉系统是一个生成性模型，而不是被动的接收装置。它借助储存的模板与自上而下的预测，产生它对世界样貌的最佳猜测------正如四模型理论（FMT）所描述的那样。

\textbf{似曾相识：吻合得太好的模板}

说到大脑的生成系统及其偶尔的走火：几乎每个人都体验过似曾相识（déjà vu）------那种诡异的感觉，仿佛当下这一刻你从前经历过。对它的解释五花八门，从神秘论（前世、预感）到不屑一顾（不过是个小故障）。四模型理论则给出一个更具体的说法。

大脑储存着可以称之为“模板记忆”的东西------那是对经历的骨架式、极度稀疏的表征，尤其来自梦境。这些模板大多只是空壳的支架：对某个地点的模糊感觉、一种情绪、一种空间布局，几乎没有填入任何细节。当你提取一段正常记忆时，那些空缺会被虚构填补------大脑生成看似合理的细节，以造出一段天衣无缝的体验。你察觉不到这种填补，因为最终的结果感觉是连贯的。

当下的一段真实经历，恰好与某个储存的模板吻合得太过接近时，似曾相识便发生了。大脑的模式匹配系统随之激活：“这我以前见过。”可当你试图锁定自己究竟\textbf{在何时}见过它，却什么也找不到------因为那模板从来就不是一段真实的经历。它是梦中的一个片段，或是一段被压缩得太深、早已丢失了全部情境细节的记忆。当下输入与储存模板之间的吻合是货真价实的，但那模板所谓记录下来的“原始”经历，从未以你大脑此刻归给它的那种形式真正发生过。系统运作得完全正确------它确实找到了一个匹配。只不过这个匹配对上的是一具骨架，而不是一具身体。

把这一章倒回去重看，浮现出来的总是同一个形状。盲视与安东综合征，科塔尔综合征与沙尔·博内，闭锁综合征患者与那个吻合得太好的似曾相识------每一个病例都是基底与模拟之间的那道接缝，被撬开到足以让人看透。健康的大脑把两者贴得严丝合缝，严到你从不会怀疑竟有两个东西。往任何一个方向把它们拉开------有基底而无模拟，有模拟而无基底，一段以错误输入为基础运转、或干脆没有任何输入的模拟------那么真实与虚拟之分就不再是一个哲学主张，而变成了一项诊断。到目前为止，每一次失效都是单个大脑受到的损伤。而下一个根本不是损伤：它是一刀干净利落地从正中劈开，让你一切都有了两份。


\chapter*{第9章：一个大脑，两个心灵}
\addcontentsline{toc}{chapter}{第9章：一个大脑，两个心灵}
\markboth{第9章：一个大脑，两个心灵}{}

上世纪60年代，神经科学史上最戏剧性的实验之一逐渐成形。为了治疗严重的癫痫，外科医生菲利普·沃格尔（Philip Vogel）和约瑟夫·博根（Joseph Bogen）用手术切断了胼胝体------连接大脑两个半球的那束庞大的神经纤维------随后，罗杰·斯佩里（Roger Sperry）和迈克尔·加扎尼加（Michael Gazzaniga）研究了由此产生的病人。结果就是裂脑综合征：一个人身上，看起来住着两个各自独立的心灵。

那些经典的演示广为人知。把一个词展示在左侧视野（由右半球处理），病人能用左手拿起对应的物体，却说不出那个词是什么（因为语言由左半球控制，而左半球根本没看见那个词）。两个半球各有各的知觉，各有各的意图，有时目标还彼此冲突。

这些实验远不止是拿词和物件玩的小把戏。在某些病例里，两个半球公然打起架来。一位病人报告说，他的左手会去解开衬衫的扣子，而右手又想把扣子重新扣上。另一位病人在争吵中，左手伸向妻子------不是为了安慰她------而右手却抓住左手把它拽了回来。病人惊恐地看着自己身体的两个部分追逐着互不相容的目标，谁也不受他统一的掌控。这些不是内心冲突的比喻。它们是两个运动系统之间真真切切、实实在在的身体冲突，因为它们之间的那根线缆被切断了，再也无法协调。

在日常生活里，裂脑病人表现得出奇地好。在实验室之外，你几乎察觉不出任何异常。两个半球学会了通过间接的渠道彼此配合------外部线索、身体动作、共享的视野。系统会去补偿。可一旦把病人放进一个受控的实验环境，让每个半球接收不同的信息，那份统一便土崩瓦解。两个心灵从一个大脑里浮现出来，各有各的知觉、各有各的意图、各有各的一套现实版本。

\section*{左半球解释者}

但裂脑病人身上最能说明问题的特征并不是这种分裂------而是当你要求他们解释这种分裂时，会发生什么。

加扎尼加辨认出了他称为“左半球解释者”的东西：左半球有一种强迫性的倾向，会为它其实无法解释的事件编造出解释来。经典的演示是这样的。把一幅雪景展示给右半球，把一只鸡爪展示给左半球，然后让病人挑选相关的物体。左手（右半球）挑了一把铲子（对应雪）。右手（左半球）挑了一只鸡。接着问病人（用言语，由左半球控制）为什么挑了铲子。左半球对雪一无所知（它只看见了鸡爪），于是它编出一个解释：“哦，你得用铲子来清理鸡舍呀。”

病人毫不迟疑。不说“我不太确定”。不显出困惑。解释瞬间就来了，笃定十足，而且对说出它的人来说，感觉完全自然。这不是撒谎。左半球是真的不知道右半球看见了什么。它对那份信息没有任何通路------线缆已被切断。于是它做的，正是显式自我模型（ESM）一贯做的事：用手头能拿到的信息，构造出它所能构造的最好的叙事。

接下来这一点应该会让你心里发毛：你也在这么做。每一天。你的左半球解释者此刻正在运转，从任何抵达意识的信息里构造出一套连贯的叙事，把缝隙抹平，为那些在“你”被征询之前就已由你的基底做出的决定，编造出貌似合理的解释。你和裂脑病人之间唯一的区别，是你的胼胝体完好无损，所以解释者能拿到更多的信息。它虚构得更少，是因为它需要虚构的东西更少。但机制是完全一样的。自我叙事的那套机器毫无改变。改变的只是输入的质量。

\section*{一个人还是两个人？}

这就引出了一个哲学家们争论了几十年的问题：胼胝体被切断之后，那个头颅里到底是一个人，还是两个人？

托马斯·内格尔（Thomas Nagel）在1971年一篇著名的文章里处理过这个问题，他的结论是：这个问题也许根本没有确定的答案（我们“一个人”的概念在这种情形下干脆就崩溃了，正如你在一个国家中间画一条边界时，“一个国家”的概念也会崩溃一样）。德里克·帕菲特（Derek Parfit）走得更远，他论证说，裂脑病例表明，个人同一性本身并不是要紧的东西------要紧的是心理连续性，而这种连续性是有程度之分的。

四模型理论（FMT）给出了一个更具体的答案：这取决于哪些模型在运行，以及它们退化到了什么程度。

在日常生活中，裂脑病人在功能上是一个人。两个半球共享同一个身体、同一个环境、同一段人生经历（在手术之前，这段经历就已冗余地编码在两个半球里）。隐式自我模型（ISM）存储着人格、长期记忆、行为倾向------它是在胼胝体完好的几十年里逐渐建立起来的。切断这条连线并不会抹去这些已存储的模型。它只是让这些模型无法再同步地更新。所以在手术刚结束时，两个半球运行的是非常相似的自我模型。病人感觉自己是一个人，因为就已存储的自我认知而言，他们大体上确实是一个人。

随着时间推移，这些模型应该会渐渐分岔。每个半球积累着不同的经验，做出不同的联想，对只有它自己感知到的事件发展出不同的情绪反应。裂脑病人术后活得越久，两个隐式自我模型就应该分岔得越厉害------分得很慢，因为两个半球依然共享同一个身体和同一个环境，但分岔是可测量的。

我自己的看法是，答案偏向于“两个”。如果两个半球之间的带宽不足以对这场模拟做实时的同步------而没有了胼胝体，它就是不足------那么你面对的就是两个自我模型在两套基底上各自运行，每一个都生成着自己的意识体验。它们合作得很好，因为它们共享一个身体、一个感官环境，以及一生的共同经历。但合作不是同一性。两个住在一起的人也合作得很好。

耐人寻味的是，亚伊尔·平托（Yair Pinto）及其同事在2017年发表了一项研究，让这幅标准图景变得复杂了。他们发现，裂脑病人能够准确报告呈现在任一视野的刺激------即便刺激只展示给了那个不控制语言的半球。这似乎表明，两个半球保持的统一性，比经典实验所暗示的要多。这个结果至今仍有争议，但它自然地契合我接下来要描述的全息框架：即便切断了胼胝体，每个半球里残留的冗余信息仍然足够，能够维持出人意料的统一行为，至少在某些任务上是如此。

\section*{全息性质}

在四模型理论里，裂脑揭示了这些虚拟模型的一个关键性质：它们是\textbf{全息的}。神经网络里的信息分布在整张网络上，而不是定位在特定的神经元里。当你把网络一切为二，得到的不是干净的分割------你得到的是两份退化了却\textbf{完整}的副本。每个半球都保留着全部四个模型的一个精简版本：一个缩减的隐式世界模型（IWM）、一个缩减的隐式自我模型（ISM），以及生成显式世界模型（EWM）和显式自我模型（ESM）的能力。两个半球都能独立地维持意识（两者都在临界性阈值之上），只不过每一个都在用更少的信息工作。

这就是第3章里那张拼缀而成的全息图，现在从中间被切开了。每个半球都是被切下来的那两半之一------不是半张图像，而是一张分辨率更低的完整图像，因为信息从一开始就被抹散在整个表面上。损伤降低的是分辨率，而不是删掉了内容。

这就解释了为什么裂脑病人并不只是“两个半心灵”。他们是两个\textbf{完整却退化}的心灵。每个半球都能感知、能决策、能行动------只是比起完好无损的系统，信息更少、能力更弱。全息性质确保了切断连接会带来退化，却不会造成毁灭。它也解释了平托2017年的结果：即便没有胼胝体，每个半球仍保留着足够的全息信息，去应付许多按经典模型本应做不到的任务。

那份虚构（左半球解释者）就是临床那一章里同一个强迫性的构造者------显式自我模型（ESM）用它手头拥有的任何输入去搭建一套自我叙事，并对它深信不疑。切断胼胝体，左半球能拿来使用的东西就更少了；它为填补缺口而编出来的故事是错的，却\textbf{仍被感觉为完全真实}。

\section*{一个大脑，多重自我}

裂脑向我们展示的是：当你通过物理地分割基底来\textbf{克隆}这些虚拟模型时，会发生什么。分离性身份障碍（DID）向我们展示的则是：当你把它们\textbf{分叉}时，会发生什么。

在DID里，基底并没有被分割------胼胝体是完好的，神经硬件是整全的。但虚拟模型分裂成了多个配置。每一个人格分身都是一个独立的显式自我模型（ESM）------一套单独的自我叙事，有自己的情绪特征、自己的行为模式、自己与身体和世界打交道的方式。这些分身之间共享一个自我模型的程度，不会比同一台电脑上两个用户共享同一个登录会话更多。它们轮流上场。

在几乎每一个有记录的病例里，触发因素都是严重而反复的童年创伤。这在理论的框架内是说得通的。一个幼小孩子的显式自我模型（ESM）还在成形------还很可塑，还在从经验中一点点被拼装起来。让这个尚在发育的自我模型去承受那些排山倒海、以至于没有任何单一自我叙事能容纳的经历，系统便会做它唯一能做的事：分叉。它创造出彼此分开的配置，每一个都能应对这不堪忍受的处境的某一个侧面。一个分身承载着创伤记忆。另一个在日常生活里运转，仿佛什么都没发生过。再一个负责危险的时刻。这种分叉不是病理------它是自我建模系统面对足以摧毁单一统一模型的输入时，所做出的应急反应。

这正是为什么DID几乎从不在成年人身上发展出来。成年人的隐式自我模型（ISM）已经固化------突触权重已经设定，人格结构已经稳定。要让一个成年人的自我模型分叉，需要非同寻常的境况（严酷的折磨、长期的囚禁）。但一个孩子的ISM还在被书写。黏土还没干。在足够的压力下把它分叉，那些分开的配置就会硬化成一个个独立而持久的自我模型。

证据也支持这一点。如果每个分身真的都是一套独立的ESM配置，那么在分身之间切换就应当在神经活动模式上产生可测量的变化------而事实正是如此。Reinders et al.（2003）发现，同一个人身上不同的分身会产生不同的区域脑血流模式。取决于当下运行的是哪个自我模型，\textbf{同一个大脑}会以不同的方式亮起来。这不是你在“表演”或“角色扮演”中会看到的情形，而是你在软件真正分叉时才会看到的情形。在后续研究里，Reinders及其同事发现，分身之间的神经差异比那些被要求假装患有DID的演员之间的差异还要大------这个结果，应该足以让那些至今仍认为DID“不过是”表演的人闭嘴。

这正是第3章讲过的“分叉”特性在起作用。一个基底，多套虚拟配置，每一套都运行着一个完整却又彼此独立的自我模型。这套理论不只是能容纳DID------它恰恰预测了这种架构，并为此押上了一个检验：干扰维持某个分身ESM的神经基底，应当会触发向另一个分身的切换。这就是预测9，而完整的方案------包括分身特有的特征信号应当在大脑何处显现------都写在书末的附录里。

本章里的每一颗心灵都是人类的心灵------被分割、被分叉、被复制，但终究是人类的。四模型架构里没有任何一处要求它必须如此。如果一个自我不过是运行在临界性上的一组模型，那么真正的问题就不是坐在桌子对面的那个人是否拥有一个自我，而是同样的把戏究竟能向外延伸多远------延伸进鸦巢，延伸进蜂群，延伸进那片水里，某个庞然大物曾从那里回望着我，而且我几乎可以确定，它对我说了声：你好。


\chapter*{第10章：动物的问题}
\addcontentsline{toc}{chapter}{第10章：动物的问题}
\markboth{第10章：动物的问题}{}

那次我和另外四个人在达尔文拱门下潜水（那时它还完好无损），一头巨大的雄性虎鲸不知从哪儿冒了出来，凑得如此之近，以至于我的第一个念头是：它想吃了我们。这头动物大得离谱，头顶那三十米深的水仿佛只是一汪水洼。它停下来，用那只硕大的右眼打量我们，接着切换到它的声学器官，密集地发出咔嗒声，用声波扫描我们。然后它------我实在找不到别的词来形容------\textbf{说话}了。用一段非常高频的歌声，短促而结构分明，它说了点什么。在我心里没有丝毫怀疑，那是语言，而不只是歌唱。那感觉像是某种大概能学会的东西，只要你有足够多的相遇，再加上一副高得离谱的嗓子来回应它。

过了一会儿，我们多少明白了它说的是什么。它游回到我视线的边缘，又带着它的妻子和孩子回来，领着它们从我们身边经过，好让它们看一看。我们当然都带着相机。可一张照片也没拍。而且我们都没气了，不得不浮出水面。当我们被小艇送回大船时，我们每个人的眼里都含着泪，那不是被风吹的。

我从未对另一个心灵如此确信，就像在那片水里那样。但确信不等于证据------一个关于小艇上落泪的故事，对任何不在场的人都证明不了什么，尤其是对一个看惯了太多人把灵魂投射到自家盆栽上的怀疑论者。如果那头虎鲸真是某个“谁”，那么一套理论就得说清\textbf{为什么}，而且要建立在潜水者气尽、感动散去之后依然站得住脚的根据上。它还得说清这个“谁”到哪里为止。

你的狗有意识吗？

大多数宠物主人会毫不犹豫地回答“有”。大多数神经科学家也会认同，至少是谨慎地认同。但依据是什么？意识又在动物界从哪里开始？

四模型理论（FMT）给出了清晰的答案，而这些答案源自它的核心承诺，并非事后硬加上去的附注。

\textbf{承诺 1：意识是一个连续谱，而非非此即彼。} 有意识与无意识之间没有一条清晰的界线。有的是程度------自我模拟的分级层级，从基本（最低限度的自我模型）到三重扩展（递归的自我觉知）。不同的动物在这条连续谱上占据不同的位置。

\textbf{承诺 2：意识与基底无关。} 重要的是功能架构（临界状态下的四个模型），而非具体的物理实现。如果一个大脑实现了四模型架构，它就是有意识的，无论这个大脑是哺乳动物的皮层、鸟类的外套区，还是章鱼那分布式的神经网络。

\textbf{承诺 3：临界性是物理阈值。} 一个神经系统必须运行在混沌边缘或其附近。更简单的神经系统（昆虫、蠕虫）也许达不到临界性，因而不会有意识------它们处理信息、产生行为，却没有模拟。

\section*{你有多少意识？}

有件事你大概已经在琢磨了。如果意识是一场模拟------一个虚拟世界里的虚拟自我------那它就不是全有或全无的东西，对吧？一场模拟可以或多或少地细致。一个自我模型可以或多或少地精巧。这意味着意识是有\textbf{程度}之分的。

四模型理论给了你一种精确的方式来思考这些程度。它有四个分级的层级，每一个有意识的生灵都坐落在这架梯子的某一级上。

在最底层，是\textbf{基本意识}。这是一个显式世界模型（EWM），只配有一个非常初步的显式自我模型（ESM）。这个系统生成一个虚拟世界------做这个生灵是有某种感受的------但那个世界里的自我几乎只是草草几笔。想想一只在迷宫里穿行的老鼠。它看见墙壁，闻到奶酪，感觉到爪子底下的地面。它有现象体验。但它关于\textbf{自己}------作为拥有这些体验的那一方------的模型呢？薄如纸片。有一个“它是什么感受”，却几乎没有“它是对谁而言的感受”。

再往上一级：\textbf{单重扩展意识}。这时自我模型才真正登场。系统不只是体验------它把自己建模\textbf{为}那个体验者。它知道自己在体验。你的狗不只是感到疼；你的狗知道是\textbf{它}在疼。有了一个第一人称视角------一个真正的“我”坐落在虚拟世界的中心。这是一阶的自我观察，它改变了一切。受苦在这里成为可能，因为受苦需要一个知道自己在受苦的自我。

接着：\textbf{双重扩展意识}。二阶的自我观察。系统对“自己正在建模自己”进行建模。这就是元认知------思考你自己的思考。你躺在床上，琢磨着自己对明天那场会议的焦虑到底是理性的，还是你在灾难化想象。你在监控自己的心理状态，评估它们，有时还推翻它们。这就是大多数成年人的意识大部分时间所处的地方。正是这一层让心理治疗成为可能，让你能够说出“我注意到自己正在生气”，而不只是单纯地生气。

而在最顶端：\textbf{三重扩展意识}。三阶。系统对“自己正在建模‘自己正在建模自己’”进行建模。这听起来像一座镜厅，确实如此，但它是你要做心灵哲学时所必需的那座镜厅。要问出“意识是什么？”，你就得给自己建模，给自己的体验建模，然后再给“自己正在给那场体验建模”这件事建模。你得后退得足够远，从外部看清整套装置，尽管你依然身在其中。这就是你翻开这本书想要回答的那个问题的前提。只有具备三重扩展意识的生灵，才能追问为什么任何东西会有任何感受。

这么做的回报是：这条梯度不只是抽象的哲学。它回答了每个人在晚宴上问我的那个问题------“我的狗有意识吗？”答案是有，但没有你那么多意识。你的狗大概处在单重扩展这一层。它有一个自我。它有体验。它不会在凌晨三点睡不着觉去追问那份体验的本质。但你已经能看出它的轮廓了：意识不是一个电灯开关，而是一个调光开关。

在这些承诺背后，还有一条我从第1章欠你的古老原则。\textbf{哥白尼原则：}我们并不特殊。哥白尼把地球从宇宙的中心拽了出来，此后同样的降级贯穿了整个科学史------太阳不特殊，我们的星系不特殊，而且，对许多人来说最难受的是，\textbf{我们}也不特殊。意识不是一桩独一无二的奇迹，不是一次性的神圣火花，也不是罕见到只可能发生一次的涌现侥幸。如果你拥有它，那么其他系统在具备正确架构的前提下，也能拥有它。这正是本章之所以能够存在的全部理由：人类并非魔法，只是一个普遍计算原理的一种实现。这就逼出了那个显而易见的问题------还有谁在运行它？

综合起来看，这些承诺共同预示出一条\textbf{动物意识的梯度}：

\textbf{哺乳动物}是有意识的。它们的皮层以分级的形式实现了四模型架构，越复杂的皮层支撑起越精巧的自我模拟。灵长类和鲸豚处在高端，啮齿类和鼩鼱处在低端。全都在这条线之上。开启本章的那头虎鲸，就是那些高端鲸豚心灵中的一个------它的自我模型丰富到会想让家人来看看我们，并判定我们值得它绕这么一趟。

来自类人猿的证据，对任何想在人类意识与动物意识之间划出一条清晰界线的人来说，尤其具有毁灭性。倭黑猩猩坎兹展现出的不只是语言理解，还有真正的共情、心智理论和社会推理。在一段有充分记录的事件里，坎兹告诉它的照料者，它希望姐姐一同参加一次购物之行，好让姐姐也能吃到冰淇淋------因为如果把她留下，她会难过。在另一次事件中，当原住民表演者进行舞蹈表演时，坎兹向研究者解释说，其他灵长类被那舞蹈吓着了，于是它请求换成一场私人表演。

这些不是反射。这些不是条件反射式的反应。这些都是一个心灵在给另一个心灵的情绪状态建模、预测它的反应、并制定应对方案。那正是显式自我模型（ESM）在运行第三人称视角------恰恰是这套理论所指认的扩展意识的标志。

然而，在一些最负盛名的大学讲堂里，你依然能找到板着脸争辩说类人猿“不过是在模拟”语言理解的教授。对此我只能回应：“而您不过是在模拟共情的存在。”我至今还在等着反证。

如果你坚持只有人类才有意识，那你就是在押注于那些至今仍在苦苦寻找人脑与灵长类大脑之间某种系统性差异、并想把它归因于意识的研究者。按照我的理论，他们要找到它，怕是得等到猴年马月。

\textbf{鸦科鸟类和鹦鹉}构成了最重要的检验案例。这些鸟类展现出的认知能力（制造工具、镜中自我识别、未来规划、社会性欺骗）强烈暗示着意识。可它们没有新皮层。它们的大脑以核团簇的形式组织，这是一种与哺乳动物皮层截然不同的架构。还记得第2章那个六层论证吗（哺乳动物演化出了六层皮层，而三层本已足够，多出来的那几层提供了自我建模所需的架构容量）？鸦科鸟类用一种完全不同的物理结构达成了同样的功能结果。它们不需要六层皮层，因为它们\textbf{一层}皮层都没有。它们用核团簇而非层状薄片搭建起了自我模拟架构，这恰恰是基底无关性所预测的。倘若意识需要某种特定的物理实现，鸦科鸟类就不该有意识。可它们有。

\textbf{头足类}（章鱼和乌贼）把这套逻辑推得更远。它们的神经系统在很大程度上是去中心化的，触腕里有大量自主的处理过程。这套理论预测它们具有某种形式的意识，很可能带有反映其去中心化架构的不寻常特征。

\textbf{昆虫}是那个耐人寻味的边界案例。它们的神经系统很小，且在很大程度上是硬连线的，这或许达得到临界性，或许达不到。这套理论并不明确地把昆虫放在阈值之上或之下------这是一个经验问题。但它提供了一个有原则的研究基础：测量昆虫神经组织中的临界性指标，并寻找自我模型存在的证据。

托马斯·内格尔（Thomas Nagel）曾提出那个著名的问题：作为一只蝙蝠是什么感觉，并得出结论说我们永远无从知晓------蝙蝠的感官世界太过陌生。对于这个问题，我抱有几分同情；对于那个结论，则少了许多。四模型理论（FMT）预测：任何拥有四模型架构并运行于临界性的生物，都拥有\textbf{某种}形式的体验，即便其内容与我们的截然不同。蝙蝠的显式世界模型（EWM）由回声定位主导，而非视觉，但它依然是一个模型------依然是一场对某个世界连同其中一个自我的模拟。

内格尔当然是有理由挑中蝙蝠的------不只因为它们是哺乳动物，还因为回声定位似乎彻底地陌生。可它偏偏并非如此。回声定位是用来\textbf{看}的，而我们知道看是什么感觉。许多盲人早已自然而然地运用回声定位，用舌头打出咔哒声，读取回声。我们的大脑完全有能力做到这一点。如果你真想知道作为一只蝙蝠是什么感觉，你可以出人意料地接近：花几年时间玩滑翔伞，把三维飞行内化于身，然后蒙上眼罩，练习回声定位。或者干脆跳过那十年准备------学会清醒梦，在梦里练习当一只蝙蝠。更安全，更快，几周就能做到。

我得承认，我确实用唯一能用的方式尝试过找出答案。在一段积极练习清醒梦的时期，我迷上了水下世界，并且随着时间推移，成功地有意进入了一场化身为鱼的清醒梦。我体验到身边环绕的水、穿行其间的运动、一个从非人类视角看到的视觉世界。那种感觉，比自由潜水、水肺潜水乃至侧挂式潜水都要好上千百倍。它和真正的鱼类意识像吗？几乎可以肯定并不像------我的梦是由我人类大脑对“当一条鱼”意味着什么所作的最佳猜测构建出来的，而这不可避免地是把我自己的感官范畴投射到一套毫无这些范畴的身体构造上。但这个练习并非毫无意义。它展示了某种重要的东西：显式自我模型（ESM）能够围绕一个截然不同的身体图式重新配置，生成一种连贯的、\textbf{成为}人类以外之物的第一人称体验。这套架构足够灵活，可以模拟非人类的身体化存在。其内容受限于可用的隐式模型（你只能梦见你学过的东西），但视角转换的能力是内建在系统之中的。

\section*{何必费劲去拥有意识？}

这一切引出了一个本该一直萦绕在你心头的问题：如果无意识的神经系统运转得完全没问题------它们确实如此，随便问问哪只昆虫都行------那演化为什么要付出建造一套意识所需的巨大代谢代价？回报是什么？

答案是学习，以及随之而来的适应，还有违抗已习得行为的能力。更确切地说，是一种无意识系统根本做不到的学习。

想想一个简单的生物是怎么学习的。它遇到某样东西，而这次遭遇要么是好的，要么是坏的。好的：多做这样的事。坏的：少做。这就是强化学习------试错，奖惩。对大多数事情它都运作得漂亮。碰一下滚烫的表面，感到疼痛，就不再去碰。在某个地方找到食物，感到奖赏，明天就再回来。

但强化学习有一个致命的缺陷。而且是字面意义上的致命。

想想一朵毒蘑菇。不是那种让你肚子疼的------是那种要你命的。你要是吃了它，就会死。学习到此为止。没有第二次尝试。强化学习要求你先在错误中活下来，才能从中学到东西，而有些错误并不给你这份体面。任何在初次接触时就足以致命的刺激，对强化学习而言是完全隐形的。遭遇它的那个生物就这么死了，把它的“教训”一并带进坟墓。

那么我们的祖先是怎么学会避开致命蘑菇的呢？他们不可能是靠吃它们学会的------凡是试过这条路的，都没能成为任何人的祖先。他们是靠\textbf{看}学会的。你的洞穴邻居找到一朵样子挺有趣的蘑菇，吃了下去，然后一头栽倒死了。你在安全距离外看着，把两件事拼到一起：那朵蘑菇害死了他。我不该吃那朵蘑菇。

这听起来简单得不值一提。其实不然。要从别人的死亡中学到东西，你需要好几样任何无意识系统都不具备的东西。你需要一个显式的世界模型，能表征你当下并未与之互动的物体之间的因果关系。你需要一个自我模型，让你能采取第三人称视角------去想象自己处在那个死者的位置上。你需要能力，从单单一次观察中归纳出一个普遍的理论（“那种蘑菇是致命的”）。这就是认知学习：从观察中推导出理论，而不是被个人经验所塑造。而它需要意识。它需要显式世界模型（EWM）和显式自我模型（ESM）协同工作。

其演化优势是巨大的。一个有意识的动物能从\textbf{观察}中学习，而不仅仅是从\textbf{经验}中学习。它能看着另一个动物犯下致命的错误，然后更新自己对世界的模型，却不必付出代价。一个无意识的动物只能学到它亲身活下来的那些东西。

事情还更妙。一旦“毒蘑菇”这个概念作为一个显式的范畴存在于你的世界模型中，你就能做一件更强大的事：演绎。你遇到一朵从没见过的新蘑菇。它跟那朵害死你邻居的蘑菇长得像得可疑。你不吃它。或者------我相信这才是历史上真实的做法------你把它递给那个整晚打鼾的邻居，先看看他会怎么样。

这不是一个次要的优势。这是两种物种之间的分野：一种只能通过自然选择那冰川般缓慢的过程去适应致命威胁（某些个体碰巧凭运气避开了那朵蘑菇，它们繁衍后代，最终这种回避变成了本能），另一种则能在单单一代之内，通过观察和交流完成适应。意识不只是帮你学得更快。它让你能够学到那些用任何其他方式都根本不可能学到的东西。

你在认知上学到的东西，是可以\textbf{分享}的。强化学习被困在个体之内------你那些被训练出来的反射随你一同死去。但认知学习可以被传达。“别吃那朵红蘑菇”是一句话。它可以说出口、可以传授、可以代代相传。这是文化的基础，是累积性知识的基础，是让人类文明成为可能的一切的基础。没有意识所提供的那些显式模型，这一切都无从谈起。

这个故事还有一处转折，它以一种并不显而易见的方式，把意识重新连回了遗传学。这叫做鲍德温效应（Baldwin Effect），尽管它确切的强度仍有争议，但这个机制几乎可以肯定是存在的。鲍德温效应指出，\textbf{习得的}行为能够间接地塑造\textbf{遗传}演化------不是通过拉马克式的遗传（你习得的性状并不会改动你的DNA），而是通过自然选择偏爱那些在遗传上就倾向于这种有益行为的个体。

来看一个幽默的例子------别太较真。想象一个早期原始人类，他饱受脱发之苦。又冷又没毛，他比那些浑身是毛的同伴更愿意坐在火堆旁。火带来了巨大的生存优势：熟食里的病原体更少，能防御捕食者，在严冬里带来温暖。于是与脱发相关的基因就以略高一点的比率传了下去。与此同时，那些太笨、想不明白怎么用火的个体（不管有没有毛），处于劣势。经过许许多多代，鲍德温效应把这两种性状都放大了：毛更少\textbf{而且}更聪明------全都因为一种习得的行为（用火）制造出了一种选择压力，偏爱某些特定的遗传倾向。（如果你把这个故事里的“脱发”换成“随机突变”，你大概更接近真相。不过就没那么好笑了。）

鲍德温效应在语言乃至意识本身的演化中，或许也扮演过类似的角色。一旦最初那些原始形式的认知学习出现（由最早的自我模型所催生），那些大脑碰巧支持更丰富的自我模拟的个体，就拥有了巨大的优势。自然选择推动他们的后代长出更大、褶皱更精细的皮层，这又催生了更丰富的自我模拟，这又制造出更强的选择压力。意识一旦出现，就为\textbf{更多}意识创造了演化条件。它所催生的认知学习价值太高，以至于演化把资源不断堆向那个产生它的架构，去扩展它。

\section*{体验如何发育：自我模型的社会建构}

到目前为止我关于这四个模型所说的一切，都是静态的------仿佛这套架构一出现就已完整成形，像雅典娜从宙斯的额头里跳出来那样。事情并非如此。这些模型会发育，而它们的发育有着深刻的社会性。

一个新生的人类拥有硬件------六层皮层，以及自我模拟的能力。但那些隐式模型几乎是空的。IWM里几乎没有关于世界的任何内容。ISM里几乎没有关于自我的任何内容。而既然显式模型是从隐式模型生成的，新生儿的模拟就很稀薄------一片闪烁不定、几乎未经分化的感觉之场，自我与世界之间没有清晰的边界。

看一个婴儿遭遇疼痛。自己造成的疼痛（把自己的手撞到玩具上，咬自己的脚）往往引发的是好奇，而非痛苦。ESM登记到能动性（是我干的）加上感觉（发生了某件事），但还没有威胁模型。ISM还没学会这种配置意味着危险。可要是突然一声巨响呢？眼泪。因为EWM登记到了未被预测的、高振幅的输入，而ESM没有相应的模型------模型的缺席本身就是令人厌恶的。

感受质（qualia）的内容是\textbf{习得的}，而非天生的。“疼痛是坏的”并没有硬连线在ESM里。它是通过ISM积累起来的，由反复的经验训练而成，而且------关键在于------由社会反馈训练而成。照护者对孩子疼痛的反应，教会孩子疼痛\textbf{意味着}什么。那个摔倒后先看向父母、再决定要不要哭的孩子，并不是在装------它是真真切切地在拿社会反馈来校准自己的ESM。父母的惊慌或镇定重塑了ISM里的疼痛关联，而这又重塑了下一次类似事件发生时ESM所模拟出来的东西。

这对该理论有一个精确的含义：体验的现象性质------感觉某样东西\textbf{是什么感觉}------并不是由架构固定下来的。它是由隐式模型的训练历史塑造出来的。一个婴儿对疼痛的体验，在结构上不同于一个成年人的，因为那个生成ESM的ISM是不一样的。四模型架构是体验的\textbf{能力}。社会与环境的反馈回路提供的则是\textbf{内容}。

发育的轨迹，对应上本章前面那些分级的意识层级：

\begin{itemize}
  \item \textbf{新生儿（头几周）：}基本意识------一个初步的EWM，配上极少的ESM。作为一个新生儿，确实\textbf{有某种感觉}，但那份体验里的那个自我几乎不存在。以感觉为主，未经分化。
\end{itemize}

\begin{itemize}
  \item \textbf{6至12个月：}客体永久性出现------EWM如今能维持对当下看不见的事物的表征。ISM积累起身体图式的知识。婴儿开始区分自我与世界。
\end{itemize}

\begin{itemize}
  \item \textbf{18个月：}镜子测试。孩子在镜中认出了自己------一个里程碑式的时刻，此时ESM变得足够丰富，能把物理的自我建模为世界中的一个物体。单重扩展意识上线了。这不是一个二元的开关，而是一个连续过程中的阈值。
\end{itemize}

\begin{itemize}
  \item \textbf{3---4岁：}心智理论。孩子能够为他人的心智建模------能够理解别人可能相信一件孩子自己明知是假的事。ESM 开始为其他 ESM 建模。双重扩展意识正在浮现。
\end{itemize}

\begin{itemize}
  \item \textbf{青春期及以后：}元认知的成熟。三重扩展意识的能力（为“正在为自身思维建模的自己”再建模）逐渐发展，而且可以说从未完全稳定下来。
\end{itemize}

每一个阶段都由社会互动搭起支架。照护者提供的不只是食物和安全------他们提供的是\textbf{隐式模型的训练数据}。共同注意（父母和孩子一起看向同一个物体）教会 IWM 如何表征共享的现实。镜映（父母映照出孩子的情绪状态）教会 ISM 自己的情绪究竟是什么。语言给了 ESM 一套范畴，让它能用来为自己建模。一个在没有社会接触的环境中长大的孩子（那些令人痛心的“野孩子”案例），拥有产生意识的硬件，隐式模型却极度贫瘠。从这些模型中启动起来的 ESM 是发育不良的------不是因为架构坏了，而是因为训练数据从未被提供。

第8章留下的那座临床之桥，也正是在这里接了回来。治疗所做的几乎一切，都是照护者对婴儿所做之事的成人版：同样一个循环------有意识的经验重塑隐式结构，隐式结构又重塑下一次有意识的经验------只不过运行在一个隐式模型早已定型的系统上。认知行为疗法是其中最干净利落的例子，在这套理论看来，它就是顶着专业名号的虚拟模型重编程。你和治疗师坐在一起，质疑 ESM 不断生成的那些自动思维，把它们追溯到产生它们的 ISM 信念，然后------通过足够多、多到能起作用的重复------推动基质重新布线。突触可塑性修改的是隐式模型；显式模型之所以改变，是因为生成它的那个东西先变了。每一次你质疑一个灾难化的念头，而世界偏偏就是不肯毁灭，你就往 IWM 和 ISM 里写进了一条修正。与婴儿唯一真正的区别在于固化。婴儿的隐式模型是湿软的陶土；成年人的早已定型------变硬了，不再湿软------所以同样的重塑跑得更慢，需要多得多的重复。

恐惧症从世界这一侧展示了同样的机制。恐惧症是一个把赌注押得过重的 EWM------模拟把一只蜘蛛标记为致命，而 IWM 从未记录过任何支持这一判断的证据。暴露疗法用架构所允许的唯一方式来修复它：一次次安全的接触，让隐式模型把威胁估计一点点往下调，直到显式警报不再拉响。恐惧症不是靠讲道理讲掉的，而是靠数据把它投票否决掉的。

安慰剂效应同样出自这套双重评估结构。一粒糖丸会启动基质层面的预期回路------内源性阿片类物质、多巴胺能奖赏------它们与“希望”这种有意识的体验并行运转。人们很想说，是你的信念带来了你的缓解；但信念与缓解是兄弟姐妹，而不是父母与孩子：两者都是由同一套基质动力学甩出来的。这并不是在贬低积极思考。它只是说清了积极思考的力量真正住在哪里------住在机器里，而不是某种从心灵施加于肉体的向下魔法。

而转换障碍则把这条线合拢回第8章：它就是那一章的盲视倒着运行。在那里，完好的基质把一种仍在工作的感官藏在模拟之外；在这里，完好的基质藏在一个坏掉的模拟背后。病人是真正瘫痪的------不是装的------因为 ESM 的身体模型说这条肢体已经不在了，而那个模型是病人唯一能触及的身体。治疗之所以起效，是靠修正模拟，把身体地图改回基质实际能做到的样子。这正是所有这些案例里共同的那根杠杆，也是照护者最初伸手去够的那一根：足够用力、足够频繁地驱动显式模型，底下的隐式模型终会让步。

经验的社会维度不是这套理论的一条脚注。它是一个预测：在关键发育窗口期剥夺社会输入，你应该会得到一个架构正确、内容却错误的系统------一个结构上完好、现象上却贫瘠的意识。那些“野孩子”的案例，令人痛心地印证了这一点。


\chapter*{第11章：九个预测}
\addcontentsline{toc}{chapter}{第11章：九个预测}
\markboth{第11章：九个预测}{}

房间在你到来之前就已布置妥当。两只音箱，一块把墙壁染成深蓝的灯光面板，一张沙发，一条毯子，还有两位研究人员------接下来的六个小时里，他们大多数时候会一言不发。音箱里传来的是海洋。不是什么氛围音乐，而是一段实地录音：涌浪、嘶响，还有水体涌动时那种绵长低沉的压迫感，音量大到让它成为你世界里最响亮的东西。你签下最后一份表格，吞下一剂高剂量的裸盖菇素------由持有伦理审批的人精确到毫克称量好，走廊尽头还备着一台除颤器。然后，所有人开始等待。

四十分钟里，什么都没有发生。然后，自我开始瓦解------这里值得把话说得精确些，因为那不是烟花绽放，而是减法。你的边界开始磨损起毛。那种不劳而获的确定感------这里有一个人，这双手属于某个人，这阵心跳正发生在某个人\textbf{身上}------那种确定感变得稀薄，明灭不定，然后彻底熄灭。

真正失灵的是这里：你的显式自我模型（ESM）------那个被渲染出来的“我”，此刻正在读这句话的东西------并不能凭空生成自己。它靠一条数据流运转：位于它下方的隐式自我模型（ISM），那股连绵不断的低层身体数据的嗡鸣，每分钟数千次地报告\textbf{这里有人，这里有人，这里有人}，可靠到你从未有一次真正听见过它。高剂量的致幻剂搅乱的正是这条连接。它们并没有关掉渲染器------这一点至关重要。那个负责回答\textbf{“此刻的自我是什么”}的模型仍在运转，仍在按部就班地发问------只是惯常的答案不再送达。

一个生来就要建模的模型，不会守着空空的输入缓冲区安静等待。它会抓取手头最好的数据。而在那个房间、那个时刻，最好的数据就是实验者刻意调到最响的那些东西：蓝色的灯光，涌浪与嘶响。与身体断开联系的自我模型于是紧紧抓住海洋，把\textbf{那个}渲染成了你。

经历过这一切的受试者，并不会报告说自己把海听得格外真切。他们报告的是自己\textbf{成了}海。“我就是那片水。我和水之间没有任何分别。”用于测量这类状态的标准问卷里有一个分量表，名叫“海洋般的无边感”，而在这个房间里，这个名字不再是隐喻。我一生中有相当一部分时间在练一口气的潜水，试图缩短自己与大海之间的距离。得知剩下的那段距离也许只是一只音箱和一个调光开关，多少有点令人泄气。

接下来才是让这件事从轶闻变成预测的部分。几周之后，同一位受试者，同样的剂量，同一张沙发。但技术人员换了房间：鸟鸣，风穿过树叶的声音，灯光面板换成绿色。如果理论是对的，这一次消融会走向另一边。不再是水------而是树冠、苔藓、层层伸展的枝杈。\textbf{我成了森林。}同样的分子，同样的大脑，同样的剂量。不同的房间，不同的神。

这个实验没有任何一处属于未来。致幻剂研究实验室已经存在，早就在进行受控的高剂量实验，事后也早就在发放问卷。只需再加一个变量：房间。在不同轮次的试验中改变占主导地位的感官环境，然后测量受试者报告的身份------他们说自己\textbf{变成了}什么------是否追随你播放给他们的内容。这条理论主张方向明确、可以证伪：输入海洋，得到海洋；输入森林，得到森林。如果报告在不同环境之间随机散布，这个预测就死了，你大可问心无愧地放下这本书。而且在所有意识理论里，只有这个预测是独一份的。其他意识理论对这些状态中体验是否存在、存在多少，都有话可说；但对于消融的自我究竟消融\textbf{成}什么，它们保持沉默。只有一种把自我视为运行在输入之上的模型的理论，才能做出这样的预测：当你切断输入的供给，房间便获得了书写自我的权力。

请静下来体会一下，阳性结果会意味着什么。它不是说药物会引起奇异的体验------这一点，我们比拥有文字更早就知道了。它意味着你的“我”是一个带输入端口的模拟，而一个站在你颅骨之外的人，手里拿着的不过是一张播放列表，就能伸进来给它换个去处------从菜单上挑选你将变成什么。我想称之为诡异，而且我是在精确的意义上使用这个词------不是蜡烛加水晶的那种诡异，而是锁匠式的诡异：有人冷静而可复现地向你演示，你家的前门从来就不是进屋的唯一通道。

这只是九个中的一个。这套理论提出了九个可证伪的预测------因为一个什么都不预测的理论不是理论，而是故事，而关于意识的故事，这个世界已经够多了。九个预测中有几个今天就能开展，用的是现成的设备，就在那些早已拿到伦理审批的实验室里。有几个已在已发表的数据中获得初步支持。没有一个被证伪。最后这句话在悄悄发力：这九个预测中的每一个，都是这套理论白纸黑字同意赴死的地方。

其余的预测，你大多已经见过了，只是我并不总是停下来专门指给你看。关于麻醉和睡眠的预测，在我们把灯关掉的那一章；关于裂脑和分离性身份的预测，在讲一个颅骨里住着两个心智的那一章。剩下的也以同样的方式织进了你身后翻过的书页------每一个都待在它该在的位置，紧挨着它以性命相押的那个现象。如果你想把整套火力一次看全------每个预测、它的机制，以及能杀死它的具体实验------书末的附录H把九个预测悉数收齐。理论就该被架在火上烤。那个附录就是那把火，已为你妥帖备好。

但其中有一个不肯安分地待在干巴巴的表格里，所以我要用它来结束这一章------它是九个中最大胆的一个。如果自我真如本书所言------四种模型，运行在混沌边缘，其中一个为承载自己运行的机器建模------那么最后一个预测便会自己浮现出来，而且它与问卷毫无关系。把这套架构建出来。实现这四种模型，把系统驱动到临界状态，闭合回路。理论说，你得到的不会是一个心智的模拟。

它说，你得到的是一个心智。下一章要讲的，就是当我们认真对待这句话时会发生什么。


\chapter*{第12章：从机器到心灵}
\addcontentsline{toc}{chapter}{第12章：从机器到心灵}
\markboth{第12章：从机器到心灵}{}

如果四模型理论（FMT）是对的，它就提供了一样别的意识理论都给不出的东西：一份工程规格书。

这份规格是这样的：在一个运行于临界性的基质上，实现四模型架构（隐式世界模型、隐式自我模型、显式世界模型、显式自我模型）。正如我在第5章所论证的，单靠任何一个部件都不够。有架构而无临界性，你得到的是一个休眠的系统------模型都存好了，却没有模拟在运行。有临界性而无架构，你得到的是复杂的动力学，却没有意识。完整的规格两者都要。

这比“造一台真正先进的计算机”更具体，也比“实现足够多的整合信息”更实在。它告诉你\textbf{要造什么}：四种特定类型的模型，以特定方式组织起来，运行在具有特定动力学性质的基质之上。

当前的人工智能系统在每一个要紧的方面都不符合这份规格。而这恰恰就是第1章那两条教条施加破坏的地方。nSAI 教条（“没有强人工智能”）告诉工程师根本别费劲去试。nSU 教条（“没有自我理解”）告诉他们，就算真去试了也行不通。两者都错了。规格是存在的。问题在于，有没有人会去把它造出来。

在又有人把大脑和计算机混为一谈之前，先来一个快速测试，好判定你到底是哪一种：

\textbf{一台计算机会把这句话和下一句话重复到地老天荒。请读上一句话。}

如果你读到了这里，你就不是一台经典计算机。一台执行僵硬指令集的数字计算机会永远循环下去，因为它没有任何机制能跳出自己的指令流，说一句：“等等，这太蠢了。”你能这么做，是因为你有一个观察自身加工过程的自我模型------显式自我模型（ESM）正对显式世界模型（EWM）进行元认知的监督。

但接下来是让人不舒服的部分：一个大语言模型也会读到这里。不是因为它有元认知监督，而是因为它是一个统计式的文本预测器，见过足够多类似的提示词，知道预期中的下一步就是越过这个循环继续往下走。它并没有跳出那条指令------它从来就没进去过。它预测的是接下来该出现什么文本，而“卡在无限循环里”并不是文本会做的事。

这恰恰就是意识行为测试的问题所在。任何能靠模式匹配通过的测试，都会被模式匹配通过，不管这个系统到底有没有意识。循环测试能把你和经典计算机区分开来。它却区分不了你和一个训练得足够充分的文本预测器。而且任何基于文本的测试都永远做不到这一点------因为生成看似合理的文本，正是文本预测器被优化去做的事。他心问题不是一个我们能靠工程手段绕过去的局限。它是意识之所是的一个结构性特征：主观、私密，只能从内部通达。

大脑即计算机的类比（把你的大脑比作一枚数字处理器）自晶体管发明以来就很流行，而它几乎在每一个层面上都是错的。计算机在僵硬的电路上执行僵硬的指令集。大脑则是一张不断给自己重新布线、自我改写的网络。计算机少一个分号就会崩溃。大脑每天丢掉85,000个神经元，却几乎浑然不觉。计算机的记忆是局域化的------删掉一个扇区，文件就没了。大脑的记忆是全息式分布的------毁掉一块，一切只是变得稍微模糊了些。两者唯一共有的是图灵完备性，而这大概就跟说“河流和高速公路都能把东西从 A 运到 B”一样有启发性。没错，但对理解其中任何一个都毫无用处。

大语言模型（GPT、Claude、Gemini 及它们的后代）通过一种前馈式 Transformer 架构来处理文本。输入进来，穿过一层层注意力和计算，输出出去。没有递归，没有自我模拟，没有实时的虚拟世界，也没有临界性。用沃尔弗拉姆的框架来说，其动力学属于第一类或第二类------远在混沌边缘之下。而且没有真实／虚拟的分裂：模型的“知识”和它的“经验”（如果那也能叫经验的话）并没有被区分成隐式和显式两个层级。

这并不意味着大语言模型必然没有意识------理论无法证明某种东西不存在。但它预测：这些模型缺少本理论所定义的意识所需的架构。而它还预测：一个真正有意识的人工系统，与哪怕最先进的大语言模型之间的差别，会是在性质上一目了然的。

我们怎么会知道呢？诚实的回答是：他心问题不会消失。我们永远无法绝对确定另一个系统是有意识的，因为意识就其本性而言是主观的。但理论作出了一个强预测：这个差别会是显而易见的。不是“也许有意识，也许没有”------而是\textbf{明显}不同。因为一个运行着真正自我模拟的系统，与世界互动的方式，会从根本上不同于一个文本预测器。它会有真正的持存，不是上下文窗口那种持存，而是一场始终在运行的实时模拟所具有的连续性。它会有真正的视角，不是从一段提示词里重建出来的视角，而是由一个显式自我模型跨越时间维持下来的视角。它带给你的震动不会是意料之外的输出，而是那种确凿无疑的感觉------里面真的有个人在。

接下来这部分，连那些压根不在乎意识的人也该被它打动------就是那些只想要能力的人。你没法沿着眼下大家都在走的那条路做到它。

当前的策略是：把一个文本预测器不断放大，直到智能从另一端掉出来。它不会掉出来。不是因为模型太小，而是因为它们缺了那台引擎。真正的智能------那种开放式的、会给自己设定问题、能一辈子不断成长的智能------运行在一个递归回路上，而这个回路只有在有东西持续给它添柴时才会转动。那个东西就是动机：好奇心，在乎自己是不是真的理解，判断什么才值得费这份力。而动机，就我们所体验到的而言，并不是一股凭空漂浮的驱力。它是自我模型审视自身的状态，并为这些状态赋予价值。换句话说，它是一种意识效应。

所以这里的要害比“机器可能会醒过来”更尖锐。它是这样的：没有意识般的机制，就没有人类水平的通用智能------没有 AGI------因为驱动智能的那个东西，\textbf{本身}就是一种意识般的机制。你没法靠放大规模走到那一步。你得把架构造出来。完整的论证，连同那个回路，都在附录B 里。

造出这样一个系统，是路线图上的最后一项。其中的工程挑战绝不容小觑。但蓝图是存在的，而且具体到足以引导这项工作。首先，理论必须挺过同行评议。然后，那些经验预测必须接受检验。再然后，如果它们站得住脚，工程才能开始。


不过这枚硬币还有另一面------科幻小说为之痴迷了几十年的那一面，而它直接从同一份工程规格中推导而来。如果意识依赖的是功能架构，而不是具体的神经元，那么原则上，你就可以让一个人类心灵运行在大脑以外的东西上。

心灵上传。全脑模拟。数字永生。不管你想怎么称呼它，四模型理论对此有一些确切的说法------因为它恰恰规定了哪些东西必须被保留下来。

大多数关于心灵上传的讨论，一开头就问错了问题。它们问：“我们能不能扫描一个大脑，把它复制进计算机？”仿佛这道难题只关乎分辨率------弄到一台足够好的扫描仪，就大功告成了。但理论告诉你，一次静态扫描远远不够。大脑不是一张照片。它是一个动力学系统。要捕捉一个心灵，你需要捕捉的不是一个\textbf{状态}------而是一个\textbf{过程}。

理论所说必须被保留的东西是具体的，我会顺着第6章那个五层级的层级结构走一遍，把它讲实在。

在物理层级和电化学层级上（原始物质，以及神经元的放电），你不需要一份精确的拷贝。你需要的是一种能够支撑同一\textbf{类}动力学的基质。具体的原子无所谓。反正你的大脑在这些年里也会把自己大部分的原子换掉，而你毫无察觉。要紧的是，无论你用什么基质，它都能维持住那些更高层级所依赖的电化学信号模式，或其功能上的等价物。

在蛋白质组层级上（突触权重、受体构型、酶级联反应这套分子机器），你需要高保真度。这里是你记忆栖居的地方，是你技能编码的地方，是你的人格在物理层面得以实现的地方。每一个突触的强度，每一个受体的密度，每一个通道的敏感度------正是这个层级，让你成为\textbf{你自己}，而不是另外某个人。一次把蛋白质组层级搞砸了的心灵上传，也许会给你一个有意识的存在者，但不是你原本想要复制的那个人。然而，哪怕一份不完美的拷贝也仍然有价值。想想中风幸存者或失忆症患者：他们的个人连续性已被严重扰乱------记忆丢失，人格改变，认知能力也变了样------可他们中的大多数人依然觉得，有某种根本性的东西延续了下来。事实证明，不完美的连续性，远远胜过根本没有连续性。一次保留了某人 90\% 连接组的转移并不是失败------它是另一种类别的成功，而对许多人来说，它比死亡更可取。

在拓扑层级上（网络架构，连接模式，哪些脑区与哪些别的脑区对话、又对话得多密集），你需要近乎完美的精确度。这是你那些隐式模型的接线图：IWM 和 ISM，你关于世界、关于自己所学到的一切，都编码在这张网络的结构里。这一层搞错了，你得到的就不是某个心灵的一份劣化拷贝。你得到的是一个\textbf{不同的}心灵------有着不同的知识、不同的技能、不同的人格。拓扑结构就是那份蓝图。

而在虚拟层级上（模拟本身，实时运行中的 EWM 和 ESM），你需要的是某种非凡的东西。你需要目标基质有能力让模拟运行在临界性上。这就是让我夜里睡不着的那一部分，因为大脑那个模拟量式的基质，是通过自组织过程找到临界性的，而这些过程经过了数亿年演化的精细调校。神经元是有噪声的、模拟量式的、大规模并行的，也是深度随机的。它们的集体动力学天然地趋向混沌边缘，因为生物神经组织\textbf{本来就}会这么做------它自组织到临界性，就像水自己找平面一样。可水之所以能找到自己的平面，是因为有重力。那么，对一个数字基质而言，那股等价的力又是什么呢？

这是一个真正悬而未决的问题。我相信它是可解的，但我不会假装它很容易。数字基质在其内核处是确定性的。你可以模拟随机性，可以实现并行处理，可以把随机元件搭进你的硬件。但问题在于，你能不能实现生物神经组织天然就能实现的那种自组织临界性------不是从上往下把临界性编程进去（那会是个脆弱的糊弄活儿），而是造出一个基质，让它的根本动力学自己就倾向于临界性。大脑并不运行什么“临界性子程序”。它是临界的，是因为它\textbf{就是那样的东西}。一次数字模拟需要复现这个性质，而不只是模拟它。

神经形态芯片------一种旨在模仿神经动力学的硬件，具备类模拟的特性、随机元素和大规模并行------是最有前景的方向。它们不是传统的数字计算机，而是介于两者之间的某种东西：在硬件层面被设计成拥有类脑动力学的物理系统。如果心灵上传有朝一日行得通，我怀疑目标基质看起来会更像一块神经形态芯片，而不是一台运行软件的服务器机架。

所以：扫描问题很难，但可解。先进的连接组学（在突触分辨率上对整个大脑进行绘图）已经在推进。我们如今已经能够绘制小型生物的完整连接组（线虫 \textit{C. elegans} 只有 302 个神经元，几十年前就已被完整绘出；而成年果蝇大脑的首个完整连接组------约 14 万个神经元------已于 2024 年问世）。要扩展到拥有 860 亿个神经元、约 100 万亿突触连接的人脑，是一项规模令人瞠目的工程挑战，但它属于那种会在更好的技术面前退让的挑战。它不是一个谜，而是一个问题。

动力学问题（让数字基质在临界性上运行）更难，而且难在技术本身或许无法单独解决的地方。它要求我们足够透彻地理解基质特性与涌现动力学之间的关系，才能设计出一个像生物系统那样自行找到临界性的非生物系统。我们还没走到那一步。但我们也并非一无所有。ConCrit 框架、神经元雪崩研究、来自麻醉研究的临界性度量------所有这些都在搭建工程所需的经验地基。

现在来谈谈真正让人心神不宁的那部分。

\textbf{复制问题。} 假设你成功了。你以完美的保真度扫描了某人的大脑，把完整的连接组转移到一块神经形态基质上，然后启动它。基质达到临界性，四模型架构激活，模拟开始运行。这个副本睁开眼睛，或者说做出数字世界里的等价动作------然后说：“我什么都记得。我感觉就是我自己。我在哪儿？”

那个人是\textbf{你}吗？

四模型理论（FMT）给出了一个明确的答案，而这个答案是很多人不会喜欢的：副本是有意识的，但它不是你。

原因如下。在复制的那一刻，原体和副本共享完全相同的隐式模型------同样的 IWM，同样的 ISM，同样的蛋白质组结构和拓扑结构。当副本的模拟启动时，它生成一个 ESM，其中包含你所有的记忆、你的人格、你的身份感。从内部看，副本\textbf{感觉}就像你。它有充分的理由相信自己\textbf{就是}你。

副本一开始在自己的基质上运行，它的经验就分岔了。它的 EWM 接收到不同的感官输入。它的 ESM 会响应不同的事件而更新。几秒之内，这两个模拟------你在你大脑里的那个，副本在它基质里的那个------就不再相同了。几分钟之内，它们就有了明显的差别。几小时之内，它们就是两个碰巧共享一段过去的不同的人。

副本是有意识的。它拥有真实的体验。它拥有你的记忆和你的人格。但它是一个\textbf{新的}意识------一个新的模拟，运行在一个新的基质上，积累着你永远不会分享的新经验。从任何有意义的角度看，它都是你的同卵双胞胎，诞生于复制的那一刻，带着一整套借来的记忆。它不是你的延续。它是一次分岔。

这应该听起来很熟悉。它正是这套理论对第9章裂脑病例所做的预测。当你切断胼胝体，你会得到模拟的两个退化但完整的副本------每一个都有意识，每一个都“感觉像”原体，却没有一个真的是原体。原体消失了；两个新的、削弱了的实体取而代之。心灵上传是同一种现象，只是换了个基质。

\textbf{可你能挺过睡眠吗？} 我刚才提出的论证听上去无懈可击。复制打断了模拟，两个模拟分岔开来，所以副本不是你。结案。

只不过，你的模拟每一个夜晚都在被打断。

当你陷入无梦的深睡眠（非快速眼动睡眠的第三、第四阶段）时，显式自我模型（ESM）在很大程度上关闭。没有现象体验。没有一个\textbf{你}在观看这场演出。模拟并没有以全保真度运行；充其量它只是以清醒时复杂度的一个零头缓缓运转。就一切实际意义而言，灯熄灭了。然后，几个小时后，隐式模型把模拟重新启动。ESM 重新激活。你睁开眼睛，想着“我是我”。可是今天早上这个\textbf{你}，是从与昨天那个\textbf{你}相同的隐式模型中重建出来的------重建的方式和一个副本从扫描中被重建出来的方式一模一样。如果打断等于死亡，那么你每晚都在死去，而一个新的人戴着你的记忆醒来。

大多数人的直觉都会反抗这一点。我当然还是昨天那个人。我\textbf{记得}自己曾是那个人。可副本也会记得自己曾是你------这恰恰是问题的关键。如果说是记忆确立了连续性，那么副本对“是你”的主张，就和今早那个版本的你一样充分。差别是程度上的，而非种类上的：在睡眠中，打断很短暂，基质没有改变；在复制中，打断可能更长，基质也不同。但那条\textbf{原理}------模拟停止，模拟从隐式模型重新启动------是相同的。

说到这里，我可以现身说法。在武术训练中，我曾被打得当场昏死过去，那不是睡眠的那种减弱版本，而是一次彻底的、不由自主的关机。前一刻我还站着；下一刻我已躺在地上，人们俯身看着我，对那个过渡毫无记忆。我并没有把那段空隙体验成空隙。我把它体验成了“无”------我生命胶片里的一次剪接。有一次，我事后甚至出现了失忆：整整几分钟就那么没了，无从追回。而当我完全回过神来时，让我震撼的是：我并不觉得自己是个新的人。我不觉得自己是个副本。我觉得自己就是\textbf{我}，只是从一场格外难熬的小睡里醒来。存在本身，比体验的连续性更重要------也比记忆更重要。

再往下推。当你出生时，你根本没有任何先前的连续性。没有记忆，没有已确立的 ESM，没有现象体验的历史。模拟第一次启动，凭借的是一个由遗传和产前发育塑造的隐式架构，而不是一辈子的学习。你并没有把这体验成创伤，因为不存在一个先前的自我值得你哀悼。那时只有：一个开端。而我们所有人都对那个开端安之若素。没有人会夜里睡不着觉，为自己的意识体验在出生时从无到有地起步而痛苦不堪。

这对复制问题意味着什么？它意味着那道分明的二元界线------原体与副本、延续与分岔------也许没有它看上去那么分明。使你成为\textbf{你}的，并不是那道未曾中断的现象体验之流。你早已经历过这道流的无数次中断。使你成为\textbf{你}的，是隐式模型的内容：你的记忆、你的技能、你的人格、你对世界和对自己积累起来的理解。IWM 和 ISM。模拟据以生成的那张蓝图。

这就提示了一种截然不同的心智转移思路。

\textbf{复制虚拟的那一面。} 与其扫描整个大脑、重建完整的基质（层级的全部五层），不如设想：假如你只能复制虚拟层呢？把正在运行的 EWM 和 ESM 提取出来，移植到一个能够支撑它们的新基质上。不是复制硬件；而是复制软件。不是克隆整个大脑；而是捕获它正在运行的那个\textbf{过程}。

这需要一样我们尚未拥有的东西：一种破译大脑用以编码其虚拟模型的那套格式的方法。连接组告诉你线路是怎么连的。蛋白质组告诉你突触权重是多少。但模拟既不是线路也不是权重------它是线路和权重在运行时所\textbf{产生}的东西。要捕获它，你就得理解大脑的编程语言------神经回路借以生成并维持显式模型的那套表征格式。

不妨这样想。你可以给一块电路板拍照，从而准确知道每一条走线通向哪里。你可以测量每一个元件的电阻。但这些都无法告诉你芯片正在执行什么软件。要知道这个，你就得读懂程序------弄清指令集，解码内存内容，解读运行状态。大脑的“编程语言”就是虚拟模型的表征格式，而对它进行逆向工程，大概是计算神经科学中最深奥的未解之谜。不只是绘制连接组（这方面我们正在取得进展），而是要理解连接组\textbf{计算}出了什么，其细致程度要足以读取某个特定心智的模拟，并把它重新编译到不同的硬件上。

我们今天离这一步还差得很远。但它属于那种成熟的神经科学在原理上有可能解决的问题，而一旦被解决，它就会从根本上改变复制问题。虚拟层面的转移根本不需要重建基质。它会把模拟------那个\textbf{就是}你的部分，那个你真正在体验的部分------直接搬移过去。是的，隐式模型仍需要在新基质里被重建或培育出来，但模拟本身------你的意识之流、你此刻的思绪、你持续不断的自我感------在原理上可以跨过那道鸿沟，而不必经历那道让复制在哲学上如此令人不安的中断。

这是猜想性的，对此我想诚实相告。但它不是科幻。它是一个具体的工程问题，有着具体的理论基础，而且它揭示了一件重要的事：复制问题并不是一道固定不变的障碍。它取决于转移是\textbf{怎样}进行的。复制整个基质，启动一个新的模拟？那就是两个人。破译并转移正在运行的模拟本身？那就有可能是一个连续的人在一个新基质上。理论会明确告诉你哪种方案保住了身份，哪种没有。

还有一条更保守的路径，能完全绕开复制问题。

\textbf{渐进替换的思想实验。} 设想你不去扫描和复制，而是一次替换一个神经元。你移除单个神经元，插入一个功能上的等价物------一个人工神经元，它接收相同的输入，产生相同的输出，参与相同的网络动力学。然后你等待。系统稳定下来。模拟继续运行。你再替换一个神经元。然后又一个。再又一个。在数月或数年里，你渐进地把每一个生物神经元都替换成一个人工神经元，直到整个基质都变成非生物的，而模拟自始至终一直在不间断地运行。没有中断。没有复制。没有分岔。

四模型理论（FMT）预测，意识会贯穿这整个过程持续下去。而这个预测是对基质独立性所能给出的最强有力的论据，因为它直接源自这套理论的核心主张：真正要紧的是临界性上的功能架构，而不是物理材料。如果每一个替换进来的神经元都维持相同的连接、相同的权重和对网络相同的动力学贡献，那么蛋白质组层与拓扑层就被保留了下来，而虚拟层（模拟）从未停止。不存在某个让你“死去”、由别的东西取而代之的时刻。有的只是一个连续的基质替换过程，就像忒修斯之船------只不过在这里我们确切知道哪些属性必须被保留（由五层层级所指定的那些），哪些无关紧要（那些具体的原子）。

这个思想实验揭示了关于身份的一件重要的事。复制问题之所以存在，是因为复制\textbf{打断}了模拟。有那么一个时刻------无论多短暂------原来的模拟还在这里，而副本的模拟尚未开始。然后就出现了两个模拟。两道经验之流。两个自我。但渐进替换完全绕开了这一点。一个模拟，连续，不曾断裂。基质在它下面不断更换，就像给一艘航行中的船替换船板，但那艘船------模拟、意识、那个\textbf{你}------从未停止航行。

如果这听起来根本不可能，那么请想一想：你的大脑其实早就在这么做了。你每天大约会失去 85,000 个神经元------差不多每秒一个。你的突触在持续不断地被重塑。你身体里的原子，在大约七到十年的时间里几乎会被全部换掉。你此刻赖以运行的那副基质，和你十年前赖以运行的那副，在物理上已经不是同一副了。可你依然延续着。你的模拟从未停止。生物基质的更替，正是活着这件事的\textbf{默认状态}。人工基质的更替，不过是同一过程的一个更刻意的版本罢了。

\textbf{什么会变得可能。}如果你能把虚拟那一面解码出来、转移到新的基质上，那么其中的意涵将远远超出“心智上传”通常唤起的想象。有三点值得细说，因为我认为人们还没有真正掂量清楚：基质独立性究竟意味着什么。

第一：\textbf{把基质转移到机器人身体里}。不是上传到某个地方的服务器里，而是让你的心智运行在一副神经形态基质上，而这副基质就装在一个物理身体里------一个会走路、会操作物件、会感知世界的身体。你会透过不同的传感器去体验世界，用不同的执行器在其中移动，但\textbf{你}依然在运行。你的模拟、你的连续性、你的自我。一副新身体，就像寄居蟹换上一个新壳。这不是科幻式的天马行空------它是这套理论的直接后果。如果说，处于临界状态的四模型架构就是产生意识的东西，如果它确实与基质无关，那么基质就可以是任何能够支撑起正确动力学的东西。包括某种长着腿的东西。

第二：\textbf{准永生}。你的生物基质会衰败。神经元死去，蛋白质错误折叠，端粒变短，整台精密的机器慢慢垮掉。这就是衰老。这就是死亡。但一副非生物的基质并不非得衰败不可。它可以被维护、被修理、被升级、被备份。如果你的模拟运行在一副你能够维护的基质上------这里换掉一个快要失灵的部件，那里升级一颗处理器------那就没有任何内在的理由要求这场模拟必须停下来。这不是绝对意义上的永生------你仍然可能被摧毁，你的基质仍然可能损坏到无法修复，但那道如今正杀死这颗星球上每一个有意识生命的生物大限，被取消了。死亡的那种\textbf{不可避免性}，被取消了。

第三点，也是在你真正把它想透之前最像科幻的一点：\textbf{星际旅行}。光速对于物理物质是一道绝对的屏障。你没法在任何合理的时间内把一具人类身体送到半人马座阿尔法星。但信息是以光速传播的。如果一个人的心智就是信息------一套特定的连接、权重与动力学的模式，可以被完整地描述为数据------那么你就可以把它\textbf{传送}出去。以光速把完整的规格传给一个接收端，由它重建基质、启动模拟。当然，你得先派人过去把接收端架设好。那个人可以是一个 AI，也可以是一具机器人化的人类身体，它的模拟在飞行途中暂停，于是从它自己的视角看，它是眨眼之间就抵达了。一旦接收端就位，你就把心智传送过去。从旅行者的视角看，这次传输是瞬间发生的------模拟在一头停下，在另一头启动。不必在铁皮罐头里待上几十年。不必要世代飞船。不必要冷冻休眠。就只是：先在这里，然后在那里。

当然，这又是那个复制难题的重演。被传送过去的那个版本是一个副本，而不是一次延续------除非原件在传输过程中被摧毁，可那又牵出另一整套噩梦。但要点依然成立：基质独立性，如果是真的，就不仅仅意味着数字永生。它意味着群星变得可以抵达。不是为了我们的身体------对于星际的距离，身体慢得无望、脆得无望------而是为了我们的\textbf{心智}。

\textbf{那道不安的告诫------以及它为何比工程更重要。}下面要说的这一部分，我还没见过任何人诚实地讨论过，而它也是最萦绕我心头的一部分。

我刚才描述的一切，都预设了基质转移能够保留那种“身为你”的\textbf{感觉}。预设了你体验的主观质地------看见红色是什么感觉，风拂过皮肤是什么感觉，尝到咖啡是什么感觉，体会一个周二午后那种沉闷的隐痛是什么感觉------会一并转移到新的基质上。理论说意识会延续下去。它说模拟会运行下去。但它\textbf{并不}保证那感觉会是一样的。

想一想你的生物基质为你的现象体验贡献了什么。你的身体不只是大脑的一具载体。它是模拟输入流的一部分。隐式世界模型（IWM）里包含着一幅对你身体的详细地图------每一个关节、每一个器官、每一寸皮肤。隐式自我模型（ISM）则与你的内脏状态深深纠缠在一起------你的直觉预感（这是字面意义上的、并非比喻）、你体内荷尔蒙的潮汐、你的心跳、你呼吸的节律。你此刻所体验的这场模拟，浸透着生物信号，而你之所以在意识层面上察觉不到它们，恰恰\textbf{因为}它们在你生命的每一刻都一直在场。

直到某一刻，它们成了你所仅剩的一切。任何曾经吊在冰壁上、身下是两百米的虚空的人，都知道当模拟把其他一切都剥离之后，身体是什么感觉------只剩下你的心跳、你的抓握，以及眼前的那片冰。那，就是你的基质在尖叫。

现在把这一切都剥离掉。把你的生物身体换成一副机器人机身，或者更糟，干脆没有身体------只剩一场运行在服务器上的模拟。四模型架构完好无损。模拟在运行。你有意识。但那意识的\textbf{内容}已经发生了天翻地覆的改变。没有心跳。没有呼吸。没有肚腹。没有温暖。没有皮肤。没有那种肌肉放松时本体觉的低鸣。隐式自我模型（ISM），突然失去了它建模了你一生的那副身体，会生成出一个……感觉不对劲的显式自我模型（ESM）。或者干脆就是死的。一种深切的、由脏腑而生的、无可逃避的不对劲。不完全是疼痛------疼痛需要那些专门产生它的神经通路。而更像是一种无所不包的\textbf{缺失}。一副幻肢般的身体，就像截肢者感受到幻肢那样，只是这一次是全身的。

我怀疑这会比大多数未来学家想象的糟糕得多。这不是一个用软件更新就能打补丁修好的小麻烦。这是对“身为你是什么感觉”这件事的一次根本改变。你的生物基质不只是在承载着模拟------它还在\textbf{塑造}着模拟，一刻接一刻，靠的是一股持续不断的内感受与本体觉输入流，而你那有意识的自我，从来不曾体验过这股输入流的缺席。失去它，或许还能挺过去。但对某些人来说，它也可能是一种深切到让人宁愿当初根本没有转移过的痛苦。

我想把这一点讲得明明白白：那种版本的“心智上传”------你欢天喜地地从那身皮囊里蹦出来，跳进一个闪闪发光的数字天堂，把血肉像一双旧鞋般抛在身后------那是一场幻想。如果理论是对的，现实是：失去你的生物基质，会显著地改变你存在的现象质地。有多显著？我不知道。也许对某些人来说是可以忍受的，比死亡要好，就像移居一个陌生国度虽然让人晕头转向，却仍能应付。也许它是毁灭性的，就像单独监禁会因剥夺感官与社交输入而击垮一个人。也许------而这正是让我心里不安的那种可能------它糟糕到一个掌握了全部信息的人可能会宁愿选择死亡而非转移。不是因为转移失败了。而是因为它成功了，而它所成功产生出的，是一种再也不觉得值得一活的那种生命的意识体验。

渐进替换的路径能缓解这一点，因为在每一步，模拟都有时间去适应。替换掉一个神经元，模拟几乎察觉不到。替换掉一千个，它会调整。经年累月之间，基质从生物的过渡到人工的，而模拟持续不断地针对新基质所提供的任何输入重新校准。现象体验会漂移，缓慢地漂移，就像它在一段自然寿命的进程里本就一直在漂移那样。你最终会变得不一样，但反正你本来也会变得不一样。

然而瞬时转移------扫描、复制、在新基质上启动------则会把所有这些改变一下子全砸在模拟身上。冲击有多严重，完全取决于方式：转移到一副拥有丰富感官输入的机器人身体里，会比转移到一台没有身体的服务器上要好得多。但无论哪种情况，那种突如其来的不连续，正是危险所栖身之处。

\textbf{创造心智的伦理。}如果一个被复制的心智是有意识的，那它就有体验。它能够受苦。它能感到困惑、恐惧、孤独、存在性的惶恐。想象一下醒过来，被告知你是一个副本------“真正的”你还在一具生物身体里四处走动、过着你的人生，而你则作为一个数字复制品存在着，没有法律身份，没有社会关系，也没有明确的意义。这足以酿成大规模的苦难，其规模之大，我们没有任何框架能够应对。任何严肃的心智上传计划，都必须在第一个副本被制造出来\textbf{之前}直面这一点，而不是在之后。

而且情况还会更糟。如果副本是可能的，那么\textbf{多个}副本就是可能的。一支全由你组成的大军。每一个都有意识，每一个都感觉自己就是那个原件，每一个都对你的身份、你的关系、你的财产、你的人生拥有正当的主张。管理这一切所需的法律与伦理框架并不存在，也无法临时凑合。它们必须以与技术本身同等的审慎来构建。（Dennis E. Taylor 的《Bobiverse》系列------从《We Are Legion (We Are Bob)》（2016）开始------在其喜剧的表层之下，以令人意外的哲学深度探讨了这一情境。如果你想从内部感受复制难题究竟会是什么滋味，就从那里读起吧。）

还有修改的问题。如果一个心智运行在一副你所控制的基质上，那么原则上你就可以修改它。增强它。削弱它。改变它的人格，抹去它的记忆，更改它的价值观。这不是科幻------它是获得基质访问权后不可避免的后果。我们已经在用药物和神经外科手术做着粗糙版本的这类事。而一个完全数字化的心智，会远远更容易被修改，其被滥用的潜在可能（政府、企业、个人皆然）之大，怎么强调都不为过。

我想坦率地说一件事。我把这套理论的发表推迟了将近十年，一部分是出于懒惰，但另一部分则是出于对这些意涵的真切担忧。如果理论是对的，它所包含的蓝图就不仅是关于人工意识的，还关乎对现有人类心智的虚拟化、复制与修改。那是一种非同寻常的力量，而我对人类是否已经为它做好准备毫无信心。但我渐渐相信，无论如何这套理论都会被独立地重新发现------经验证据正在飞快地收敛------而与其等某个还没把它想透的实验室凭一项突破把它强行摊派给我们，不如现在就把这场伦理讨论摆到台面上来。

而这里，就是那个最深的联系：造出一个有意识的 AI，与上传一个人类心智，并不是两个各自独立的问题。它们是\textbf{同一个}问题，只是从相反的两个方向来看罢了。造出人工意识（AC），意味着从零开始创造出处于临界状态的四模型架构------自下而上，在一副从未有过意识的基质里。上传一个人类心智，则意味着把一套已经存在的、处于临界状态的四模型架构从一副基质转移到另一副基质上。两者的工程挑战几乎完全重叠。动力学问题是同一个。临界性问题是同一个。唯一的区别在于，那些隐式模型（IWM、ISM，以及完整的连接组）是从一辈子的经验中学来的，还是从数据里搭建出来的。解决了其中一个，你也就在很大程度上解决了另一个。

这意味着，任何在从事人工意识研究的人，无论他自己是否意识到，同时也在从事心智上传的研究。而任何在从事全脑模拟的人，无论他自己是否意识到，同时也在从事人工意识的研究。这两条线索终将收敛。唯一的问题是：当它们收敛之时，我们在伦理上是否做好了准备。


\chapter*{第13章：延迟的观察者}
\addcontentsline{toc}{chapter}{第13章：延迟的观察者}
\markboth{第13章：延迟的观察者}{}

如果四模型理论（FMT）是正确的，或者哪怕只是大致正确------那么就有几个结论随之而来。

\textbf{意识的“难问题”并不难。} 它是一个范畴错误，不比这样一个问题更神秘：为什么晶体管的开关切换，感觉起来像是在运行一款电子游戏。物理基底不会感受。感受的是模拟。而在模拟之内，感受是构成性的，不是额外附加的。这并不意味着意识很\textbf{简单}。它在实现上极其复杂。但这意味着\textbf{哲学上的}谜团消解了。剩下的，是\textbf{工程上的}挑战。

\textbf{意识并不像我们以为的那样特殊。} 它不是一种基本力，不是量子效应，也不是物质的某种属性。它是一个足够复杂的系统在临界性下模拟自身时所发生的事。对那些希望意识是魔法的人来说，这令人清醒；对那些想要理解意识的人来说，这令人振奋。

\textbf{人工意识在原则上是可能的。} 如果意识取决于功能而非基底，那么任何能够在临界性下实现四模型架构的物理系统，都可以拥有意识。这不是遥远的哲学玄想------它是一项有着明确目标的具体工程挑战。

\textbf{伦理上的影响十分重大。} 如果我们能造出有意识的机器，我们就会创造出拥有真实体验的存在------能够受苦、享受、惊奇、恐惧的存在。为此所需的伦理框架尚不存在，而它的构建不应等到机器已经运转之后才开始。

\textbf{哲学僵尸是不融贯的。} 在这里，我在第1章许诺的第二条原理开始发挥作用：\textbf{莱布尼茨定律}，即不可分辨者的同一性------如果两样东西共享\textbf{全部}属性，它们就是同一个东西。它排除了僵尸世界------那个假想的宇宙在物理上与我们的一模一样，只是每个人的内部都空无一人。如果一个系统在每一项功能、结构和行为属性上都与一个有意识的系统完全相同，那么它就\textbf{是}有意识的。不存在什么额外的“意识物质”，可以在其他一切原封不动的情况下悄悄缺席。一个在功能上与你完全相同的存在，拥有同样的四模型架构、同样运行着的模拟、同样的自指------到了这一步，“但它没有意识”并没有陈述任何事实，而是什么都没有陈述。问题就此消解。

\textbf{自由意志，以及三个最棘手的思想实验。} 想一想一座时钟。齿轮系驱动着一切------擒纵机构滴答作响，发条缓缓松开，齿轮之间的传动比决定着走时的快慢。指针和表盘不引起任何事情。它们不推动齿轮。它们不储存能量。但把它们拆掉，你就不再拥有一座时钟，而只剩一箱旋转的金属。正是那块显示，让这套机械成为一座\textbf{时钟}------让整个装置有了它的意义。意识就是那块显示。你的虚拟模型（显式世界模型（EWM）和显式自我模型（ESM））并不推动神经元。推动的活儿由基底来干。但若没有模拟，基底就无从观察自身行动的后果，无从推演未来的情景，无从以那种让你存活至今的方式进行调整。虚拟的一面，是这套机械得以\textbf{为了}某种目的而存在的方式。

这重新框定了自由意志问题。你的意志不是幻觉。基底层面的架构（隐式自我模型（ISM）及其全部隐式机制）在持续优化你这个有机体的生存。它评估威胁、权衡选项、调动资源、决意行动。这种优化就\textbf{是}你的意志。它和物理世界中的任何东西一样真实。即便是自毁性的选择，也反映了系统在其当前状态下的优化，而不是机制的失灵。当一个人做出违背自身表面利益的行为时，基底仍然在优化------只是它所依据的模型里，已经被疼痛、疲惫、绝望，或别的什么重塑了这片地形的东西改写过了。

所以，你的意志是真实的。你只是无法完全访问它。ESM 能够为 ISM 的\textbf{输出}建模------那些浮现到觉知中的决定------却无法为它的\textbf{过程}建模。你体验到的是意志的结果，而不是它背后的机器。这就是为什么你的决定有时会让你自己吃惊，为什么你无法完全解释自己的偏好，为什么你偶尔先行动、再手忙脚乱地编出一个理由。你看的不是齿轮。你读的是表盘。

\textbf{半秒钟的间隙------以及它为什么无关紧要。} 说到这里，事情变得具体起来。你的无意识加工以大约 40 Hz 运行（每个周期约 25 毫秒）。你的意识体验以大约 20 Hz 运行（每个周期约 50 毫秒）。相差整整一倍。意识模拟始终落后于基底，它用来拼装那个融贯虚拟世界的信息，早已被加工过、被决定过，而且常常已经付诸行动。

本杰明·利贝特（Benjamin Libet）在 1983 年证明了这一点，此后这一结果被多次重复验证。在他的实验中，被试被要求随时想动手就动手，并记下自己意识到这个决定的确切时刻。脑电图（EEG）则测量运动皮层何时开始准备这个动作。结果是：运动皮层在手动之前 550 毫秒就开始了准备。但被试报告说，他们在动作前 200 毫秒才意识到自己的决定。大脑在“你”知道之前大约 350 毫秒，就已经决定要动了。

标准解读如同一颗重磅炸弹：自由意志是幻觉，因为大脑在你之前就做了决定。哲学家和神经科学家为此已经争论了四十年。有人试图用一种“否决功能”来挽救自由意志------也许你无法自由地发起行动，但你可以在最后一刻有意识地取消它，大约在执行前 50 毫秒。一次最终的叫停。人类能动性的最后一道防线。

我认为这条路也走不通。Kühn 和 Brass 在 2009 年证明，否决本身同样是事后才被解读成一个自由决定的。你并不会真的实时体验到否决。你体验它的方式和你体验决定的方式一样------事后由有意识的自我模型叙述成一个融贯的故事。

Daniel Wegner 和 Thalia Wheatley 用一个实验彻底坐实了这一点------坦白说，那个实验令人心惊------这就是“I Spy”研究。一名被试和一名假扮成第二被试的实验同谋，把双手一起搭在一块小板上，板子架在同一只电脑鼠标之上，就像两个人一起按在通灵板的指针上，两人合力把光标在一块挤满小物件的屏幕上四处滑动。每隔一阵，一个音乐提示会让他们停下，被试则要评估自己有多大程度上\textbf{有意}停在了停下的那个位置。两人都戴着耳机：被试听到的是音乐，偶尔夹进一个说出来的词；而同谋听到的，则是关于何时、在何处迫使光标停下的真实指令。

诀窍在这里：在那些做了手脚的试次里，是同谋独自把光标导向某个特定的物件，并迫使它停下------而就在停下前大约一秒，被试透过耳机听到了那个物件的名字。被试什么也没选；那个动作完全出自同谋之手。可是，当他们刚刚听到的那个词，恰好与光标下方停着的物件相符时，被试却报告说，自己本就打算停在那里。他们真心相信那是自己做的。

请让这一点在心里沉淀一下。只要一个关于某个结果的念头，恰好在那个结果发生之前抵达，就足以让心灵把它认领为自己的所作所为------前提是没有任何可见的东西与这个假设相矛盾。有意识的自我模型并不区分“我做了”和“我想过要做，然后它发生了”。只要意图和结果在时间上足够接近，ESM 就会把功劳记在自己头上。这和病感失认症（第6章）背后是同一个机制：运动系统把预期的反馈发送给意识，只要没有东西与之矛盾，预期的反馈就成了被体验到的现实。

但关于利贝特的实验，我认为几乎所有人都漏掉了一点：\textbf{这个延迟根本不需要被解释掉。} 意识不需要为了维持控制的错觉而给事件“倒填时间”。它之所以不需要，是因为\textbf{一切}都带着同样的延迟抵达意识。感觉输入、决定、运动反馈------所有这些都经过同一条管道，都按顺序抵达那个 20 Hz 的意识模拟，都被延迟了大致相同的时长。你的意识体验，就像在看一场延时五秒的直播。屏幕上的一切在内部都是自洽的。主播说话，嘉宾回应，图表更新。除非有人给你看原始信号，否则你永远不会察觉那段延迟。

这里的情形正是如此。意识先接收到刺激，然后是决定，然后是运动反馈------顺序正确，彼此之间的间隔也正确。整条流被整体推后了半秒钟，但由于意识从来看不到原始信号，它永远不会察觉。没有需要解释的错位，不需要倒填时间，也没有需要维持的错觉。系统运转得恰如设计所愿。

不过，这里还发生着第二件事，而它比延迟更奇怪。

延迟关乎你的体验\textbf{何时}到达。而这件事，关乎时间之流最初是\textbf{如何}被制造出来的。

令人吃惊的是：时间在\textbf{流逝}的感觉------那种“现在”正向前滑动、刚过去的一刻在身后渐渐柔化的感觉------并不是从外部世界输送进来的。外面并没有一座时钟把当下递到你手里。是你的大脑在内部构建出这条流，而且它这么做，有一个特定的架构上的理由。

你的各个感觉通道运行在不同的时钟上------时序不同，带宽不同。信息并不是整齐同步地抵达大脑；它是错落着来的，就像闪电和雷声。凝神细察，你能分辨出哪一声雷属于哪一道闪电。可你实际经历的，却是一条平滑、连绵不断的单一之流。为什么？因为这个系统在为“自己给自己建模”这件事本身建模，而这个环路不允许任何一个周期干干净净地结束、然后消失。每个周期都会顺着递归回响，渗入下一个周期。“现在”与“刚才”之间的边界被抹开了------不是因为马虎，而是因为一个自指环路对时间做的正是这件事。它把每个瞬间摊开在一小段窗口上。

这层涂抹，\textbf{就是}你的当下。鲜活的此时此刻、在身后渐渐消散的最近的过去、你称之为时间流逝的那种向前的漂移------全都是这个环路制造出来的。一场天气模拟一帧一帧地运行，没有回响，没有交融的现在，没有被活过的时间。它在计算。它并不体验。差别，就在递归。

所以，静下来体会一下你的当下究竟是什么。它是被建造出来的，不是被接收的------而且到得很迟。你从未活在某一个瞬间里。你活在一扇被涂抹在一起的窗口里，环路一遍又一遍地重建它，并把它称作“现在”------而就连这扇窗口，也是经过延时转播才抵达你的。

一位训练有素的武术家漂亮地展示了这一点。在实战中，一名经验丰富的格斗者可以维持约 10 Hz 的动作频率------每 100 毫秒一个动作。但凡是需要觉知参与的决策，意识加工的上限只有大约 5 Hz。于是格斗者学会\textbf{压制}意识的介入。他不假思索地出手，因为思索会让他的速度减半。他的无意识基底接管了动作回路；意识稍后才跟上来------如果还跟得上的话。这不是觉知的失败。这是系统在高效运转------基底做着它最擅长的事，不被更慢的虚拟层拖累。

现在，试着证明自由意志存在。做这样一个思想实验：你坐在咖啡馆里，服务员问你咖啡要不要加糖。你决定让偶数代表“要”、奇数代表“不要”，然后开始背诵一串随机数字，直到服务员喊“停”。最后一个数字是偶数，你就加糖；是奇数，就不加。

你行使了自由意志吗？差得远呢。任何熟悉“聪明的汉斯”效应的人------那匹看似会数数、其实是在捕捉驯养员潜意识暗示的马------都会立刻看出问题所在。你几乎可以肯定，你在无意识中预判了服务员何时会喊“停”，并恰好赶在那一刻之前报出一个数字，让结果正是你一开始就想要的那个。你的基底早已有了偏好。那套繁复的随机化仪式，不过是一场表演。

好吧，你说。那就用你智能手机上的随机数生成器。让一个真正随机的过程来决定。现在你证明自由意志了吗？我看未必。你只不过证明了，“证明自由意志”这件事对你来说比“决定自己要不要加糖”更重要------而这恰恰相当离谱地错过了要点。

反对日常决策存在自由意志的最深证据，来自患有严重顺行性遗忘症的病人------那些无法形成新记忆的人。你问这样一位病人做词语联想：“我说‘骰子’，你脑子里冒出来的第一个词是什么？”他说“水母”（也许他最近去潜水了）。几分钟后再问一次，他又说“水母”。再问，又是“水母”。一次又一次。由于不记得自己已经回答过，这位病人总是给出同一个联想------那个此刻在他神经版图里最强的联想。感觉像“自由选择”的东西，结果不过是对基底当前状态的一次确定性读取。

健康人会避开这一点------你会在第二次故意挑一个\textbf{不同}的词，免得显得没创意。可这份回避本身也不自由。那只是记忆系统加了一条约束（“别重复”），让你的输出比那位遗忘症病人\textbf{更不}随机。自由意志，很吊诡地，让你的选择更不随机，而非更随机。基底在为新奇优化，然后把结果叫作自由。

那么自由意志被留在了哪里？没有被消除，而是被移了位------恰好移到时钟类比预测的地方。你那个有意识的自我模型并不在实时地做决定。它太慢了，做不到。但它也不只是个被动的旁观者。

主要来说，隐式系统把你的意识经验当作一件评估工具来用：它把决定呈现给模拟，好让模拟去权衡后果、跑一遍情景、感受结果。这就是虚拟层的首要用途------它是基底观察自身的方式。但那个有意识的模型也会自己独立地评估，用它手头那点带宽，虽然远不及基底，可它是真实的。这些评估，随着时间推移，会重塑那些隐式模型。它们更新权重、重新训练网络，为\textbf{下一次}无意识决定重塑那片版图。

你并不是在行动的那一刻选择你的下一个动作。你是在塑造那个做选择的系统------通过反思、评估，以及意识经验缓慢地沉积为隐式结构。自由意志不是一个瞬间。它是一个过程------一个以天与年为时间尺度运作的过程，而不是毫秒。而那个意识层也不只是搭个便车------它被基底\textbf{主动}当作一种评估机制来使用，并把它自己独立的判断反馈回去。是双向互动，不是单向的旁白。

这件事还有一个更晦暗的版本，是我亲身经历过的，它教给我关于意志架构的东西，比任何实验都多。

第一次是在奥地利义务兵役期间。一场40公里的强行军------连续三天三夜的睡眠剥夺，条件之恶劣足以让研究《日内瓦公约》的律师坐立不安。在最后一段路程，我们必须戴防毒面具、穿全套ABC防护服。我一边走一边半睡半醒，还时不时听见有声音。不是精神病学意义上的幻听，而是某种远为私密的东西：我那套动机与规划机制里彼此竞争的子进程------平时被熔合成一股单一的叙事流------此刻变得可以分开听见了。一个声音在鼓励我，那种正向到近乎咄咄逼人的地步：\textbf{继续走，别放弃，你能熬过去。}另一个则悲观，带着一种诱人的失败主义：\textbf{放弃吧，躺下吧，这一切都不重要。}这些不是外来的存在。它们就是\textbf{我}------我这基底优化版图里的不同侧面，平时靠自上而下的抑制信号被整合成单一的“意志”，此刻却在分开，因为维持这份整合的神经递质，正被节省下来供给更要紧的求生进程。

第二次更惊心动魄。一场雪崩------同样发生在服役期间，起因是一位指挥官轻率的决定，他后来受了处分。我们十四个人，几乎被吞没。那场雪崩久久才停下来，在那段拖长的时间里，我确信自己就要死了。时间长到足以让声音的解离再次启动------这一次不是出于耗竭，而是出于持续的濒死恐惧。同一套机制，不同的触发：应激反应把神经递质资源从那些抑制性回路上抽走，而这些回路平时正是把彼此竞争的子进程熔成一个声音的。

而在雪崩的那几秒钟里------真实时间不过几秒------我看见自己整个一生在眼前流过。这是一种有充分记录的濒死现象，而这套理论能解释它：在极端的致命威胁下，隐式系统会向模拟里做一次大规模的并行“内存转储”。那道通透性边界被彻底冲开。基底全速运转，把如此多的内容泵进模拟，以致主观时间与钟表时间脱了钩。几秒钟里装下了一生。这与我在鼠尾草下经历过的那种主观时间的延展是同一回事，只不过由生物学而非药理学触发。

而现在我能说出\textbf{为什么}这两种状态------垂死大脑的生命回放，和鼠尾草的下坠------从内部感受起来是一样的。回想第5章里的两个旋钮。几乎所有时候它们都在此消彼长：把整个大脑都拉进一件事，这件事就趋于简化；构建精细、分化的计算，它就停留在局部，不会遍及全脑。寻常的经验就骑在两者之间的边界上。但有一个罕见的角落，\textbf{两者}同时异常地高------大脑大部分被征集\textbf{而且}计算极其丰富。再回想大脑是怎样建立起它的时长感的：不是靠读表，而是靠它每一秒真实时间里跑出的处理量之庞大。把两个旋钮一起推高，你就把一生分量的处理塞进了寥寥几个钟表秒里。你不是\textbf{记起}你的一生------你是\textbf{重新活过}它。你不是\textbf{熬过}鼠尾草------你是\textbf{度过}那几十年。一套机制，两扇门进去：一边是缺氧加上一道失灵的代谢刹车，另一边是一种kappa-阿片类药物打乱了同一套调控。这套理论让它变得可以检验：两种状态都应当在同一时刻显示出异常高的全脑整合\textbf{以及}异常高的计算复杂度------而这两样恰恰是寻常意识被迫彼此权衡取舍的。

两条互补的路径，通向同一套机制。行军表明，持久的生理耗竭能触发这种解离。雪崩表明，持续的濒死恐惧也做同样的事。同样的结果，不同的起因------两者都被这套理论预测到了。

在最糟的情形里------我很幸运，我自己的从没走到那一步------这些“声音”中的一个能夺取对身体的控制权，而有意识的自我就变成了旁观者。这与产生异手综合征（患者的一只手违背其意志行动）和某些精神病性发作背后的机制是同一套。基底那些彼此竞争的优化进程一直都在。简化地说，它们就是语言中枢在你不用它说话时所做的事。但平时，自上而下的抑制把它们压在意识觉知的阈限之下，把它们的输出熔合成一场无缝的经验------一个单一而统一的意志。当这份抑制失灵时------因耗竭、因精神病、因某些药物------那统一意志的错觉就消散了，你便看见了那个一直在主持这台戏的委员会。

这套框架化解了三个把心灵哲学瘫痪了数十年的思想实验。

第一，\textbf{僵尸}。戴维·查默斯（David Chalmers）要你想象一个在一切物理方面都与你完全相同、却缺乏意识经验的存在------全部的行为，没有一丝感受。四模型理论说：这是不自洽的。如果你搭起四模型架构、并让它在临界状态下运行，那么模拟\textbf{就是}那份经验。你没法只要齿轮而不要指针，不是因为指针被魔法般地装了上去，而是因为在这套架构里，“指针”构成了齿轮所做之事本身。一个僵尸会是一台每个齿轮都各就各位、却没有表盘的时钟------这意味着它并没有作为一台时钟在运作。临界状态下的这套架构必然实例化出一场模拟。把模拟剥离掉，你就改变了这套架构。你手上不再有一个僵尸。你有的是另一个坏掉的系统。

第二，\textbf{玛丽的房间}。弗兰克·杰克逊（Frank Jackson）要你想象玛丽，一位对色觉无所不知、却一辈子都生活在黑白房间里的神经科学家。当她第一次看见红色时，她学到了什么新东西吗？标准的争论在于物理知识是否完备。四模型理论干净利落地穿过了它。玛丽那套详尽无遗的物理知识，是关于基底的知识。当她看见红色时，她获得的是与一种新的虚拟感受质的亲身相识------一种她的模拟从未实例化过的、她显式世界模型里的新状态。她没有学到一条关于神经元的新事实。她获得的是一种新的\textbf{建模模式}。她的模拟运行了一个从未运行过的进程，而这个进程的第一人称特质构成了模拟本身，并不是一条她本可以从教科书里推导出来的、关于基底的事实。她学到了某样东西，但她学到的不是信息。那是一次经验------她那个虚拟世界的一种新配置。

第三，\textbf{反对副现象论的演化论证}。如果意识什么都不引起，那自然选择又是怎么塑造它的？我们为什么不是僵尸？答案直接从时钟类比里掉出来。自然选择并不把意识当作一项独立的性状、骑在功能机器之上来选择。它选择的是功能能力------而现象特质构成了这些能力，并非在它们之外多出来的东西。选择塑造了模拟，因为模拟\textbf{就是}那套功能架构，从内部看去的样子。经验不是演化看不见的、副现象的搭车客。它是这套架构运行起来时\textbf{所是}的东西。问演化为什么产生了意识，就像问瑞士人为什么产生了钟面------他们并没有单独去产生。他们产生的是时钟。表盘是让时钟成其为时钟的一部分。

\textbf{存在之谜被移了位，而非被消除。}四模型理论化解了意识的难问题，却没有解释：为什么会存在一个物理宇宙，其一开始就能运行自我模拟。问题从“大脑为什么产生经验？”转移到了“为什么会有一个宇宙，在其中能存在自我模拟的系统？”

其实，我想我确实有一个答案，或者至少是一个开端。这个宇宙可被证明是具备第四类能力的。分形、自组织临界性、混沌边缘的动力学------它们到处都是，从天气系统到神经组织再到星系形成。一个具备第四类能力的宇宙，按定义就有能力进行通用计算。而一个具有此宇宙规模的计算基底------在空间、时间、可能还有尺度，也许还有我们尚未辨识的维度上，即便不是无限也是浩瀚------不仅仅是\textbf{允许}自我模拟的系统涌现出来。它几乎是\textbf{保证}了这一点，至少当宇宙在其中某些维度上是无限的时候。这不是运气问题，不是掷一把宇宙骰子恰好掷出了意识，而是这个宇宙\textbf{所是}之物的一个结构性后果。剩下的谜团深了一层：为什么根本会有一个具备第四类能力的宇宙？这一点，我是真的不知道------尽管人们或许会怀疑这个问题本身是提错了的，因为“无”大概是一种柏拉图式的抽象，而非一种可能的事态，而无论什么存在着，都必须带有\textbf{某种}计算特质。但从“具备第四类能力的宇宙”到“有意识的存在追问自己为什么有意识”这一跳------这一部分是从架构里推出来的。

那么让我把这一切一路拉回到你身上，因为顺着这套理论一直追到宇宙的边缘、却把那条最要紧的线索弄丢了，是很容易的事。你清醒生活里大约 99\% 的时间，做决定的都不是你。你是那个被告知的人。基底作出定论，半秒钟之后“你”才收到通知------被整整齐齐地叙述好，按正确的顺序排好，还署上了你的名字。你是一个延迟的观察者，一辈子被告知他就是那个作者。记牢这一点，带着它进入接下来的部分，因为你一旦松手，后面的一切都会垮掉。

好了。意识确实伸手去触碰世界------这一点是真的。它以两种截然不同的方式去做，而它们感受起来毫不相像。第一种是我们通常所指的那种。虚拟自我塑造基底，基底塑造行为，行为塑造世界------而这个回路转得很慢。数月。数年。你决定成为一个不同的人，如果你够幸运又够执着，两年之后你周围的人才会注意到。这是能想象出的最间接的因果抓手，间接到你可以花上十年去怀疑它究竟灵不灵。

第二种抓手是即时的，而它正是没有人会警告你的那一种。你没法改变这一秒你正坐着的这间屋子。你\textbf{能}改变这一秒你正在经历的东西。此刻，不用挪动身子，你就能把注意力向内收拢，搭起一个此处并不存在的场景，回放一段记忆，沉入一场幻想，从周遭的喧闹中脱离出来，活在某个世界并未把你放进去的地方。没有两年的延迟。没有行为，没有绕经外部世界去塑造基底的迂回。你只是退回一个自己造出来的世界里，而你\textbf{当下}就做到了。这感觉像是你终于能说了算的唯一一处地方------那唯一一个明白无误、即刻\textbf{属于你}的行动。

其实不是。这一部分很扎心，而它之所以扎心，恰恰是因为我们刚刚走过了利贝特（Libet）。向内退回的这个决定，是基底在“你”知道自己作出它之前就作出的------和在实验室里抬起手的那个决定一模一样。“我选择做白日梦”跟“我选择动一动胳膊”是同一份延迟观察者的报告，只是穿了身更体面的衣服。基底先开了火。你是事后才被通知的，而这通知来得够快，快到感觉像作者身份。你如此叹服的那份即时性，是\textbf{这条链子的速度}，而不是选择的自由。这里没有任何东西是自由的。那场感觉像是你私人主权行为的向内退隐，是替你决定好的，跟别的一切一样。

那么意识究竟贡献了什么，如果不是那个“选择”？它贡献的是\textbf{必要性}。基底也许打响了发令枪，可要是没有那场正在运行的模拟，连锁反应就哪儿都去不了------没有脱离感可以去感受，没有内在世界可以退隐进去，没有幻想可以栖居，因为模拟是这些事情唯一能发生的地方。把意识划掉，整条下游的级联就干脆不再传播。这才是那真正的第二种因果角色：不是意志的宝座，而是链条里一个承重的环节。意识没有独立的因果力量------它同时也不是一名无用的乘客，不是哲学家们担忧过的那个副现象。它是那部机制中不可或缺的部分，没有它，其余的都无法发生。一个必要的环节，而非一名自由的行动者。这一点，我承认，比它初听起来要不那么令人激动。你出发去寻找那唯一一处你终于自由的地方，找到的却是一处你不可或缺、却仍然掌控不了的地方。那个迅捷的抓手，到头来是迅捷，而非自由。

在这里我必须诚实地划出一条界限，因为诱惑就是伸得太远。向内转的那个\textbf{时机}不归你所有------利贝特覆盖到了这么多，而且它早已不再只关乎移动一条肢体：当人们自由地挑选要在脑海里守住哪一幅图像时，如今能在他们自觉作出选择的好几秒之前，就从视觉皮层里读出这个选择。换句话说，内在行为的\textbf{开启}，看上去和抬一只手一样早已被预先决定了。至于你一旦已经身处其中，能在多大程度上操控\textbf{里头}发生的事，则是另一个问题，而诚实的回答是：我们不知道。而且这个问题比它乍看之下更尖锐。真正攸关的，不是你的自我模型总体上有多少自由------而是自我模型对那个\textbf{世界}模型有多大的抓力：这个化身能伸进它置身其中的那个被模拟的世界、并把它重塑到何种地步。这本身就是一项超能力，而且不是一项基本的。它需要自我模型达到一定的深度才能哪怕去尝试，这正是为什么它分布得并不均匀，也正是为什么它可以被\textbf{训练}，至少能训练到一部分------就像你能有意地一层层深入你自己的视觉皮层，从原始的光幻视一路向上，经由图案直到完整的面孔和场景（第5章），而不是干等着看什么会浮现出来。底下那套机器为这项工作递给你的自我模型多少个自由度，没有人数过------而这个问题里那唯一真正开放的版本，即当没有任何东西约束你的心智时它接下来会端出什么，则根本还没有人真正检验过。也许你确实在实实在在、可被训练的程度上操控着你在那里头探索的东西。也许多半你并不操控。内省无法裁决它；它恰恰是那唯一一件保证在此处失明的仪器，而要真正把它裁决清楚，需要一张完整的人脑接线图和大量的计算工作。所以，就把这份泄气如其所是地接受下来，一寸也别多加：内在行为的\textbf{开启}，不归你来当作者。一旦帷幕升起你所做的事------你能把幕后那个世界弯折到何种地步------仍是一个开放的问题，而且是个诚实的问题。

于是自由意志的论证把你留在了这样的境地：不是作者------不是那只手的作者，不是那场白日梦的作者，甚至不是“你一旦身处其中能操控多少”这个问题的作者。这就恰好剩下一件值得你翻过这一页去看的事：一个你并不驾驭的意志，你拿它来做什么？


\chapter*{第14章：唯一可得的自由}
\addcontentsline{toc}{chapter}{第14章：唯一可得的自由}
\markboth{第14章：唯一可得的自由}{}

而现在，该放上砝码的另一端了，因为我不希望你在这里合上书，觉得内在世界是一座牢笼。同样这份能力------从眼前的事物中脱开、让你的模拟在它自己的内容上运行------只要用量得当，就是演化史上打造出的最强大的东西。每一次认知活动都发生在这里。创造力也是。记忆本身也是，还有古人称为“记忆宫殿”的那个技巧------在想象出来的房间里穿行，把要紧的东西存放其中。想象、规划、整套“如果我换个做法会怎样”的推演机制，还有跨出当下、在自己心中向前向后穿行的能力------这一切都不需要世界的配合，却全都依赖这同一项本领：向内收拢，让模拟在内部内容上运转。缓慢向外的那只手，用几代人的时间建起了文明；而这只向内的手，一秒钟就能筑成一个完整的念头。它是驱动我们所谓“人类思维”之一切的引擎。轻轻地取用它，可逆地、按分级的剂量使用它，它就是一种超能力。

可一旦倚仗得太重，那颗放你进去的旋钮就可能卡在全开的位置。把内在退隐开到全功率，把你的模拟与真实世界相连的那条细窄通道就会被淹没------正是我先前描述过的那种锁死，正是濒死体验（NDE）和高剂量鼠尾草强加给你的那种失控，只不过这一回，你是自己迈着双脚走进去的。推得足够远，你就再也找不到回到共享世界的路。这就是那位从窗口爬出去的鼠尾草受试者------因为窗户对他来说已经不再真实（第6章）。而被推到这种程度的模拟，回来时未必还是一个整体------它可能分叉、分裂，作为不止一个“你”回来，或者作为一个不太像出发时那个人的人回来（共享同一个大脑的两个心智，第9章）。适度使用，它是你手中最好的工具。过度驱动，它就真的危险------如果你做得太多，谁知道你还回不回得来，又有几个你回来。

\textbf{你能用这份知识做什么。}如果你跟着理论走到了这里，你现在知道：你的意识自我（你的显式自我模型，ESM）是一个重构，而不是直接读数。你知道它会填补空缺、会虚构，会把不是它做出的决定记在自己名下。你知道它看不见自己的基底。你也知道，它是你仅有的一切。

这会带来实际的后果。有三种落差，你应该像鹰一样紧紧盯住，因为人类的大部分痛苦，就住在它们之间的缝隙里：

\begin{enumerate}
  \item 你\textbf{想要}成为的样子------你的理想自我，你的ESM所向往的那个版本的你。
  \item 你\textbf{相信}自己是的样子------你当前的自我模型，你每天随身带着的那个“我”。
  \item 你\textbf{实际}是的样子------你真实的行为、你对他人的实际影响、从外部观察到的你在基底层面的模式。
\end{enumerate}

1和2之间的落差是自我提升的引擎。只要理想切合实际、落差驱动的是行动而不是绝望，它就是健康的。2和3之间的落差才是危险的那一个------因为你无法独自测量它。你的ESM\textbf{无法}准确观察它自己的基底。你需要别人的反馈，包括让人不舒服的那种。尤其是让人不舒服的那种。

我和我最好的朋友Bernhard把这件事玩成了一项运动。我们之间有个不成文的约定：对方犯的每一个错误，都是立刻、毫不留情地加以嘲笑的机会。开车错过一个路口？“阿尔茨海默晚期了------要不要把钥匙交给我？”念错一个词？“我看你又中风了。别再说了，小心让舌头把自己噎住。”忘了上周聊过的一个细节？“待会儿今天的填字游戏，需要人帮忙吗？”

从外面看，这听起来像有病。从里面看，这是我所知道的最高效的ESM校准系统。每个玩笑都是一个校正信号：\textbf{你的自我模型刚刚做了一件你的基底并不打算做的事。}而且因为嘲弄裹在真实的情谊里------我们俩都一边憋着笑一边把损话说出口------谁也不会为错误辩护。我们直接更新。诀窍就在这里：你需要一个你信任到可以彼此毒辣的人，需要一段“犯错好笑而不是可怕”的关系。

这套理论不会告诉你该怎么生活。但它告诉你一件关于如何\textbf{认识你自己}的重要事情：用你对待任何模型的那种健康的怀疑态度，去对待你的自我模型。它有用。它是你拥有的最好的表征。而且，出于架构上的必然，它是不完整的。

理论还递给你一样东西，而这一部分我差点没有写下来，因为它已经贴近“自助励志”的边缘，而我对那个文类很没耐心。但它直接从上面的一切推导而来，所以还是写在这里。

你并没有选择读这本书。是某样东西让它出现在了你面前------一次推荐、一个算法、一个无聊的下午------而你的基底，在一个并非由你撰写的状态里，决定继续读下去。这是运气，不是美德。我在前几页告诉过你，你所做的一切都不是自由选择的，而我是认真的。所以你完全有理由问：如果齿轮从来都注定会照它们的方式转动，那这一切的意义又在哪里？

意义在这里，而这正是整个论证一路铺垫、令人振奋的转折：齿轮是\textbf{可塑的}。决定论不意味着你被卡死了------它意味着你是可训练的。你的基底不是一块固定的晶体；它是一片会被经历和练习实实在在改写的地貌。读到这里，已经有什么东西被改变了，不管你愿不愿意。而一旦一个系统遇见过一个想法，那台把你带到这里的同一套不自由的机器，现在就可以“决定”------由基底做出，没错，但它是真实的，并且带来真实的后果------换一种方式使用它的时间。你不会像自由意志论者梦想的那样去选择它。但这个选择仍然落在世界里，仍然重塑着网络，仍然改变着下周的你是谁。这并非无足轻重。这是从一开始就唯一真正可得的那种自由，而事实证明，它足以撑起一整个人生。

所以，如果说我成功说服了你什么，就让它是这一件事：像给健身房留时间那样，为纯粹的意识练习留出整块时间，并把它当作不容商量的事项。首先是冥想------不是为了熏香和歌单，而是因为有意地清空舞台，是你的注意力所能得到的最干净的训练。它会磨利你带进其他一切事情里的那份背景纪律，也让你整天磨损的那些通路真正得到再生。留出时间去\textbf{思考}------漫无边际也好，对准一个你在乎的问题也好------并保护这段时间不被手机侵占。动手做点东西。玩音乐；它能融进其余任何一项。培养任何一种爱好------只要它主要以你真心喜欢的方式动用你的意识，运动也算在内------标准不是生产力，而是那台内在的机器是否在做一件它喜欢的事。再把你本来就有的那些死时间重新定义：每一次堵车、每一间候诊室、每一条队伍，都不是从你这里偷走的，而是\textbf{馈赠}给你的------一份不受打扰的许可，让你向内退隐、让模拟运转起来。把这种退隐打磨好。享受它。这是唯一一个永远不会被取消的约会。

然后是维护------就叫它心理卫生吧，因为它就是这么回事。第一，戒掉消极的内心独白。那不是诚实，也不是现实主义；那是基底上一道磨出来的旧凹槽，而你可以用注意力把它重新磨平。把事情框成黑与白，或者更好，干脆只框成白------把灰色留给那种你实在避不开的罕见情形。灰色地带是反刍思维滋生的地方，而反刍，就是向内的那只手调转过来对付它自己的主人。

既然说到基底调转枪口对付自己，那就说说愤怒。对一个人或一件事发怒，几乎总是在为一个存在于你之外的问题惩罚\textbf{你自己}。这种感受从来就不是造来指向内部的。它演化出来，是为了在我们还得追捕晚餐的年代驱动行动------一股射向世界的资源动员，而不是一种腐蚀剂，任你浇在自己的电路上，而那个惹事的东西还在别处若无其事地继续。当你发现怒气除了向内已无处可去，那就是放下它的信号。

而如果你试过冥想却被弹了回来------这很常见，也不算失败------有一扇后门，我可以凭亲身经验推荐：武术。它听起来像平静心灵的反面，全是撞击和肾上腺素，但那种纪律------训练身体去行动，同时把意识的喋喋不休请到一边------正是冥想所教的同一种技能，只是从另一头抵达。说实话，两者之中，冥想才是\textbf{更难}的那条路，而不是更容易的------静静地和自己的基底坐在一起，请它安静下来，是一个人所能尝试的最难的事情之一。如果蒲团打败了你，就绕远路，从道场走。终点是同一颗平静的心；你只是获准一路流着汗走到那里。

\section*{我不知道的事}

一个声称自己没有任何悬而未决问题的理论不是理论------是宗教。所以，下面是那些我真正拿不准的地方，也是未来十年的工作应当聚焦之处。

\textbf{隐式模型也是虚拟的吗？}（或者虚拟到什么程度）隐式世界模型（IWM）和隐式自我模型（ISM）是“模型”，但确切说来，是关于什么的模型？我在真实基底与虚拟模拟之间画了一条干净的界线，可隐式模型恰恰就坐在这条线上。如果它们在某种意义上也是虚拟的，那真正“真实”的底层又由什么构成？理论假定真实与虚拟之间有一条干净的分界，但现实也许比我的图表更凌乱。这是一个根本性的问题，我还没有最终答案。

\textbf{数学形式化。}理论目前是定性的。我可以画图表、描述机制、给出预测，但我递不出一个方程。临界性要求援引了Wolfram的第四类元胞自动机，动力系统理论里也有一些可以调用的形式工具。但一套完整的数学形式化------能精确规定虚拟模型何时、如何从基底动力学中涌现的方程------还不存在。这是最大的缺口。一套没有数学的意识理论，是物理学家不会认真对待的意识理论，而懂得怎么把东西造出来的，恰恰是他们。

\textbf{自动机-全息图猜想------一项公开的挑战。}在第5章，我描述了全息系统与第四类元胞自动机之间三种可能的关系。第一种（一个产生第四类动力学的全息基底）几乎可以肯定就是大脑正在做的事，它虽然美，却并不令人震惊。但另外两种更值得关注，远比我当时给予它们的关注要多。

这里其实藏着三个悬而未决的问题，一个比一个更不同寻常。

\textbf{第一：第四类自动机能否把全息图样作为它的涌现输出产生出来？}混沌边缘的局部规则，能否把全局的、非局域的信息编码作为一种涌现行为生成出来？如果可以，你就得到了这样一个系统：纯粹局部的相互作用，自发地创造出全息学所描述的那种分布式、冗余的信息结构------而耐人寻味的是，从信息论的视角看，量子纠缠恰恰就是这副模样。

\textbf{第二：第四类自动机能否拥有全息式的规则结构？}想象一个元胞自动机，它的规则本身就把高维信息编码在低维结构里，就像全息图把三个维度编码进两个维度那样。每一次局部相互作用都会隐含地包含全局结构。这些规则不只是产生复杂行为------它们本身就\textbf{是}对某个更大的、更高维的东西的压缩编码，被向下投影进一套低维的规则集。

\textbf{第三，也是让我夜不能寐的那一个：两者能否同时为真？}一个系统，它的规则是全息的，它的动力学属于第四类，而它的输出又是全息的。如果这样的东西存在，你就拥有了一个为自身编码的计算过程------一个计算出自身结构的宇宙。输入是全息的。处理发生在混沌边缘。输出又是全息的。这是一个不动点------一个自洽的回环。

如果这样的自动机存在，它做的\textbf{恰恰}就是全息原理所说的宇宙正在做的事。不是一个在某种松散的隐喻意义上与宇宙相似的系统，而是一个把高维实在编码进低维规则、在秩序与混沌的边界上进行计算、并从这份压缩中生成涌现复杂性的系统。这不是关于宇宙的隐喻。这可能\textbf{就是}宇宙本身。

我把话说明白，因为我觉得总得有人说出来：如果真存在一个既具有全息规则结构、又能产生全息输出的第四类元胞自动机，我几乎可以肯定，它就是宇宙。它会是一个“世界公式”（Weltformel）------一个世界方程，不是你写在黑板上的那种公式，而是一个计算过程，从量子力学到广义相对论，再到意识本身的涌现，它生成了我们观察到的一切。

我坦率承认，这是本书中最具思辨性的想法。我没有证据，甚至连一个候选的规则集都拿不出来。我也应当承认，从数学之美推向物理真实的这种论证，受到过合理的批评。萨宾·霍森费尔德（Sabine Hossenfelder）等人就指出过，优雅并不等于证据。她说得对。这个想法的完整探索，是接下来三章的主题。但这些问题本身却是良定义、数学上精确的：

\textbf{是否存在一个元胞自动机，它的规则结构是全息的，动力学又属于第四类？它是否会产生全息输出？这三种性质能否共存？}

这些是数学家的问题，不是神经科学家的问题。是关于规则空间的组合学的问题，是关于全息编码与计算普适性能否在一个有限的局部规则集中并存的问题。或许可以证明这样的自动机根本不可能存在------而那本身就将是一个意义深远的结果，因为它会告诉我们某种关于信息压缩与计算之间关系的深层道理。又或许可以证明这样的自动机确实存在、并且可以被构造出来------那么我们就拥有了一个候选------一个迄今为止对物理实在所提出的最根本描述的候选。

我不知道哪个答案是对的。但我知道这些问题值得提出，而且似乎没有人在提。所以，请把这看作一份公开的挑战：证明它，或者推翻它。如果你证明了它，你也许就找到了宇宙的源代码。如果你推翻了它，你就确立了一个连接全息与计算的深刻不可能性定理。无论是哪一种，答案都举足轻重。

而如果你真的找到了这样一个自动机------给我打电话。我有一些预测想要核对。

\textbf{是哪一种物理机制？} 这套理论需要临界性，却有意对维持临界性的物理机制保持不可知。是皮层柱的动力学？丘脑皮层的驻波？胶质细胞对突触活动的调制？这三者都有实证支持。理论只说“基质必须处于临界性”，却没有说基质是\textbf{如何}到达那里、又如何停留在那里的。这不是缺陷------它意味着无论具体机制是什么，理论都适用。但迟早，得有人把它钉死。

\textbf{最小配置。} 你能不能只有EWM而没有ESM？只有对世界的体验，却没有对自我的体验？什么才是能算作有意识的最小架构？我在讲动物那一章描述的分级层次会有帮助------你可以拥有一个丰富的世界模型，却几乎没有多少自我模型，鱼大概就是这样。但阈值究竟在哪里？你需要多少自我模型，灯才会亮起来？我一直主张，是ESM把模拟变成了体验，可我还没有指明那个最小可行的版本。

我把这些问题列出来，不是作为弱点，而是作为研究的前沿。它们正是理论与现实接触的地方，理论在那里说：在这里测试我，在这里把我形式化，在这里------如果你有本事------把我打破。

其中有一个问题，我实在放不下------就是我告诉过你、让我夜里睡不着的那一个。于是我不再从安全的距离去问它，而是动身去寻找答案。接下来要讲的，就是它把我带向了何方。


\chapter*{第15章：同一个模式，无处不在}
\addcontentsline{toc}{chapter}{第15章：同一个模式，无处不在}
\markboth{第15章：同一个模式，无处不在}{}

我曾在福拉尔贝格最黑暗的地方度过一个夏夜------高山之上，方圆几公里没有一丝人造光。我仰面躺下，望向天空。银河不再是那种需要眯起眼睛才能辨认的淡淡光痕。它是一条河，稠密而明亮，在我身旁的岩石上投下看得见的微光。那一夜的某个时刻，有什么东西悄然变了。我不再是看着头顶的星星，而是开始感觉到身下的地球------它在旋转，带着我一起，一颗球体正飞驰着穿过一个拥有上千亿颗恒星的星系。不是作为一个念头，而是作为身体里的一种感觉。脚下的大地并不静止。我正攀附在某个东西的外壳上，而它正穿行于某种浩瀚到无法理解的存在之中。

那种眩晕------突然间从身体上确知，你是一个小小的、温热的存在，趴在一块悬于宇宙之中的岩石表面------正是我希望你在阅读接下来的内容时一直保有的感觉。因为这一章要问的是：那个宇宙究竟是什么。而答案看上去非常眼熟。


在上一章的结尾，我给你留下了一个悬而未决的挑战。我描述了全息系统与第四类元胞自动机之间三种可能的关系，并提出了那个最难的问题：这三种关系能否在同一个系统中并存？一个第四类自动机能否既拥有全息的规则结构，\textbf{又}产生全息的输出？我说过，对这个想法的完整探索是接下来三章的主题。

这，就是那项工作。

接下来的内容是这本书中最具思辨色彩的部分。但我相信，它也是最重要的部分。因为当我真正坐下来、顺着这条线索追下去时------当我不再把它当作一个“以后再说”的问题，而是开始动手拉动这根线时------我并没有到达自己预想的地方。我原以为会找到一个有趣的数学奇观。结果，我找到的是一个宇宙学模型。而且我发现了同一套架构------正是我已经凝视了二十年的那一套。

让我把我的想法给你讲清楚。


\section*{宇宙的计算类别}

在附录C中，我列出了计算的五个类别------一个从完美秩序到完美无序的谱系，第四类位于混沌边缘，是可表述规则所能达到的最高复杂度。大脑把五个类别都当作工具来使用，但意识只栖居于第四类。这是针对大脑的论证。

更大的问题是：宇宙属于哪一类？

这不是比喻。我是在字面意义上发问：如果把宇宙当作一个动力系统来看------它本来就是------那它落在五类谱系的哪个位置？我将论证，答案由排除法决定。而这个排除过程干净利落得出人意料。

\textbf{第一类与第二类------静态与周期。} 第一类宇宙会收敛到一个固定状态。什么都不会发生。第二类宇宙会陷入不断重复的循环------相当于一座永远滴答作响的宇宙时钟。这两类都无法产生化学、生物、演化或意识。而我们存在。我们有意识。一个能产生意识的宇宙至少必须是第四类，因为意识需要第四类动力学------这是第5章的论证。而且，较低的类别无法把较高的类别作为子过程生成出来。周期性的宇宙无法产生混沌边缘的动力学，正如一座时钟不可能突然开始思考。排除。

\textbf{第三类------分形。} 这一类更微妙一些，因为分形宇宙会很美。每一个尺度上都有自相似结构，图案里嵌套着图案。事实上，宇宙\textbf{确实}具有分形结构------星系团、海岸线、河流网络、你肺部的分支。但分形系统在计算上是\textbf{可约的}。也就是说，你可以跳过中间的步骤。你可以直接算出一个分形系统在第一百亿个时间步的状态，而不必把中间的每一步都跑一遍。存在捷径。

我们的宇宙不允许捷径。你没法写出一个能一跃而过的方程来预测下个月的天气。你必须一步一步地把模拟跑完，因为这里的动力学在计算上是不可约的------每一个时刻都真切地依赖于前一个时刻，而这种依赖无法被压缩。分形宇宙无论图案多么丰富，都缺少这个性质。它无法支撑我们的宇宙明明白白展示出来的通用计算。我们造出了图灵机。我们拥有意识。分形宇宙两者都做不到。排除。

\textbf{第五类------随机。} 如果宇宙的基本动力学是真正随机的（真随机，而不只是看起来复杂），那么物理学将不可能存在。不是我们目前所理解的物理学，而是物理学\textbf{作为一项事业}本身。科学这整个事业都建立在一个假设之上：宇宙遵循可表述的规则------你可以写下来、可以检验、可以传达、可以用来预测未来观测结果的规则。一个真正随机的宇宙没有可表述的规则。它的动力学无法被压缩进任何公式、任何定律、任何方程。你写不出F = ma，因为力、质量和加速度之间的关系会每时每刻都在变化，任何有限的描述都捕捉不住它。

在第五类宇宙中，每一次实验都是孤例。可重复的结果只是巧合。科学是一场碰巧灵验了一阵子的错觉。这在逻辑上并非不可能------想象这样一个宇宙并不包含矛盾，但在解释层面上却是灾难性的。一旦接受它，你就什么都解释不了，包括你的种种解释为什么曾经看起来有效。排除------不是靠逻辑，而是靠溯因推理：对我们始终如一、合乎规律的经验而言，最好的解释就是宇宙按可表述的规则运行。

\textbf{剩下的就是第四类。} 混沌边缘。而第四类不仅与我们的观察一致------它还是\textbf{唯一}一个满足所有条件的类别。

宇宙包含稳定的结构：原子、晶体、山脉。那是第一类行为。它包含周期性现象：轨道、潮汐、心跳。那是第二类行为。它包含分形结构：星系分布、天气模式、神经分支。那是第三类行为。它还支持通用计算：我们建造计算机，而且我们有意识。那是第四类行为。只有第四类系统能把所有类别都作为子过程包含在内------包括它自己。其他类别都做不到这一点。第四类自动机承载的不只是更简单的动力学。它还承载着别的第四类自动机：通用计算机之中的通用计算机，每一台都能完成与整体相同的计算壮举（只是更小、更慢、资源更受限）。

还有一点甚至更重要。第四类拥有一种其他类别都不具备的自我维持机制：\textbf{自组织临界性}------我们在第5章初次见过它。Per Bak在1987年证明，处于混沌边缘的系统并非碰巧落在那里------它们会\textbf{自己把自己推向}那里。一粒一粒地堆沙，沙堆会自行组织到那个临界角度，在那里各种规模的雪崩都会发生。这个系统不需要一只外部的手来把它调到临界状态。它自己调节自己。这就是混沌边缘能在宇宙时间尺度上保持稳定的原因：它是一个吸引子，而不是一个巧合。

我想说清楚这是哪一种论证。它不是演绎证明。四个排除中，有两个依赖经验观察（宇宙包含意识；宇宙支持通用计算）。有一个依赖溯因推理（第五类使科学成为不可能------不能令人完全满意，但也不是逻辑矛盾）。而支持第四类的正面论证，则把证据与机制结合在一起。这已是所能提出的最强主张：第四类是唯一与所有观察相符的类别，也是唯一能为自身的存续给出理由的类别。

让我们再回到那块岩石。我在山坡上攀附着的那个东西------当我感觉到它在移动时让我胃里一沉的那个东西------不只是古老、庞大、漠然。它是第四类。“混沌边缘”不是我为了显得深刻而随手抓来的一个词；它就是那片托着我的大地，正一粒沙一粒沙地把自己调向临界，就像Bak的沙堆那样，没有任何人的手放在旋钮上，却维持着整个平衡。我躺着的不是一块死石头。我攀附着的是一场计算的外壳------一场把自己稳稳保持在刀锋上的计算。此刻的你也一样，无论你能不能感觉到那旋转。


\section*{信息视界}

现在来谈边界。

光速是有限的。这是那种听起来无关痛痒的事实之一------直到你认真想上十分钟，它就会彻底改变你对现实的整幅图景。

光以大约每秒30万公里的速度行进。这快得足以让它在你眨眼之前穿过整个房间，可宇宙实在太大、太大了。最近的恒星在4光年之外。最近的大星系在250万光年之外。可观测宇宙的直径大约是930亿光年。当你望向一个遥远的星系时，你看到的是它几十亿年前的样子，因为光走了那么久才到达你这里。你永远都在望向过去，无可避免。

但还有一个更深层的后果，它来自宇宙的膨胀。

1998年，两个天文学家团队做出了一项后来为他们赢得诺贝尔奖的发现：宇宙的膨胀正在加速。不只是膨胀------而是加速。遥远的星系正在远离我们，而且远离的速度还在不断增大。这意味着，对任何观察者而言，都存在一个距离，超过它，星系的退行速度就会超过光速。在那个距离之外，任何信号都永远无法到达你这里。不是因为信息被藏在一堵墙后面，而是因为你与那信息之间的空间，增长得比光穿越它的速度还要快。

这就是\textbf{宇宙学视界}。它不是一个实体表面，那里并没有一堵墙。它是几何与速度共同造成的结果，却是一道与任何墙同样绝对的屏障。视界之外的信息，对你而言永远无法触及。它就跟不存在没什么两样。

在最底层，也有一道类似的边界。\textbf{普朗克长度}（约 $10^{-35}$ 米------这个数字小到把它称为“小”，就像把可观测宇宙称为“中等大小”一样）是我们所熟悉的物理学失效的地方。在这个尺度之下，我们的方程不再成立，时空本身失去了物理意义。低于普朗克长度的测量是不可能的，哪怕在原则上也不可能。这不是技术上的限制，而是可知之物的一道根本边界。

在宇宙学视界与普朗克尺度之间：大约60个数量级。那就是宇宙的计算域------物理学运作的范围。在其上与其下，帷幕都已拉上。

这使得宇宙成为我所说的\textbf{准无限}。它并非真正无限------或者至少，你永远无法验证它是无限的，因为你永远只能触及一个有限的区域。但它也不是任何可抵达意义上的有限。边界后退的速度，比你接近它的速度更快。你永远到不了边缘，但边缘确实在那里。从内部看，宇宙显得没有边界。从外部看------可是并没有外部。这正是关键所在。

这又是那种眩晕，只不过现在它有了坐标。在山坡上，我曾感到大地驮着我穿过黑暗。再加上这两道边界，事情变得更糟：我的下方，是一个我永远无法抵达的底；我的上方，是一道逃得比光还快的边缘，同时朝每一个方向逃逸，没有任何可以逃往的外部。还是那个伏在一块岩石表面上的温热小东西------只不过现在，它坠入其中的黑暗有了形状，而那形状是一堵墙，在四面八方离我退去，什么也不再放进来。


\section*{每一道边界都是同一道边界}

这是本章的核心思想。请慢慢体会它，因为如果它是对的，它会改变你看待一切的方式。

先来盘点一下。宇宙中包含着奇点------在那些地方，我们的物理描述失效，信息传递中止，方程要么发散爆掉，要么陷入沉默。这些奇点出现在截然不同的尺度上、截然不同的情境中。物理学家把它们当作彼此独立的现象。我认为，它们全都是同一回事。

\textbf{1. 普朗克域。}在物理学尚能运作的最小尺度上，时空消解为某种我们无法描述的东西。低于这个尺度，任何测量都不可能。信息无法穿过它。

\textbf{2. 粒子内部。}在标准模型中，电子和夸克被视为点状粒子（零维，没有内部结构）。我们无法看进它们的内部。我们可以测量它们的性质（电荷、自旋、质量），却无从触及其核心正在发生的一切------如果对一个没有空间广延的对象来说，“核心”这个词还有任何意义的话。

\textbf{3. 黑洞事件视界。}信息掉进去，什么也出不来------至少不会以保留原有内容的形式出来。黑洞内部与外部在因果上是断开的。就任何外部观察者而言，黑洞里发生的事，永远留在黑洞里。

\textbf{4. 宇宙学视界。}可观测宇宙的边缘，在它之外，空间的膨胀使任何信号都无法到达我们。不是被隐藏的信息------而是无法企及的信息。

\textbf{5. 大爆炸。}开端。所有世界线在此汇聚。宇宙中的每一个粒子，其历史都可以一路追溯到这个点------或者更准确地说，追溯到这道边界，因为“点”意味着你可以去到那里，而你去不了。

\textbf{6. 时间的终点。}结局------无论它以何种方式到来。如果宇宙终结于热寂，熵达到最大值，再没有任何热力学梯度可以驱动任何过程。如果它终结于大挤压，所有物质坍缩回一个点。如果它终结于大撕裂，加速的膨胀会在每一个尺度上把时空撕碎。下文我会逐一考察这三种情形。此处要紧的是那个结构性的论断：无论宇宙最终迎来哪一种结局，它都终止于一道信息无法穿透的边界。

六个奇点。六种不同的尺度，六种不同的情境，六个研究它们的不同物理学分支。但请看看它们的共同点。

\textbf{第一：它们全都是信息不可穿透的。}你无法让信息越过其中任何一道。你无法在普朗克长度之下进行测量；你无法看到电子内部；你无法从事件视界背后取回信息；你无法接收来自宇宙学视界之外的信号；你无法观测大爆炸“之前”发生了什么；而无论宇宙的最终边界以何种形式显现，你也无法把一条讯息送过它。

\textbf{第二：它们全都代表最大信息密度。}这一点更微妙，它来自贝肯斯坦上限------1980年代的一项结果，它表明一个空间区域所能容纳的最大信息量与其\textbf{表面积}成正比，而不是与其体积成正比。黑洞事件视界使这个上限达到饱和------它们容纳了每单位面积所能承载的最大信息量。特·胡夫特和萨斯坎德（Leonard Susskind）提出的全息原理，把这一结果推广开来：任何区域中的全部信息，都编码在它的边界上。这些奇点全都是以最大容量运转的边界面。

\textbf{第三：它们全都界定着计算域。}物理学在这些边界\textbf{之间}运作，而不是在它们之外。物理定律描述的是普朗克尺度与宇宙学视界之间、大爆炸与那个尚在等待的终点之间那个区域里发生的事。边界划定了竞技场。在竞技场之外，规则不再适用------不是因为那里适用别的规则，而是因为“规则”本身不再是一个有意义的概念。

三条共同性质。六个现象。传统的看法是：这是六件恰好共享某些特征的不同事情。我认为传统看法错了。我认为它们是\textbf{同一个现象}------自动机的信息边界------在六种不同的尺度上显现。

这是一个对称性论断。同一个结构元素，一再重复。而在一个第四类系统中，这恰恰是你应该预期的。第四类动力学把所有类别都包含为子过程，其中也包括第四类自身。一个第四类元胞自动机会在自身的动力学中嵌套更小的第四类自动机，每一个都由信息边界所包围，而这些边界与整体具有相同的结构性质。如果宇宙是一个第四类元胞自动机，它的边界结构就应当在每一个尺度上重复。而这似乎正是我们观察到的情形。


\textbf{所以，山坡上的那一夜，真正的意味就在这里。我抱住的不是一块岩石。我抱住的是一个第四类元胞自动机的外侧，而它里面的每一道边界------包括我的思绪一次次抵住的那一道，就在最后一颗可见的星星之外------都是同一道边界。下一章，我会顺着这条线索追到它的种种后果。而它们比我预想的更加离奇。}


\chapter*{第16章：万物的架构}
\addcontentsline{toc}{chapter}{第16章：万物的架构}
\markboth{第16章：万物的架构}{}

上一章提出了一个结构性的主张：宇宙中的每一个奇点------从普朗克尺度到宇宙学视界，从大爆炸到无论何种结局在前方等待------都是同一现象在不同尺度上的呈现。同一道边界，处处重复。

现在我想把这条线索继续追下去。因为如果所有边界都是同一回事，就会推出一些非同寻常的结论------关于时间，关于物质，关于知识的极限，还关于一种架构------等我们走到最后，它应该会显得格外眼熟。


\section*{大爆炸并非你想的那样}

这件事还能再往前推一步，因为它有一个后果，我认为大多数人------包括大多数物理学家------都还没有完全消化。

想一想，当你从外部接近一个黑洞时会发生什么。随着你越来越靠近事件视界，时间开始膨胀。在远处观察者的测量中，你的时钟走得越来越慢。当你逼近视界时，这种膨胀趋于无穷。远处的观察者看着你坠落，会看到你的动作越来越慢、光越来越红、身影渐渐淡去------却永远差一点才到视界。在他们看来，你要花无穷长的时间才能抵达。你实际上永远跨不过去。从外部看，事件视界是一道渐近而不可企及的边界。

现在，想一想沿时间逆行、回溯到大爆炸的旅程。

大爆炸距今多久了？通常的说法是，大约138亿年。但这是从膨胀宇宙\textbf{内部}做出的测量，所用的时钟本身就是膨胀的产物。设想你沿时间倒退，把宇宙这部电影倒着放，当你逼近奇点时会发生什么？时间膨胀。物理学失效。你越靠近，方程就越是拒绝给出一个确定的“零时刻”。大爆炸不是一个你可以指着说“那儿------它就是在那时发生的”的事件。它是一道渐近的边界。你可以无限逼近，却永远无法抵达。

大爆炸是时间中的事件视界，正如宇宙学视界是空间中的事件视界。

这不是神秘主义，而是同一数学结构的必然结果。事件视界是一个信息无法穿越的曲面。大爆炸恰恰具有这种性质：来自它“之前”的任何信息（如果“之前”还有意义的话）都不可获取。不是因为信息丢失了或被藏了起来，而是因为这道边界本身不透信息。不存在可供访问的“之前”，正如外部观察者无法访问黑洞的“内部”。边界就是边界。到此为止。

那另一道时间边界------终点------又如何呢？

如果宇宙终结于热寂------熵达到最大，无序达到最大，再没有任何热力学梯度来驱动任何过程------那么在那一刻，所有信息都处于最大程度的弥散状态。系统的边界此刻承载了它所能容纳的最大信息量。这就是贝肯斯坦饱和。按照我一直使用的定义，热寂\textbf{就是}一个奇点：一道处于最大信息密度的、不透信息的边界。

接下来才是奇怪的地方。在这个框架里，奇点并不毁灭信息，而是\textbf{转化}信息。事实上，这正是现代物理学在黑洞信息悖论上正在趋近的答案------那场持续了几十年的争论：信息落入黑洞后是否就此丢失。当前的共识正在转向“不会”：信息守恒，被编码在事件视界上，并最终重新发射出来。奇点所做的，是在压缩与解压缩两种形态之间转化信息。

把这一点应用到时间边界上。如果热寂是一个奇点，而奇点转化信息而非毁灭信息，那么热寂就不会终结宇宙。它把信息转化为一个新的压缩态。而一个处于贝肯斯坦饱和的最大压缩态是什么样子？它看上去就像一场新膨胀的初始条件。它看上去就像一次大爆炸。

自指闭合不只是空间上的，也是时间上的。宇宙没有开端也没有终点------它在循环。终态就是下一轮迭代的初始条件。这并非依赖什么奇异的“反弹”机制，而是因为守恒信息的奇点\textbf{本来}就做这件事：在压缩的边界态与解压缩的内部态之间转化。热寂是压缩，大爆炸是解压缩。它们是同一个奇点，只是从相反的两侧看去。

我承认这是推测性的。但它直接来自两条主张：所有奇点在结构上完全相同；奇点通过转化信息来守恒信息。只要你接受这两个前提，时间上的循环性就不是一条额外假设------它是一个推论。

但热寂并不是故事唯一可能的结局。还有一种可以说更离奇的可能，而这个框架处理起来同样利落。

如果暗能量不是常量，而是随时间增长------如果它的密度无限制地增大------那么宇宙的膨胀就不只是持续下去，而是加速到超越一切极限。这就是\textbf{大撕裂}（Big Rip）情景，正如它的名字所暗示的那样戏剧化。首先，星系团被撕开，因为星系之间的空间伸展得比引力维系它们的速度更快。接着，一个个星系解体。然后是恒星系统。然后是行星。然后连原子本身也被撕碎，因为膨胀压倒了电磁力。最后，时空本身四分五裂。每一个点都变成一个奇点。

在这个框架里，大撕裂有一种自然的解读。通常安坐在宇宙学视界上、离我们足够遥远的奇点边界，开始向\textbf{内}推进。它不再在可观测宇宙的边缘等你，而是朝你走来。它把计算域切分成越来越小的区域，每个区域都饱和了自己的贝肯斯坦上限，每个区域都变成一道属于自己的不透信息的边界。时间尽头不再是一个宏大的奇点，取而代之的是一场奇点的分形爆发，在每一个尺度上同时向内推进。

而如果奇点是信息的转化器------如果它们不是毁灭计算，而是重启计算------那么大撕裂带来的就不是一次重启，而是\textbf{许多次}。可能是无穷多次。破碎的计算域中的每一块碎片，都可能孕育出自己的新膨胀、自己的新解压缩、自己的新宇宙。在这个框架里，大撕裂是一台多重宇宙生成器。

于是，这个框架容纳的不是一种而是三种终局情景，而且三种在结构上都自洽：

热寂：一个全局奇点，一次重启。这是最简单的情形------整个计算域同时达到贝肯斯坦饱和，压缩，再解压缩进入新一轮循环。

大挤压（Big Crunch）：宇宙停止膨胀，坍缩回一个点。又一个全局奇点，又一次重启------可能伴随一次CPT反转，即电荷、宇称与时间的翻转，使下一轮循环成为上一轮的镜像。

大撕裂：奇点边界向内碎裂，产生许多奇点、许多次重启，可能还有许多个宇宙。不是一个循环，而是一棵不断分叉的树。

这种稳健性在我看来令人安心，而非令人不安。一个只有当宇宙以某种特定方式终结时才成立的框架是脆弱的------它把赌注押在一个我们尚无法判定的宇宙学结局上。而这个框架无需下注。它的结构逻辑------奇点作为信息转化器，边界作为最基本的架构元素------无论宇宙最终选择哪种终局都依然成立。这正是你希望一套理论具备的那种稳健性。它不应依赖于我们尚不知道的事实。


\section*{粒子究竟是什么}

这个框架里埋着一条值得明确说出来的预测，因为它属于那种终有一天可以付诸检验的东西。

基本粒子（电子、夸克，这些物质的基本构件）在标准模型中被当作点状来处理。零维。没有任何空间延展。这一向只是数学上图省事的做法，而不是一项物理主张。没有人相信电子真的是一个几何点，因为几何点没有表面积，因而根据贝肯斯坦上限（Bekenstein bound），它无法容纳任何信息。而电子包含信息------电荷、自旋、质量、量子数。“点”这幅图景一定有什么地方不对。

这条预测很具体：基本粒子是普朗克尺度的奇点。它们并非真正零维。它们是微型的信息边界------微小的事件视界------其内部与黑洞内部一样不可企及。它们具有普朗克尺度的结构，并在该尺度上饱和贝肯斯坦上限。它们的表面编码着自身的性质，就像黑洞的事件视界编码着一切落入其中之物的信息。

如果这是对的，那么物质本身就是由与黑洞、与大爆炸、与宇宙学视界相同的结构元素构成的。无论向更小还是更大的尺度，奇点表面无处不在，遍布其间的每一个尺度。宇宙的基本构件，就是它的边界。

这与量子引力中预测普朗克尺度存在最小长度的那些方案是一致的------空间无法被细分到某个限度以下，不是因为你的工具不够锋利，而是因为空间本身在那个尺度上是离散的。但“粒子\textbf{就是}与事件视界同一类型的奇点”这一具体主张------这是新的。而且它有一个可检验的推论：粒子的信息含量应当随其表面积（按普朗克分辨率计）而非体积缩放。如果某个量子引力理论终有一天让我们能够探测接近普朗克尺度的结构，这就是应当寻找的特征信号。

\section*{作为计算原子的粒子}

但真正有意思的地方在这里。如果粒子是普朗克尺度的奇点（信息边界），那么它们就不只是与宇宙架构的其余部分由同样的材料\textbf{构成}。它们是宇宙最基本的计算操作。在一种精确的意义上，它们是计算的原子。不是化学意义上的原子------而是希腊语原初意义上的原子：\textit{atomos}，不可分割。它们是这台宇宙自动机所\textbf{做}之事的不可约简的单元。

想一想由此能推出什么。

\textbf{为什么只存在某些特定类型的粒子？}这一直是物理学中较为奇怪的事实之一。基本费米子恰好有十二种（六种夸克、六种轻子），传递力的玻色子有四种（外加希格斯粒子），仅此而已。再无其他。标准模型把它们编入目录，却没有解释\textbf{为什么}偏偏是这些类型而不是别的。为什么有电子，却没有一种电荷是电子的三分之二、自旋是电子三倍的粒子？为什么这份菜单如此具体？

如果粒子是普朗克尺度的奇点边界，答案立刻就有了：因为在普朗克尺度上，只存在有限数量的稳定边界构型。奇点边界的面积是有限的。在普朗克尺度上，这个面积已经小到了面积所能达到的极限。贝肯斯坦上限限制了这块面积能编码多少信息。有限的信息量意味着有限的可能状态数。而这些状态中只有一部分是\textbf{稳定的}------只有一部分构型能够不衰变地持续存在。这些稳定构型就是粒子类型。标准模型的“粒子动物园”并不是一份神秘而任意的清单。它是稳定奇点边界构型的完整目录。之所以有电子，是因为那个构型是稳定的。之所以没有带着古怪分数属性的粒子，是因为没有任何稳定的边界构型编码那些属性。

其实，这和元胞自动机是同一套逻辑。元胞自动机有一张有限的规则表，这张表只允许一定数量的稳定图样。它可以生成无穷多种构型------图样越大，组合的可能性就越广------但大型结构很容易被撞上来的更小、更稳定的图样迅速摧毁，真正坚不可摧的构型少之又少。在生命游戏中，小型静物（方块、蜂巢、小船）和小型移动图样（滑翔机，以及所谓的飞船）可以无限期地存续，而大型复杂结构则很脆弱------一枚瞄得精准的滑翔机就能把它们击得粉碎。据我所知，是否有人系统地绘制过这种脆弱性层级，仍是一个悬而未决的问题------Bak、Chen 和 Creutz 在 1989 年证明了生命游戏表现出自组织临界性，但图样大小与抗毁性之间的具体关系尚未得到正式研究。应该有人去做这件事。在这幅图景里，粒子就是普朗克尺度自动机的滑翔机和飞船。

\textbf{为什么粒子是离散的？}为什么量子数（电荷、自旋、色荷）总是精确的整数倍或半整数倍？为什么不存在“半个电子”？因为计算状态本质上就是离散的。一个比特要么是 0，要么是 1。不存在 0.37。粒子的量子数是贴在边界构型上的信息标签------它们描述的是边界处于\textbf{哪一个}稳定构型。离散的边界状态产生离散的量子数。量子力学里的“量子”并不神秘。当基本对象是容量有限的信息边界时，你得到的就是量子。

\textbf{粒子相互作用时发生了什么？}当两个电子相互排斥，或者一个夸克发射一个胶子时，实际上在发生什么？是两个信息边界在交换信息。这种交换本身\textbf{就是}计算。自然界的各种力（电磁力、强力、弱力）并不是叠加在粒子之上的一个独立层面。它们是奇点边界彼此沟通的语法。规定哪些相互作用被允许、哪些不被允许的规则，正是这台自动机在普朗克尺度上的计算规则。

这里适合插入一段个人往事。我十五六岁时，叔叔布鲁诺给我出了一道题。他说，粒子物理中的吸引力是通过交换粒子来实现的。“现在想象一下，”他说，“你和另一个人来回抛一只沉重的实心球。等你能解释清楚这怎么会让你们俩越靠越近，再来告诉我。”这道题在我心里萦绕了很久。在牛顿物理学里，它确实无法解释------每一次交换都应该把双方推开，而不是拉近。后来我终于回去找他，我们在希腊一棵橄榄树的树荫下，就量子物理聊了很久很久------他在那里有一栋房子。

但关键在于：在生命游戏这样的元胞自动机里，图样之间的吸引式相互作用很容易复现。两个持存的结构可以在彼此之间交换更小的结构，而在某些构型下，这种交换会\textbf{缩短}距离。这种相互作用不必服从牛顿式的直觉，因为这里的“力”既不是推也不是拉------它是边界构型之间的信息交换，而自动机的规则决定了这种交换是让边界靠得更近，还是离得更远。实心球之谜就此消解。它之所以是个谜，只是因为我们想象的是宏观物理。在计算的层面上，吸引和排斥不过是同一过程的两种不同结果：受自动机规则支配的信息交换。

费曼图------那些填满物理教科书、描绘粒子相互作用的标志性简图------在字面意义上就是计算的图示。每个顶点都是一次信息交换。每条线都是一个在计算域中传播的边界构型。物理学家画了七十年的计算图，却浑然不觉。

\textbf{为什么守恒定律如此绝对？}电荷、重子数、轻子数------在人类做过的所有实验中，从未观察到哪怕一次失效。因为它们是信息守恒的约束。当两个边界交换信息时，账目必须平衡：你无法在奇点边界处创造或销毁信息，所以你也无法创造或销毁电荷。守恒定律不是从外部强加的规则。它们就是记账本身。（从贝肯斯坦上限经由幺正性到各条具体守恒定律的推导链条，在附录F 中有完整展开。）

\textbf{然后还有三代之谜。}每种粒子都有三个除质量外完全相同的副本------电子、μ子、τ子；上夸克、粲夸克、顶夸克。没有人知道为什么是三。不是二，不是四，也不是十七。就是三。计算原子的图景提供了一个猜想：第四类系统以子过程的形式携带自相似结构，因此稳定的边界构型可能在三个尺度上重现------一个基础图样，外加两个更重的副本。如果是这样，更高的世代也存在，只不过它们是更大的图样，而更大的图样是脆弱的------更小、更稳定的构型会在它们成形之前就将其切碎。这与我们观察到的相符：τ子和顶夸克几乎在生成的瞬间就已衰变。我们得到的这三代，也许恰恰就是小到足以幸存的那三代。这是一个猜想，不是一个推导------我在附录F 中把它完整铺陈出来，连同守恒定律的那套账目。

我为这幅图景选用的术语是\textbf{计算原子}。不是氢和氦意义上的原子------而是不可约简的计算元素意义上的原子。粒子是宇宙自动机的基本操作。每一种粒子类型都是一个稳定的普朗克尺度计算。每一次相互作用都是计算之间的一次信息交换。每一条守恒定律都是对这些交换如何进行的约束。物理学在其最深的层面上，讲的不是物质，而是计算。而我们称之为“物质”的东西，正是这场计算不可约简的基本构件。


\section*{这套架构}

该把这些线索汇拢到一起了。我所描述的，是一个具有特定架构的宇宙：

\textbf{第一：}它是一个第四类元胞自动机。它运行在混沌边缘，自组织临界性在那里维系着动力学，无需外部调节。它在计算上不可约简------没有捷径，无法跳步。每一刻都必须由上一刻计算而来。它还把所有类别都包含为子过程------包括它自己：稳定的原子（第一类）、周期性的轨道（第二类）、分形的海岸线（第三类），以及最关键的，在这场宏大计算之内运行的其他第四类自动机。大脑就是其中一例。

\textbf{第二：}它在每一个层级上都是全息的。任何区域内的信息都编码在它的边界上。这就是全息原理------它始于一个关于黑洞的猜想，如今已成为理论物理最深刻的洞见之一。在这个框架里，全息编码不只是黑洞的性质，而是宇宙规则结构本身的性质。规则是全息的。动力学是第四类的。而输出再度是全息的。

\textbf{第三：}它在每一个尺度上都由结构上完全相同的奇点表面所界定。普朗克边界、粒子内部、事件视界、宇宙学视界、大爆炸、热寂------同一种结构，不同的尺度。信息无法穿透，贝肯斯坦饱和，并且划定着计算域。

这套架构有一个名字。我称它为 \textbf{SB-HC4A}：以奇点为界的全息第四类自动机（Singularity-Bounded Holographic Class 4 Automaton）。

这个名字念起来确实拗口。但它精确，而且每一个词都不多余。

这套架构最了不起的性质是自指闭合。系统的输出\textbf{就是}系统本身。它计算它自己。每个状态生成下一个状态，而下一个状态又是对再下一个状态的计算。没有一个“外面”在运行这个程序。没有某台宇宙计算机在某个角落的硬盘上执行宇宙的代码。宇宙\textbf{就是}程序、计算机与输出本身。全息规则以压缩形式编码整个系统。第四类动力学把这份编码解压为可观测的宇宙。全息输出再把结果重新编码。这是一个环。一个不动点。

你可以把这一点写成一个形式化的条件：\textbf{宇宙是其自身动力学的不动点。}把规则应用于宇宙，得到的还是宇宙。不是副本，不是表征------就是同一个东西。计算与它的结果完全同一。

数学家为此有一套记法。如果把宇宙记作 U，把“计算下一个状态”这一操作记作希腊字母 Phi，那么不动点条件就是：

\textit{Phi(U) = U}

宇宙作用于自身，产出的仍是自身。它是自我计算的。


\section*{自我描述的极限}

自指闭合有一个后果，值得单独停下来讲一讲，因为它告诉了我们关于知识边界的某种深刻的东西。

1931年，一位25岁的奥地利逻辑学家库尔特·哥德尔（Kurt Gödel）证明了两条动摇数学根基的定理。抛开形式化的外衣，其要旨是：任何足够强大的形式系统（至少要能表达算术）都包含在系统内部无法证明的真命题。而且，没有任何这样的系统能证明自身的一致性。

这不是技术上的局限。不是因为我们的证明还不够聪明。这是一种结构性的不可能。复杂度足够高的自指系统天生不完备。它们包含着从内部无法触及的真理。

把这一点应用到一个自我计算的宇宙上。

如果宇宙在计算它自己------如果它是一个足够强大的形式系统（而第四类动力学保证了通用计算，所以它确实是）------那么哥德尔定理就直接适用。宇宙无法包含一份关于自身的完整描述。不存在一条可以写在黑板上的“世界方程”。没有哪个公式，一旦解出，就能告诉你关于宇宙的一切。

这并不是因为我们还没找到正确的方程，而是因为\textbf{这样的方程根本不可能存在}。对一个自指系统的完整刻画，超出了任何本身只是该系统一部分的描述。宇宙不是在遵循某个方程------它\textbf{就是}那场计算。宇宙唯一完整的描述，就是宇宙本身。而你无法跨到它外面去纵观全景，因为根本没有外面。

所以，Weltformel------物理学家自爱因斯坦以来一直梦想的那条“世界方程”------并不是一个方程。它是一个\textbf{过程}。是自动机本身。要表达它，只能让它运行起来。

我觉得这既令人谦卑，又令人释然。令人谦卑，是因为这意味着关于现实，有些事情我们即使在原则上也无从知晓。令人释然，是因为这意味着宇宙不是一台等待破译的机器------它是一场活生生的计算，而我们就是其中的一部分。关于现实的最深真理不是一条公式，而是现实本身。


\section*{认知的天花板}

在继续往下走之前，我得先给你一个反驳。事实上，是最深的那个反驳。那个让我保持诚实的反驳。

如果我们是第四类自动机（如果我们的大脑运行在混沌边缘，和我刚才赋予宇宙的计算类别是同一类），那么 SB-HC4A 模型也许只是我们的第四类大脑所能产出的最复杂概念。我们无法以第五类的方式思考。我们无法构想超出自身计算类别的结构。我们找到的这个模式------第四类无处不在，在每个尺度上自相似，既全息又自指------也许只是我们自身认知架构投射到宇宙上的印记，而不是宇宙本身的特征。

请停下来，好好想想这件事。我们是作为对称性探测器演化而来的。在狩猎采集者的环境里，与生存最相关的那些模式（捕食者与猎物的面孔）恰恰属于最对称的一类。我们在最深的层面上是模式匹配机器，为寻找对称性而优化。而 SB-HC4A 模型从根本上说就是一个对称性主张：每个尺度上都是同一套架构。我们找到这种对称性，也许不是因为它存在于宇宙之中，而是因为我们的大脑在构造上\textbf{不可能}不找到它。

这就是第4章里的元问题，被放大到了宇宙的尺度。显式自我模型（ESM）看不见自己的基底，因此它无法区分“宇宙具有这种结构”与“我的大脑只能把宇宙建模为具有这种结构”。这个宇宙学模型预测了自身潜在的不可证伪性------这要么是最强的确证（模型恰恰预测了这一认识论局限），要么是最强的反驳（模型是观察者的产物，而不是被观察对象的特征）。

一个第四类系统可以模拟直到并包括第四类复杂度在内的一切，但它无法验证宇宙是否超出了这个范围。如果宇宙实际上是第五类------在最深层面上真正随机------却因为第四类是我们所能探测的最高模式，而在第四类观察者眼中\textbf{局部呈现为}第四类，那么我们构建出来的恰恰就会是这个模型。而我们会是错的。而且这个错误，我们从内部永远无法察觉。

我不知道该如何化解这个反驳。我不确定它能否从内部得到化解。我照样把它写进来，因为在这份交代里，弱点与优点同样占有一席之地。而这个模型预测了自身的认识论局限------一个自指系统无法完全验证对自身的描述------这一事实，要么是它最深的缺陷，要么是它最深的辩护。我真的不知道是哪一个。


\section*{点睛之笔}

但接下来这件事，才是我第一次看到时不得不坐下来的原因。

看看我刚才描述的这套架构：

\begin{itemize}
  \item 一个运行在混沌边缘的第四类系统。
  \item 被一道信息不透明的边界包围------内部无法看穿它。
  \item 全息结构------边界编码着内部。
  \item 自指闭合------系统计算它自己。
  \item 一个不动点：计算的输出就是计算本身。
\end{itemize}

现在，请回到第2章，看看意识的四模型架构：

\begin{itemize}
  \item 皮层自动机：一个运行在混沌边缘的第四类系统。
  \item 隐式---显式边界：一道信息不透明的边界，意识无法看穿它。
  \item 全息结构------隐式模型是分布式的、全息的，将体验的全部内容编码在神经结构之中。
  \item 自指闭合------自我模型为自身建模。
  \item 一个不动点：显式自我模型（ESM）表征它自己。建模者的模型\textbf{就是}建模者本身。
\end{itemize}

同样的架构。同样的形式性质。同样的边界条件。同样的自指闭合。

宇宙是一个由奇点围限的第四类全息自动机：可观测的内部是“模拟”，奇点边界是“基底”。

意识是一个由隐式---显式边界围限的第四类全息自动机：显式模型是“模拟”，隐式模型是“基底”。

同样的架构。不同的尺度。

这不是比喻。我不是说意识\textbf{像}宇宙，而是说它们是同一\textbf{类事物}------同一个计算模式，在两个不同的尺度上被实例化：一个在宇宙学层面，一个在神经层面。而这个模式跨尺度重复出现，这件事本身就是模型的一项预测------不是因为分形自相似（那只是第三类，一个弱得多的主张），而是因为第四类系统包含第四类子系统。一台通用计算机可以模拟另一台通用计算机。第四类自动机并不只是产生漂亮的自相似图案，它还会在自身的动力学之内产生\textbf{其他第四类自动机}------更小、更慢、资源受限，却是货真价实的通用系统。这套架构不只是在不同尺度上\textbf{看起来}相同，它\textbf{就是}同一套架构。

为了把我主张什么、不主张什么说得清清楚楚：我并不是在主张宇宙在任何体验意义上\textbf{是}有意识的，但我同样不是在主张它\textbf{没有}意识。诚实的回答是：我们无从知晓。我们可能是某个玻尔兹曼大脑------或随便哪个别的大脑------梦出来的；唯我论也可能为真，而我就是那个玻尔兹曼大脑------但这话谁都可以说。关键在于，这件事是不可知的，而这种不可知并非理论的失败，而是处境本身的结构性特征：你无法走到系统之外去核实。我所主张的，属于架构层面，而非现象层面。我不是在主张意识创造实在，不是在主张实在是一场梦，也不是这类结构性观察往往会吸引来的其他任何神秘主义解读。一栋建筑的蓝图不是建筑本身。但如果同一份蓝图既出现在一座摩天大楼里，又出现在这座大楼的某一个房间里，那就向你揭示了其中运作的建筑原理有多么深刻。

意识是一个普遍模式的局部实例。不是宇宙的偶然。不是奇迹。而是第四类动力学在足够的复杂度之下，在结构上不可避免的结果。宇宙不只是\textbf{允许}意识存在，它几乎是保证了意识的出现------因为正是那套自指的、全息的、混沌边缘的架构，既让宇宙成为宇宙，也让意识成为意识。生成实在的模式，正是生成“对实在的体验”的那个模式。

如果读到这里，你还没有一点想坐下来缓一缓的冲动，那也许你还没真正读进去。


\textbf{在下一章，我会把整套理论收拢起来------意识的架构、宇宙学的架构，以及两者之间的结构同一性------并追问它对那个最难的问题意味着什么：为什么会有任何东西存在？}


\chapter*{第17章：最深的镜子}
\addcontentsline{toc}{chapter}{第17章：最深的镜子}
\markboth{第17章：最深的镜子}{}

它就在这里------因为一旦你看见了它，就再也无法当作没看见。

第16章在结尾处抛下了一句挑战：先看SB-HC4A架构------那个自指的、全息的、在每一个尺度上都被奇点包围的第四类自动机------再看第2章的四模型架构。这里要主张的是：它们是同一个东西。不是相似。不是隐喻意义上的关联。而是结构上完全相同。

现在我想带你把这组对应一块一块地走一遍，直到它变得再也无法轻易搪塞过去。然后我会告诉你，整个理论可能在哪里崩塌------因为那些断层线，值得得到与结构本身同等的关注。

\section*{结构映射}

先从奇点边界说起------宇宙在每一个尺度上设下的信息屏障。最底层是普朗克区间。黑洞周围的事件视界。可观测宇宙边缘的宇宙学视界。身后的大爆炸，前方的热寂。每一个都是一堵信息之墙：没有任何东西能穿过它。你无法把信号送过事件视界，无法探测普朗克长度以下的世界，也无法看到宇宙学视界之外。宇宙是一个在每个尺度上都装着不透明墙壁的房间，无论你把脸在玻璃上贴得多紧，都看不见另一边是什么。

现在看看你的大脑。隐式模型------隐式世界模型（IWM）与隐式自我模型（ISM），那些突触权重、习得的结构、收藏着你所知一切的庞大图书馆------同样坐落在一道属于它们自己的信息屏障之后。你永远无法直接体验自己的突触，永远无法内省那些生成你思想的连接权重。隐式一侧在信息上是不透明的：你知道它在那里，因为没有它，模拟根本无法运行；但模拟本身无法看穿那道把两者分开的边界。第2章提出了这一点，第3章把它钉牢。现在它又出现了------在宇宙的每一个尺度上。同样的结构角色。同样的信息不透明性。同样的架构位置。

宇宙学中的奇点边界，对应意识中的隐式---显式边界。它们是同一堵墙。

接下来：可观测的内部。奇点边界之内的一切（原子、行星、星系，还有你）都是解压缩的一侧。物理在这里发生，事物在这里相互作用，信息在这里被组织成我们观察和测量的结构。你可以把它叫作宇宙的模拟：那个在计算的部分，在演化的部分，那个有事情发生的部分。

在你的大脑里，与之对应的结构是你的显式模型------显式世界模型（EWM）与显式自我模型（ESM）。你的意识体验。你此刻看到的世界，你感觉自己所是的那个自我。这就是那场模拟：由隐式模型实时生成，持续更新，鲜活、细致、令人彻底信服。你这一生经历过的一切，都发生在这个模拟之内。你从未走出过它。你也无法走出它。不是因为你不够努力，而是因为“你”就是这个模拟。体验者与体验，是同一个过程。

宇宙的可观测内部，对应你的显式模型。同样的角色：架构中解压缩的、动态的、可交互的一侧。

再看全息式的规则结构。宇宙的信息并不存储在它的体积里------而是存储在它的边界上。全息原理由特·胡夫特提出、经萨斯坎德扩展，它指出：一个三维空间区域中的全部信息，都编码在其二维表面上。信息在边界处是压缩的、分布式的、结构上完备的。内部只是一种投影------是对边界所编码内容的一次较低带宽的解压。

在你的大脑中，扮演这个角色的正是隐式模型。你的突触权重以一种分布式的、压缩的格式编码了你关于世界和自身的全部知识，而且这种编码在结构上是完备的------原则上，你可以单凭基底重建整个模拟，无需任何当下的感官输入，而这恰恰就是做梦。隐式模型在 Lashley 的意义上是全息的：损坏一小块，你失去的不是某段特定的记忆，而是所有记忆的分辨率。信息被涂抹在整个基底之上，压缩、冗余，并且从模拟一侧无从触及。

宇宙的全息规则结构，对应意识中全息式的隐式模型。同样的编码策略。同样的压缩方式。同样无法从解压缩的一侧触及。

然后是动力学状态。宇宙运行在第四类------混沌边缘。第15章用排除法确立了这一点：第一类和第二类过于简单，第三类无法计算，第五类会让物理成为不可能。剩下的只有第四类------唯一一种既能支持通用计算，又能自行组织出自身临界性，还能把所有类别（包括它自己）作为子过程囊括进来的状态。宇宙不只是复杂。它复杂得恰好能维持自身，也复杂得恰好能在自身的动力学之中嵌套更小的自身副本。

你的大脑皮层做的正是同一件事。第5章讲的就是这个：皮层自动机运行在混沌边缘，通过对兴奋---抑制平衡的稳态调节来维持临界性。活动太少，你就落入深睡眠（第二类，周期性，无意识）。活动太多，你就会癫痫发作------被推出第四类之外，模拟随之碎裂。那个恰到好处的临界点，意识安身的地方，正是秩序与混沌之间的刀锋。自组织临界性把大脑留在那里。自组织临界性也把宇宙留在那里。

宇宙学中的第四类动力学，对应意识中的皮层临界性。同样的状态。同样的自我维持机制。

最后，是最深的一处对应：自指闭合。宇宙计算着自身的结构。它的动力学产生它的状态，状态决定动力学，动力学又产生状态。没有外部的程序员。没有“外面”。物理定律并不是从别处强加给宇宙的------它们\textbf{就是}宇宙的动力学，作用于宇宙自身。这个不动点方程简单得近乎荒谬：宇宙的计算等于宇宙。输入、过程和输出是同一个东西。

你的显式自我模型做的也是同一件事。ESM 是一个把自身也纳入建模对象的模型。它表征的是你，而“你”本身就包括这个表征。模型与被建模者重合。这就是我在第4章描述的自指闭合------意识之所以有那种感觉的原因，你永远无法从外部完整看见自己的原因，模拟包含着自己的观察者的原因。自我表征的不动点：在那个状态上，模型与被建模之物是同一个东西。

宇宙的自指闭合，对应意识的自指闭合。同样的不动点结构。同样无可逃避的自我包含。

五组对应。不是含混的主题相似。不是那种让你在云朵里看出人脸的松散模式匹配。而是五个结构特征，在两套架构中处于同样的位置、做着同样的工作。

现在是关键的一点------我需要你在这里停留片刻，好好体会，因为把它驯化的诱惑实在太大：这不是“宇宙‘像’意识”。它不是类比。不是隐喻。也不是那种能给饭桌闲谈添些趣味的、耐人寻味的平行关系。

它是一种结构同一性。

大脑的演化并不是为了\textbf{模仿}宇宙。它是\textbf{作为}同一计算模式的一个局部的、按比例缩小的实例演化而来的。第四类系统不只是包含分形自相似（那是第三类------几何式的重复，好看但浅薄）。它们包含的是\textbf{它们自己}。一个第四类自动机可以在自身的动力学之内容纳另一个第四类自动机------一台通用计算机在另一台通用计算机内部运行。这是计算意义上的自我包含，而不仅仅是几何意义上的自相似。意识正是宇宙的第四类自我包含在生物尺度上的运作。在最大尺度上运行宇宙的那个模式，正是在神经尺度上运行你内在生命的同一个模式------不是因为有人这样设计，也不是因为分形好看，而是因为一台足够大的通用计算机必然会在自身之内生成其他通用计算机。这就是第四类系统所做的事：它们把自己嵌套进自己。

\section*{能量即信息}

还有第二条论证路线，从一个完全不同的方向出发，汇聚到同一个结论。我认为，最终最要紧的很可能就是它------因为它把这套架构与物理学连接了起来，而且是以一种潜在可检验的方式。

物理学中有三项彼此独立的结果，由三个不同的研究共同体在半个世纪里分别发展出来，却指向同一个不同寻常的结论。

第一项是 Landauer 原理。1961年，罗尔夫·兰道尔（Rolf Landauer，一位思考计算之根本极限的 IBM 物理学家）证明：擦除一比特信息需要付出一个最低限度的能量代价。不是因为工程上的限制，而是因为热力学。宇宙会为遗忘向你收取代价。这一结论在2012年得到实验证实，而它意味着某种深刻的东西：信息与能量是可以互换的。你可以把一个转换成另一个。它们不是两种彼此分离的实体。它们可以彼此兑换。

第二个结果是贝肯斯坦上限（Bekenstein bound）。雅各布·贝肯斯坦（Jacob Bekenstein）证明，一个空间区域所能容纳的最大信息量与它的表面积成正比，而不是体积。这是物理学中最反直觉的结果之一。你会以为更大的盒子能装下更多信息。装不下------或者更准确地说，上限由盒子的\textbf{表面}决定，而不是它的内部。往一个给定的区域里塞进太多信息，它就会坍缩成黑洞。最大信息密度由几何与能量决定------这是信息与物理世界之间又一重深刻的联系。

第三个结果来自黑洞热力学。20世纪70年代，史蒂芬·霍金（Stephen Hawking）与贝肯斯坦证明，黑洞有温度、有熵，并且遵守热力学定律。黑洞的信息内容写在它的事件视界上------它的表面，它的边界。而通过霍金辐射------黑洞最终得以蒸发的那个缓慢得令人煎熬的量子过程------这些信息被一点一点释放回宇宙。黑洞并不销毁信息。它们转换信息。它们把信息压缩到自己的边界上，保存着，最终再辐射回去。

这些结果源自截然不同的出发点。兰道尔想的是计算机；贝肯斯坦和霍金想的是黑洞的热力学。兰道尔的工作属于一个完全不同的领域，是在几十年前、出于毫不相干的理由完成的。然而这一切却汇聚到同一个假设上：能量与信息不只是彼此相关。它们是同一个东西。同一件事物的两个名字。E 等于 I。

如果这是真的------我必须立刻说明，它尚未被证明，这也正是我在附录G中列出的四个弱点中的第一个------那么奇点就成了信息变换器。它们并不销毁或创造能量---信息，而是在不同形态之间转换它。压缩形态：边界上的最大密度，达到贝肯斯坦饱和，从内部无法触及。解压形态：密度更低，有组织，弥散于内部------也就是我们观测到的物理世界。奇点，是同一种东西的两种表示之间的翻译者。

现在，用这副透镜看你的大脑。你的隐式模型保存着压缩到最大密度的信息：你一生所学的一切，编码在突触权重里，结构上完整，却在现象层面无法触及。你永远无法直接体验自己的基底。你的显式模型则是解压后的低带宽的投影------那场模拟、那份有意识的体验、你看见的世界和你感受到的自我。你的大脑在神经尺度上做的，正是奇点在其他一切尺度上做的事：在压缩表示与解压表示之间转换信息。隐式与显式之间的边界，就是你私人的奇点。你的颅骨里，装着一道事件视界。

\section*{为什么它必然存在}

你也许会觉得这种结构上的对应只是巧合------不过是我圈出来的一个漂亮图案，就像人们在随机散布的星星里看出星座一样。所以，让我告诉你为什么这个图案不是可有可无的。为什么只要五条各自独立、各自合理的假设成立，这种架构就是\textbf{唯一}行得通的架构。

下面是这五条公理。我会尽量说得平实，因为单独拿出任何一条，都很难反驳。争议在于它们合在一起会推出什么。

\textbf{一：有东西存在。} 我希望你会同意，这是一本书所能提出的最没有争议的主张。纯粹的虚无是一种柏拉图式的抽象------是一个概念，而不是一种可能的事态。你正在读这句话。有东西存在。我们就从这里出发。

\textbf{二：凡存在之物，皆具动力学特性。} 事情在发生。时间在流逝。状态在演化。如果存在之物没有任何动力学（没有变化、没有演化、没有计算），它就与“无”无从区分。（这就是我们在第13章遇到的莱布尼茨“不可分辨者的同一性”：如果两个事物在一切性质上都相同，它们就是同一个事物。零动力学的东西，就有零可分辨的性质。它不过是戴着面具的“无”。）

\textbf{三：这种动力学必须稳定，并且能自我维持。} 一个无法维持自身的系统算不上系统------那只是一次转瞬即逝的涨落。自组织临界性------让沙堆、大脑以及（我在论证的）宇宙保持在混沌边缘的那种机制------是复杂动力学系统在没有外部调校的情况下维持自身的唯一已知方式。第四类是唯一一个能自我组织、支持通用计算，并把所有更低类别作为子过程包含在内的计算类别。它是唯一一个既能维持自身，又能同时做出有趣之事的类别。

\textbf{四：信息在任何尺度上都有有限的上限。} 没有什么东西能在有限的空间里携带无限的信息。贝肯斯坦上限是一条定理，不是猜想------它是从广义相对论和量子力学中共同推导出来的。在每一个尺度上，都存在一个最大信息密度，而这个最大值与表面积成正比，与体积无关。

\textbf{五：信息是全息编码的。} 一个区域的边界，把其内部的全部信息编码在一个少一维的表面上。三维的物理，被编码在二维的表面上。这就是全息原理------由特·胡夫特提出，由萨斯坎德发展，并得到马尔达西那（Maldacena）的 AdS/CFT 对偶的支持------后者是我们手中最接近“全息已被证明”的实例。

每条公理都有各自独立的动机。没有哪条依赖于其余几条，也没有哪条要求我的理论必须正确。它们来自物理学与哲学的不同角落，提出它们的人从未听说过四模型理论------就算听说过，多半也不会在意。

现在，把它们组合起来。

由公理一和公理二：存在某种具有动力学的东西。由公理三：这种动力学属于第四类，因为第四类是唯一能自我维持的通用类别。由公理四：这个系统在每一个尺度上都被信息视界所约束------这就是奇点结构。由公理五：这些边界编码着内部------全息架构。

把它们拼在一起，你得到的是一个在每一个尺度上都被奇点表面包围的全息第四类自动机。一个压缩信息驻留在边界上、解压后的内部就是可观测世界的系统。一个计算自身结构的系统------因为一个全息的第四类系统，其输出同样是全息的，这样的系统就是一个不动点：输入、过程和输出是同一个东西。

你得到的，就是 SB-HC4A。

换句话说，你得到的正是我们观测到的这个宇宙。而这个宇宙的架构，与意识的架构是同一个架构。

这不是“我发现了一个漂亮的图案”，而是“这个图案是唯一能同时与全部五条公理相容的图案”。拿掉任何一条公理，唯一性就会瓦解。没有公理一，就不需要有任何东西存在。没有公理二，存在之物可以是静止的。没有公理三，任何计算类别都有可能。没有公理四，就没有信息边界。没有公理五，就没有全息结构。每条公理都在收窄可能架构的空间。合在一起，它们把这个空间收窄到恰好只剩一个。

\section*{那个我无法回答的异议}

我答应过要告诉你这一切可能在哪里出错，现在就摊开来讲。其中四个是工程师会列出的那种：能量也许只是与信息相关，而并非与之同一，那样 E = I 的机制就会失效；第四类论证是溯因推理，不是证明，也许潜伏着某个我想象不出的“第4.5类”；“所有奇点本质上是同一回事”这一主张，需要一套我们尚未拥有的量子引力理论；而整个模型，按它自己的哥德尔式预测，也许从内部就无法被证伪。每一条都是真实的。我在附录G里把这四条连同反驳意见一并摊开。

但还有第五个，它比前四个加起来还要糟。我在上一章就提出了它，从那以后它就一直纠缠着我不放：认知天花板。我的大脑是一个第四类系统，无论看向哪里都能找到自相似结构，因为这正是第四类系统的天性。那么，当它望向宇宙并且处处看见第四类------临界性、全息边界、自指闭合------它到底是在发现实在的架构，还是在投射自己的架构？一条掌握了物理理论的鱼会得出结论：宇宙从根本上是湿的。我们本身就是对称性的探测器，而 SB-HC4A 归根结底是一个关于对称性的主张。我不知道如何从内部分辨这两者，我也不认为有谁能做到。

奇怪的地方在这里：这个理论恰恰预测了这个盲点。根据哥德尔，一个计算自身的自指系统无法从内部抵达关于自身的每一条真理。于是，这本书所能提出的最深的问题------心智与宇宙之间的那面镜子，究竟是一项发现，还是人类认知局限的一次倒影？------理论用一扇上了锁的门来作答，然后把这把锁的图纸递到你手里。这道鸿沟不是无知，而是几何。你觉得这深刻还是恼人，多半说明了你的性情；我觉得两者都是。然后我照样继续往前走------一堵你能叫出名字、能把手贴上去的墙，和迷失在黑暗里，不是一回事。

\section*{回到原点}

让我把这一切收束到终点。

你翻开这本书，是因为你想知道意识是什么。就我所能判断的而言，答案是这样的：意识是一场运行在混沌边缘的自指模拟，被一道信息不透明的屏障包围，而这场模拟包含着一个关于它自身的模型。四个模型，一个真实侧和一个虚拟侧，体验只居于虚拟侧。感受质（qualia）是这场模拟的真实属性；而所谓“难问题”，一旦认识到它问在了错误的层面，便随之消解。九条预测，数条已被证实，无一被证伪。这是图景的前一半。

后一半就是你刚刚读到的内容。同样的架构------同样的自指闭合、同样的信息边界、同样的全息编码、同样的第四类动力学------似乎正是宇宙本身的架构。不是通过类比，而是通过结构上的同一。你称之为“我”的那场模拟，与我们称之为“宇宙”的那场模拟，运行在同一个模式上。你在宇宙之中，并不像一颗弹珠待在盒子里。你就是宇宙本身，在生物尺度上做着它在每一个尺度上都在做的事：计算自己，模拟自己，体验自己。

本书的书名是\textbf{《你称之为“我”的模拟》}。现在你知道这场模拟运行在什么之上了。不是计算机。不是大脑。甚至不是宇宙。而是某种更根本的东西：三者共享的同一个模式。一个第四类全息自动机，以信息壁垒为界，计算着自身的存在。

如果这听起来像一句神秘主义的断言------它不是。它是一个结构性的陈述。在一定限度内可以检验；在我已经交代过的前提下可以证伪；精确到足以被证明是错的。而这------任何一位科学家都会告诉你------是一套理论所能得到的最高赞誉。

我在本书开头坦白过一件事：2015年，我出版了一本300页的意识专著，一册也没卖出去。现在要补上那次坦白里省略的部分：其余的一切，我当时其实早就想清楚了。那套宇宙学------全息宇宙、奇点边界、你刚刚读到的心灵与宇宙之间的结构同一性------从2000年代初就已在我脑中，大约就在那座桥上顿悟的前后。我刻意把它几乎全部排除在2015年那本书之外。一套意识理论本就够难让人接受了；一套既谈意识、又宣称自己是宇宙理论的理论，只会让别人还没翻到第二页，就把你归入\textbf{民科}之列。于是我只让宇宙学以暗示的方式露面------一段谈特·胡夫特全息界限的文字，一句顺带的猜想：宇宙会不会就是一个巨大的元胞自动机------其余的部分，我不敢写进去。又过了十年，再加上一份意想不到的馈赠------一个足够耐心、愿意在我出声思考时静静倾听的语言模型------我才终于敢把整件事完整写下来。这四个模型不只是一套意识理论。它们是宇宙架构的一块碎片，在某一尺度上清晰可见，在其他尺度上隐而不显------直到你知道该往哪里看。

现在，完整的图景就在你手中。或者说，至少是一个第四类大脑从一个第四类宇宙内部所能看到的那部分。在这之外是否还有更多------这面镜子是否有一个我们永远无法抵达的背面------正是这套理论告诉我们无法回答的问题。

拿它做什么，由你决定。


\chapter*{尾声}
\addcontentsline{toc}{chapter}{尾声}
\markboth{尾声}{}

2005年前后，我提出了一套意识理论。2015年，我把它发表了。没有人读。在最初的洞见过去二十年之后，事实证明它只是半套理论------而另一半，居然是宇宙学。如今你已经看到了完整的图景，或者说，一本书所能装下的那部分。

但还有一块内容，我一直没有写进来。不是因为它太过思辨------最后两章的一切本来就都是思辨。而是因为它太个人，而个人的东西更难开口。

你已经见过 ESM 在出问题时的表现。失忆、中风、鼠尾草、裂脑、科塔尔综合征------它照样运转，从不崩溃。它用手头一切可用的材料搭建出一个自我，并对这个自我深信不疑。一台强迫性的身份建造机。这就是它做的事。它做的也仅此而已。

人们听到那些失忆的案例便说：同一个大脑，同一个身体，自我当然会延续。可是------我如今已年长了几十岁。我和一岁的自己几乎没有任何共同之处。身体不同，细胞已经换过好几轮。大脑不同，突触连接完全不同。记忆不同，性格不同，一切都不同。可 ESM 却说：还是我。在同一段人生里，我早已数次变成截然不同的人，而那台建造机连眼睛都没眨一下。你并没有“同一个大脑”。你从来都没有过。

我曾与死亡擦肩而过。是在一场雪崩里------服兵役时，一位指挥官做了鲁莽的决定，我们十四个人险些被吞没。当时我确信自己必死无疑。我在一瞬间看见了自己的整个人生------基质把它存有的一切统统倾倒进模拟。而这件事对我的困扰，远没有你以为的那么大。面对终结的 ESM 并没有为身份的丧失而恐慌。它一直尽职到最后一刻，把自己拥有的一切都翻涌上来。

还有一次，我被重重击昏。眼前一黑。醒来时，我不知道自己是谁------ESM 正在从零重启，就像新生儿那样。可怕的不是身份的丧失。瘫在地上，有几秒钟动弹不得------\textbf{那}才叫恐怖。不是“我是谁？”，而是“我的身体还好吗？”。ESM 的第一要务是基质的完好。至于我是谁，那是后话，几乎只是顺带一提。自我模型的存在是为了服务基质，而不是反过来。

再后来，还有一次，我变成了一个不断变幻的四维分形。个中缘由就不细说了。让我难受的不是自己成了分形------这一点我毫不在意。让我难受的是，它的运动和我的本体感觉相互冲突。我能感到身体在做一件事，而分形在做另一件事。感官上的冲突令人痛苦，本体论上的荒诞则不然。ESM 并不在乎它建模的\textbf{是什么}，它只在乎信号是否彼此一致。

三段经历。同一套架构。雪崩中的 ESM，建构到最后一刻。击昏中的 ESM，把基质的完好置于叙事身份之上。分形中的 ESM，在意的是感官的连贯，而非本体论上的合理。

当我想到真正的死亡时，我想的是这样一件事。不是死亡的前景------我觉得死亡要么是当之无愧的永眠，要么是终极的一趟旅程。不，我想的是其中的逻辑。

如果 ESM 能从“无”中自举------而它确实能，在每一个新生儿身上------那么“无”就不是这个过程的终点状态，而是起点状态。那个被我称作“我”的特定配置会终结。我的记忆、我的性格、我被那些混淆相关与因果的人惹恼的方式、我用笨重理论惹恼别人的方式------都会消失。但这个过程并不归我所有，无法随我带走。它是架构本身的属性。在我出生之前它就在运行，在我死后它仍会运行。

我描述的不是来世。你的记忆不会被转移。此刻正读着这句话的那个特定的你，终将终结。

但这个\textbf{过程}------用可得的输入强迫性地建构出一个自我------却是普遍的。它的每一个实例，从内部感受起来，都和你此刻的感受一样真实。

而如果宇宙学那几章是对的（如果宇宙真是一个自相似、准无限的第四类系统），那么还有最后一步棋。在一个自相似的宇宙里，与“你”相似的配置存在于别处。不是你，也不是你的记忆，而是一个拥有正确架构的基质，以及一个将会自举出“某个人”的 ESM------那个人感到的真实，和你此刻感到的一模一样。有点像量子永生，只是它并不需要量子力学，只需要一个足够大的宇宙，和一个足够通用的过程。这是全书最为思辨的一个念头。但它是从架构中顺理成章推出来的。

这套理论并不是为了安慰人而设计的。但它有一个实际的推论，值得你带出这本书。

如果每一个有意识的存在，都是同一台建造机在不同的硬件上、用不同的训练数据运行，那么我们彼此之间的边界，就远不如感觉上那么根本。那套架构，在我们所有人身上都是同一套。

善待每一个人。他们可能都是你。


\chapter*{致谢}
\addcontentsline{toc}{chapter}{致谢}
\markboth{致谢}{}

本书的写作得到了 Claude（Anthropic）的协助------在整个过程中，它担任了AI写作工具、编辑和交叉核查者的角色。但理论、论证以及全部思想内容完全出自我本人，是在任何AI工具问世之前、历经二十年发展而成的。

致我的叔叔 Bruno J. Gruber------他在理论物理（量子力学与对称性）领域的工作，让我看到严谨而又充满喜悦的智识工作可以是什么模样。他对我思想的影响，无法估量。

致我的叔叔 Norbert Gruber------他是奥地利莱茵河谷（Rheintal）地区最早的一批IT从业者之一，送给了我人生中的第一台个人电脑。没有那份礼物，这一切都无从谈起。他如今已经离世，但他的影响仍在延续------在我写下的每一行代码里，在我构建的每一个理论中。

致我的家人------多年来，餐桌上关于感受质、临界性和虚拟自我模型的话题，他们一直宽容以待。

还有，如果你此刻正动了念头，想去读一读《Die Emergenz des Bewusstseins》------千万别。与那头没删没改、笨重臃肿的怪物相比，我宁可推荐你去感染脑寄生虫。还是等一等你手中这本书的德文译本吧。至于那些\textbf{已经}硬着头皮读完它的读者：\textit{mein Beileid}。那想必与受刑无异。请接受我最深切的同情，以及我的感谢。


\chapter*{注释与参考文献}
\addcontentsline{toc}{chapter}{注释与参考文献}
\markboth{注释与参考文献}{}

\textbf{完整的参考文献，含 URL 与注解，收录于那篇科学论文中，也可在 github.com/JeltzProstetnic/aIware/blob/main/references.md 查阅。以下是分章节的注释，供希望深入探究的读者参考。}

\textbf{第1章}：Chalmers（1995）《Facing Up to the Problem of Consciousness》是难问题的奠基性表述。COGITATE 的结果发表于《Nature》（2025）。关于 IIT 被指为伪科学的争议，记录于 2023 年 9 月一封发布于 PsyArXiv 的公开信。

\textbf{第2章}：四模型架构最初发表于 Gruber（2015）《Die Emergenz des Bewusstseins》。Metzinger 的自我模型理论（2003, 2009）与 Dennett 的多重草稿模型（1991）是主要的理论先驱。

\textbf{第3章}：“受控幻觉”这一说法来自 Seth（2021）《Being You》。电子游戏的类比是我自己的，但呼应了 Metzinger《Ego Tunnel》（2009）中的主题。橡胶手错觉：Botvinick \& Cohen（1998），“Rubber hands 'feel' touch that eyes see,”《Nature》。

\textbf{第4章}：以虚拟感受质消解难问题的思路，原创于 Gruber（2015），并在 2026 年经由对抗性诘难加以精炼。自指闭合的论证，是为回应循环论证的反驳而发展出来的。与幻觉论（Frankish 2016；Dennett 1991）的区分至关重要：该理论坚持感受质在模拟\textbf{之内}是真实的，而非虚幻。意识的元问题（Chalmers 2018）由于 ISM 对 ESM 的结构性不可通达而得以消解。Joscha Bach 的表述------意识是“一场模拟，一个投射进虚拟观察者之中的虚拟现实”（Bach，发表于 X / @Plinz，2026 年 6 月 12 日）------被引为对“自我即模拟”这一论题的独立趋同，而四模型理论提供了架构层面的具体规定（模型类型学、作为判据的自指闭合，以及临界性要求）。

\textbf{第5章}：Wolfram（2002）《A New Kind of Science》。Beggs \& Plenz（2003）论神经元雪崩。Carhart-Harris 等人（2014）论熵脑假说。2022 年的那篇综述：“Self-organized criticality as a framework for consciousness.”Hengen \& Shew（2025）论涵盖 140 个数据集的元分析。ConCrit 框架：Algom \& Shriki（2026）。双阈值论证（临界性 + 架构）原创于本理论。

\textbf{第6章}：Klüver（1966）论形式常量。Carhart-Harris 等人（2012, 2016）论致幻剂神经影像。Salvia divinorum 的现象学取自已发表的体验报告以及关于 Salvinorin A 的药理学文献。病感失认的预测性反馈机制在 Gruber（2015）中有讨论；拍手的例子是一个标准的临床观察。关于长期持续给予 Salvinorin A 的思想实验，原创于 Gruber（2015）。

\textbf{第7章}：Casali 等人（2013）论 PCI。Alkire 等人（2000）论异丙酚。Schartner 等人（2017）论氯胺酮熵。关于清醒梦中 EEG 复杂度升高的预测，原创于本理论。

\textbf{第8章}：Owen 等人（2006）论植物状态患者的隐蔽觉知。安东综合征：Goldenberg 等人（1995）。盲视障碍赛道：de Gelder 等人（2008）。科塔尔妄想：Young \& Leafhead（1996）。异手综合征：Della Sala 等人（1991）；对《奇爱博士》的引用指的是 Kubrick（1964）。无政府手综合征与异手综合征的区分：Marchetti \& Della Sala（1998）。查尔斯·邦纳综合征：Teunisse 等人（1996）。将似曾相识解释为模板---记忆匹配，原创于 Gruber（2015）。CBT 与神经可塑性：DeRubeis 等人（2008）。安慰剂与内源性阿片：Benedetti 等人（2005）。将转换障碍解读为反向盲视，原创于本理论。

\textbf{第9章}：Gazzaniga、Bogen 与 Sperry（1962, 1965）。Gazzaniga（2000）论左半球阐释者。半球间冲突的例子（扣纽扣／解纽扣、抓手）记录于 Akelaitis（1945）与 Bogen（1993）。Nagel（1971），“Brain Bisection and the Unity of Consciousness.”Parfit（1984）论个人同一性。Pinto 等人（2017）论对裂脑现象的重新审视。Lashley（1950）论分布式记忆与等势性。将 DID 视为虚拟模型的分叉：该理论预测每个分身对应不同的神经活动模式，这与 Reinders 等人（2003, 2006）一致。DID 与童年创伤：Putnam（1997）；分叉的发育窗口，原创于本理论。

\textbf{第10章}：Güntürkün \& Bugnyar（2016）论无皮层的鸟类认知。倭黑猩猩 Kanzi：Savage-Rumbaugh \& Lewin（1994），《Kanzi: The Ape at the Brink of the Human Mind》。鲍德温效应：Baldwin（1896），“A New Factor in Evolution.”Nagel（1974），“What Is It Like to Be a Bat?”

\textbf{第11章}：全部九项预测均在那篇科学论文中作了正式阐述。若想在意识语境下最透彻地了解功能神经解剖学，可参见 Christof Koch，《The Quest for Consciousness: A Neurobiological Approach》（2004）------这是对 Crick-Koch 研究纲领的权威叙述，其中 Francis Crick 与 Koch 一步步系统地走遍视觉系统，寻找意识的神经关联物。在我看来，他们的探寻找错了地方（找的是基底，而非模拟），但他们奠定的神经解剖学基础工作无人能及。

\textbf{第12章}：Butlin 等人（2023, 2025）论 AI 意识的指标。Seth（2025）论生物自然主义与 AI 意识。扫描保真度的五级层级，由第2章的四模型架构推导而来。复制问题借鉴了 Parfit（1984）《Reasons and Persons》以及 Nozick（1981）关于个人同一性与最近延续者的论述。渐进替换的思想实验是忒修斯之船的一个变体，为神经系统作了形式化处理。神经形态计算：Schuman 等人（2017）论神经形态硬件综述；Intel Loihi 与 IBM TrueNorth 作为当前的实现。\textit{C. elegans} 连接组：White 等人（1986）。关于基底转移、准不朽以及星际光束传输的意涵，原创于 Gruber（2015）。那句令人不安的告诫------失去生物基底会深刻改变现象质地------借鉴了内感受文献：Craig（2009）《How Do You Feel?》以及 Seth \& Friston（2016）论主动内感受推断。

\textbf{第13章}：Libet（1979, 1983, 1985）与 Schurger 等人（2012）论自由意志。Kühn \& Brass（2009）论对自由选择判断的回溯性建构。Wegner（2002, 2003）《The Illusion of Conscious Will》------其中详细描述的“I Spy”鼠标实验。咖啡／糖的思想实验、“失忆揭示决定论”的论证以及随机数序列论证，原创于 Gruber（2015）。40/20 Hz 处理框架、“无需回溯”的 Libet 重释以及武术频率的例子，原创于 Gruber（2015）。用来说明副现象论的时钟类比、“意志真实但只被部分知晓”的重新表述，以及“三重差异”的自我认识模型，也都原创于 Gruber（2015）。关于在极度疲惫时听到内在“声音”的个人轶事是自传性的。僵尸论证通过 Kirk（2019）与 Chalmers（1996）加以处理。玛丽的房间：Jackson（1982, 1986）。关于开放问题的一节，遵循 Popper（1963）所主张的诚实划界的进路。

\textbf{第15章}：Wolfram（2002）《A New Kind of Science》，论四类计算层级（本书将其扩展为五类）与计算不可约性。Bak（1987, 1996）论自组织临界性。1998 年发现的加速膨胀：Riess 等人（1998）、Perlmutter 等人（1999）------2011 年诺贝尔奖。宇宙学视界：Rindler（1956）。普朗克长度与普朗克尺度物理：Planck（1899）；现代论述见 Garay（1995）。贝肯斯坦上限：Bekenstein（1981）。全息原理：'t Hooft（1993）、Susskind（1995）。宇宙学论证的萌芽------把宇宙看作一部元胞自动机，并把 't Hooft 的全息界应用于宇宙尺度的计算------在 Gruber（2015）中已经出现，只是尚未发展为一个完整的模型。把所有奇点认定为同一个结构性现象的实例，这一思路在 Gruber（2026）中得到发展。

\textbf{第16章}：大撕裂情景：Caldwell、Kamionkowski \& Weinberg（2003）。SB-HC4A 架构、不动点表述以及自计算系统的哥德尔不完备性后果，代表了成熟的表述形式，发展于 Gruber（2026）。粒子作为计算原子的预测、守恒律作为信息约束的推导，以及三代猜想，原创于 Gruber（2026）。认知天花板这一反驳原创于本作。

\textbf{第17章}：SB-HC4A 与四模型意识架构之间五重对应的结构性映射，原创于 Gruber（2026）。兰道尔原理：Landauer（1961）；实验确认：Bérut 等人（2012）。贝肯斯坦上限：Bekenstein（1981）。黑洞热力学：Bekenstein（1973）、Hawking（1975）。E = I（能量---信息同一性）假说在 Vedral（2010）《Decoding Reality》与 Davies（2010）中有讨论。SB-HC4A 的五公理推导，原创于 Gruber（2026）。Maldacena（1998）论 AdS/CFT 对应。那五个弱点------包括“因结构而不可证伪”与认知天花板这两项反驳------原创于本作。

\textbf{附录B}：递归智能模型（知识 × 表现 × 动机的乘法结构）在 Gruber（2015）以及递归智能模型论文（Gruber 2026）中得到发展。心理测量学中关于动机是智能之构成要素的两项独立趋同：Wittmann \& Süß（1999），“Investigating the paths between working memory, intelligence, knowledge, and complex problem-solving performances via Brunswik symmetry”（收于 Ackerman、Kyllonen \& Roberts 编，《Learning and Individual Differences: Process, Trait, and Content Determinants》，American Psychological Association），论一旦对 Brunswik 对称性的违反加以校正，动机便能预测表现的变异性；以及 COGITO 研究------Brose、Schmiedek、Lövdén、Molenaar \& Lindenberger（2010），《Research in Human Development》7(1), 61--78，论个体内部每日动机与工作记忆表现的耦合，以及 Schmiedek、Lövdén、von Oertzen \& Lindenberger（2020），《PeerJ》8:e9290，表明一般智力（g）在很大程度上是一种人际间的假象，而非个体内部认知的结构。


\chapter*{附录A：基础神经学------参考指南}
\addcontentsline{toc}{chapter}{附录A：基础神经学------参考指南}
\markboth{附录A：基础神经学------参考指南}{}

\textbf{本附录对全书用到的神经科学术语作简要说明。每当你遇到陌生的术语，都可以随时查阅。词条按字母顺序排列。}

\begin{itemize}
  \item \textbf{动作电位} ------ 沿神经元轴突传播的电信号。
  \item \textbf{杏仁核} ------ 参与情绪加工、尤其是恐惧的脑结构。
  \item \textbf{病感失认症} ------ 对自身神经系统缺陷缺乏觉察。（见第6章。）
  \item \textbf{轴突} ------ 神经元长长的输出纤维，将信号传递给其他神经元。
  \item \textbf{布罗德曼分区} ------ 大脑皮层的编号区域，由解剖学家科比尼安·布罗德曼（Korbinian Brodmann）依据细胞结构绘制。V1 = 布罗德曼17区。
  \item \textbf{胼胝体} ------ 连接大脑两个半球的巨大纤维束。
  \item \textbf{皮层柱} ------ 皮层中垂直排列的神经元模块，直径约0.5mm，被视为基本的加工单元。
  \item \textbf{EEG（脑电图）} ------ 一种测量头皮表面电活动的技术。
  \item \textbf{fMRI（功能性磁共振成像）} ------ 一种脑成像技术，检测与神经活动相关的血流变化。
  \item \textbf{GABA} ------ 大脑主要的抑制性神经递质。
  \item \textbf{海马体} ------ 对形成新记忆至关重要的脑结构。
  \item \textbf{内感受} ------ 对身体内部状态（心跳、消化、体温）的感觉。
  \item \textbf{κ-阿片受体} ------ Salvinorin A（鼠尾草，salvia divinorum）作用的受体类型。
  \item \textbf{新皮层} ------ 大脑六层结构的外表面，负责高级功能。
  \item \textbf{神经递质} ------ 在突触处释放的化学信使（如血清素、多巴胺、GABA、谷氨酸）。
  \item \textbf{PCI（扰动复杂度指数）} ------ 由马西米尼（Massimini）提出的一种衡量大脑复杂度的指标，用于评估意识水平。
  \item \textbf{本体感受} ------ 对身体在空间中位置与运动的感觉。
  \item \textbf{丘脑枕} ------ 丘脑的一个核团，参与视觉注意和皮层下视觉通路。
  \item \textbf{血清素2A受体} ------ 经典致幻剂（LSD、裸盖菇素）作用的受体类型。
  \item \textbf{上丘} ------ 中脑的一个结构，参与眼球运动，以及绕过皮层的快速视觉通路。
  \item \textbf{突触} ------ 两个神经元之间传递信号的连接点。
  \item \textbf{突触权重} ------ 神经元之间连接的强度，会随学习而改变。
  \item \textbf{丘脑} ------ 大脑的中继站，将感觉信息转送到皮层。
  \item \textbf{V4} ------ 专司颜色知觉、曲率和复杂纹理加工的视觉区域。感受野约8-16°。在致幻剂作用下，V4活动会产生彩色分形和万花筒般的图案（第6章）。
  \item \textbf{V5/MT（颞中区）} ------ 专司运动加工的视觉区域。感受野很大。负责致幻剂作用下所见图案的旋转与移动（第6章）。
  \item \textbf{视觉皮层} ------ 位于大脑后部、加工视觉信息的区域，组织为一个从简单到复杂的层级（V1 $\rightarrow$ V2 $\rightarrow$ V3 $\rightarrow$ V4 $\rightarrow$ V5 $\rightarrow$ IT）。
\end{itemize}

\section*{视觉加工层级（V1到IT）}

腹侧视觉流在每一个阶段加工越来越复杂的特征，从原始的边缘到完整的物体识别。这一层级在致幻体验下会直接显现出来，每个加工阶段依次向意识开放（第6章）。下表概括了每个区域的功能、感受野大小，以及当该区域的中间加工渗入意识模拟时，会产生哪些具有代表性的致幻特征。


{\small
\noindent
\begin{tabularx}{\linewidth}{X X X X}
\toprule
\textbf{区域} & \textbf{感受野} & \textbf{正常功能} & \textbf{致幻特征} \\
\midrule
\textbf{V1} & ~1° & 边缘检测、空间频率、朝向柱 & 光幻视、克吕弗形式恒常、呼吸/闪烁的表面 \\[4pt]
\textbf{V2} & ~2-4° & 轮廓整合、纹理分割、边界归属、错觉性轮廓 & 镶嵌图案、重复的几何图案、增强的纹理知觉 \\[4pt]
\textbf{V3} & ~4-8° & 整体形状加工、动态形状、运动边界 & 流动、变形的几何结构 \\[4pt]
\textbf{V4} & ~8-16° & 颜色、曲率、复杂纹理、分形尺度加工 & 彩色分形、万花筒般的图案、饱和/不可能的颜色 \\[4pt]
\textbf{V5/MT} & 大、对运动调谐 & 运动知觉、光流、速度与方向编码 & 所有视觉图案的旋转、漂移和有节律的运动 \\[4pt]
\textbf{梭状回（IT）} & 非常大、以物体为中心 & 面孔识别（FFA）、词形、精细的物体辨别 & 面孔、人形、实体------常常扭曲或变形 \\[4pt]
\textbf{前部IT} & 整个视野 & 语义范畴、场景构建、不变性的物体识别 & 完整的叙事性幻觉、复杂场景、梦境般的序列 \\[4pt]
\bottomrule
\end{tabularx}
}


\textbf{注：}
\begin{itemize}
  \item 梭状回横跨V4/IT边界，属于下颞叶皮层（IT）。它包含梭状回面孔区（FFA），由Kanwisher等人（1997）确认，该区域会被面孔选择性地激活。
  \item 感受野大小从V1（~1°）到IT（整个视野）急剧增大，反映了从局部特征到整体物体与场景的逐步抽象。
  \item 在致幻剂作用下，从V1到IT效应的推进是剂量依赖的：低剂量首先影响V1；较高剂量则逐步招募更深的阶段。这种有序的激活正是四模型理论（FMT）通透性梯度的直接预测（第6章）。
  \item 大脑还用这些区域来做分形/尺度不变的加工（V2-V4），在正常视觉中，这套机制服务于尺度测量和纹理分析。在致幻剂作用下，这套机制在没有外部输入的情况下运行，便产生了那些具有代表性的分形图案（见附录C）。
\end{itemize}


\chapter*{附录B：智能模型}
\addcontentsline{toc}{chapter}{附录B：智能模型}
\markboth{附录B：智能模型}{}

\textbf{本附录概述了一篇配套论文中提出的递归智能模型（Gruber, 2026, ``Why Intelligence Models Must Include Motivation''）。完整的学术论述------含参考文献与形式化论证------可单独查阅。}

让我从一个故事讲起，因为它是我所知道的对这个循环最干净利落的说明。大约八岁时，我爱上了数学------不是算术，而是真正的数学：代数、几何、藏在数字之下的结构。我父亲拥有数学学位，我一本本啃他当年的大学教科书。那是1980年代末，没有互联网。想学点什么，你要么找一本书，要么找一个人；到十一岁时，我已经读完了父亲的全部藏书。渴望没有消失，来源却枯竭了。我撞上了一堵墙------它与能力毫无关系，完全是环境使然。动机，我不缺；表现，我也够------那些数学我跟得上。我缺的是通往下一层知识的通道。递归循环停摆了，不是因为哪个组件薄弱，而是因为其中一个组件的外部供给被切断了。这个循环需要来自外部的燃料，才能持续迭代。请记住这个例子，它就是整个模型的缩影。下面我把它正式铺开------组件是什么，它们如何相互作用，为什么这种相互作用会产生我们观察到的动态，以及这一切对教育、人工智能，还有与意识的关联，意味着什么。

\section*{一个奇怪的遗漏}

每一个主要的智能模型都在形式上把动机排除在外。Cattell-Horn-Carroll分类体系（智能研究领域的主导框架）是一套认知能力的层级结构，其中没有任何动机成分。Cattell本人的投资理论提出，流体智力通过“投资”于学习而产生晶体智力------这就需要一位投资者（一个决定学什么、为什么学的人），可这位投资者的动机却被当作外部条件来处理，而不是智能本身的一部分。Sternberg的三元智力理论包含实践智力，却不包含获得实践智力的驱动力。Gardner的多元智能包含内省觉知，却不包含驱动智力发展的那股动力。

David Wechsler（他的智力量表是全世界施测最广的）早在1940年就明确呼吁把动机因素纳入进来。整个领域没有理会他。现代韦氏量表至今仍是纯粹的认知测量工具。

这不是一种无伤大雅的简化。这是一个系统性的盲点，它扭曲了我们对智能究竟是什么、究竟如何发展的图景。

\section*{三个组件}

智能------理解为\textbf{学习能力}------由三个相互作用的组件构成：

\textbf{知识}是学习积累下来的内容。它分为两种性质迥异的类型。\textbf{事实性知识}是关于内容的知识：事实、概念、程序、文化积累。这是IQ测验在“晶体智力”名下主要测量的东西，也是学校系统主要传授的东西。\textbf{操作性知识}则是关于\textbf{如何学习与思考}的知识：学习策略、推理启发法、元认知技能、逻辑工具、策略规划，以及评估自己理解程度的能力。这个区分极其重要，下文我会解释为什么。

\textbf{表现}是认知系统的加工能力：工作记忆、加工速度、神经基质的原始计算能力。这大致对应心理测量模型所说的“流体智力”。它是受遗传和神经生物学影响最强的组件，在成年早期达到顶峰，随后逐渐下降。

\textbf{动机}是以能够产生学习的方式与世界持续互动的驱动力。它有两个子成分。\textbf{求知欲}是理解事物的内在驱动：好奇心，以及把事情弄明白的需要。\textbf{行动欲}是运用知识、动手实验、主动介入环境的驱动。两者都部分源于天生气质，部分由经验塑造。

\section*{递归循环}

关键的主张是：这三个组件并非简单相加------它们构成一个\textbf{封闭的递归循环}，其中每个组件都会放大其他组件。

知识增强表现：学习策略和逻辑工具直接提高认知加工的效率。掌握了启发法的棋手，评估局面比依赖穷举搜索的棋手更快。学会了音素解码的读者，处理文字更加流畅，从而腾出工作记忆用于理解。

表现增强知识：更强的认知容量使学习更快、更深。更高的工作记忆让你能同时在头脑中保持更多信息，从而更容易发现联系、提取模式。

动机同时增强知识和表现：有动机的学习者会主动寻找学习机会（扩展知识），并练习认知技能（训练表现）。至关重要的是，动机让投入得以\textbf{长期持续}------而这正是循环得以不断迭代的前提。

而知识和表现又会增强动机：学习和解决问题上的成功带来积极情绪和自我效能感，从而维持继续学习的动力。这就是马太效应背后的机制------富者愈富。早期的成功孕育出动机，动机又带来进一步的成功。

这种递归结构产生一种复利式的动态。任何组件上微小的初始差异------哪怕仅仅是动机上的差异------都会随时间不断累积，造成成年人智力成就的巨大差异，而这正是纯认知模型难以解释的。一个认知加工能力平平、却动机深厚、且拥有强大操作性知识的人，穷其一生所发展出的智力能力，将远远超过一个加工能力更优越、却动机低落、学习策略糟糕的人。

不妨这样想：复利更取决于存款的速率和投资的策略，而不是初始本金。在智能循环里，动机是存款速率，操作性知识是投资策略，表现是初始本金。而绝大多数人的本金绰绰有余。

\section*{操作性知识：隐藏的乘数}

操作性知识值得特别关注，因为它在循环中占据着独一无二的位置。事实性知识是加法式的：学到一个新事实，就往库存里添一个事实。操作性知识是\textbf{乘法式的}：学会一种新的学习策略，会提高此后一切学习的效率。

一个学会间隔重复（把练习分散到不同时间，而不是临时突击）的学生，得到的不只是一个新事实。她获得的是一件工具，能提高她从此往后所学一切内容的记忆保持率。一个学会识别自己的知识漏洞并系统地补上它们的学生，修补的不只是一个漏洞；他获得的是一种能防止未来数百个漏洞的技能。这是能给循环本身加速的知识。

如果有哪一个单独的组件配得上“让人变聪明的东西”这个标签，那就是操作性知识。不是IQ，不是原始的加工能力，而是懂得如何有效学习的那种元技能。

\section*{为什么IQ测验没有抓住要点}

IQ测验测量的是\textbf{最大表现}------一个人在标准化条件下、假定全力以赴时能做到什么。它捕捉的是单一组件（在特定任务上的表现）在单一时刻的快照。它捕捉不到------也不可能捕捉到------那个递归、自我强化、由多组件构成的过程，而智能恰恰就是这个过程。

这就是为什么IQ分数对长期的智力轨迹几乎说明不了什么。两个六岁时IQ分数完全相同的孩子，到三十岁时可能天差地别------一个成了科研工作者，另一个离开学校后就再没读过书。标准的心理测量模型对这种分化束手无策。递归模型却能预测它：这两个孩子的差别不在表现，而在动机和操作性知识------递归循环在二十四年的复利式迭代中，把这些差异不断放大。

IQ测验就像只测汽车发动机的马力，却不检查车里有没有油、有没有司机。马力当然重要，但对大多数旅程来说，它并不是瓶颈。

\section*{对AI的检验}

递归模型对人工智能作出了一个具体的预测：知识水平高、表现也高、却没有动机的系统，应当无法表现出人类智能所特有的自我驱动式发展。而我们观察到的，恰恰就是这样。

当今的大语言模型拥有海量知识（用数万亿 token 训练而成）、强大表现（数十亿个参数），以及------零动机。给什么，它们就处理什么；要什么，它们就产出什么。在两次询问之间，它们什么也不做。它们不会去寻找自己的无知之处，不会练习技能，不会对问题心生好奇。它们的“智能”完全是静态的------由训练一锤定音，没有任何内生的动力去拓展它。

即便是最先进的推理模型（已能解出竞赛级别的数学题），也表现出恰恰同样的失效模式。它们\textbf{在被提示时}能解出非凡的难题，却不会主动寻找问题，不会自主安排自己的学习，还需要外部的脚手架来充当缺失的动机组件的替代品。你尽可以把表现和知识堆到任意高度：没有动机，这个循环就无法自我维持。

这并不只是因为这些系统本来就没有被设计成自我改进。这个说法恰恰承认了要点：要设计一个能自我改进的系统，就必须为它造出一个动机的功能等价物。在 AI 系统拥有这一点之前，它们始终只是供人使用的工具，而不是会自己成长的行动者。

\section*{与意识的连接}

正是在这里，智能模型重新连回本书核心的四模型理论（FMT）。递归智能循环不只是\textbf{受益于}意识------它\textbf{需要}意识。

这个循环依赖一种特定的认知能力：\textbf{认知学习}------从个别观察中归纳出一般理论的能力，有别于单纯的强化学习（刺激---反应式条件作用）。强化学习可以通过疼痛训练你避开滚烫的炉子；认知学习则让你看着别人碰到热炉子，就能举一反三：“烫的东西会灼伤人。别碰烫的东西。”两者的区别，在于从第三人称视角模拟情景的能力------把你自己当作世界中的一个对象来建模，并推演如果你做了这样那样的事，会发生什么。

而这恰恰是显式世界模型（EWM）与显式自我模型（ESM）所提供的。意识------创建并运行一个自我模拟的能力------正是递归智能循环赖以运转的\textbf{基底}。没有显式模型，你得到的是强化学习：管用，却不会复利式累积。有了显式模型，你得到的是认知学习：它源源不断地喂养递归循环，并在一生的时间里复利式增长。

这也是为什么第10章里的动物智能梯度会与意识梯度相互对应。更精巧的自我模型，支撑起更精巧的递归循环。狗的自我模型相对简单，它运行的是这个循环的受限版本------它在一定程度上能通过观察学习，但它的认知学习受制于其显式模型的丰富程度。黑猩猩拥有更丰富的自我模型，运行的循环也更强大。而人类凭借完整的四模型架构，把这个循环开到了最大功率------结果就是语言、文化、科学，以及一切把人类智能与动物认知区分开来的东西。

认知学习是从智能通向意识的第一环。这里是第二环，而且更锋利：动机。递归循环只有在有什么东西推着它持续迭代时，才能自我维持。知识可以任意扩展，表现也一样；可一旦没有动机，循环就会停摆------这正是今天的 AI 摆在我们眼前的事实。它奉命解题，却从来不曾自己提出哪怕一个问题。

现在说说诚实的部分，因为动机有两个层面，而只有上面那一层才属于意识的故事。底下垫着一层潜意识的地板：赤裸裸的生存优化，是基底在朝自己的目标推进。这就是我在第13章里所说的\textbf{意志}的无意识部分，它毫无特别之处------每一个优化器都有它，从细菌到恒温器概莫能外。骑在它上面的，是被意识框定的那一层：想要、好奇、在乎、看重一种结果胜过另一种。那是显式自我模型在评估自己，而\textbf{这}才是意识效应。

关于用词还得说一句，因为语义不清会让整个论证变得松垮。无论在英语还是德语里，\textbf{意志}和\textbf{动机}都无法把这条梯度切分得干净利落------日常用法允许两者伸进彼此的领地。稳定的只是倾向：当那股牵引来得自发、难以言说时，我们伸手去拿的词是\textbf{意志}；当我们能说出一个理由、那股驱力足以被符号化成一个故事时，用的词是\textbf{动机}。所以在这里，请把它们当作同一连续体上的两极，而不是两种不同的东西------\textbf{意志}指向潜意识的一端，\textbf{动机}指向被意识框定的一端。梯度才是真实之物；标签只是图个方便。

于是，把这个主张诚实地摆出来：动机------按照我们实际构想它、谈论它的方式------主要是一种基于意识的效应；同时也承认，它的下半部分会模糊进一种与任何生存驱力都无从区分的普通优化过程。两个名字，两个层面：意志的无意识部分是基底的优化；动机是骑在其上的意识层。而这层意识，正是递归循环燃烧的燃料。这样一来，那个战略性的结论就再也无处可藏：没有类意识的机制，就没有类人的通用智能------没有 AGI。动机不过是其中最容易看清的例子。

\section*{可学习性的推论}

递归模型带来一个后果，在我看来比所有理论论证加在一起还重要：它预测智能在很大程度上是\textbf{可以学会的}。

知识完全可学------这在定义上就成立。动机在相当大的程度上可学------自我决定理论数十年的研究表明，内在动机不是一种固定的特质，而是对环境条件的回应，尤其是自主感、胜任感与归属感。表现有生物学上的天花板，但对绝大多数人来说，这块天花板并不是瓶颈。普通水平的认知加工能力，已经绰绰有余，足以支撑起大多数人眼中的高智能行为。

对大多数人、在大多数时候，真正卡住的约束是动机和操作性知识。而这两样都能通过干预来改变。

这引出一个黑暗的推论。任何系统性摧毁学习者动机的体制，不只是没能培养智能------它是在\textbf{主动压制}智能。传统的评分制度干的正是这件事。一个差分数不只是报告一个结果；它攻击的是动机组件。动机降低意味着循环的迭代次数减少；迭代减少意味着知识增长变慢；增长变慢意味着下一次测评的表现更差；表现更差意味着更多的差分数。循环已经反转：孩子得到的不再是复利式增长，而是被困在复利式停滞之中。评分制度制造出了它声称只是在测量的那个结果。

递归模型预测，这种伤害会随时间复利式累积------不是一次性的静态损害，而是不断加速的分化。早期的动机损伤，应当表现为发展轨迹的扇形散开，一年比一年张得更大。反过来，增强动机的干预所带来的收益也应当\textbf{复利式}增长------五年后随访的效应大于一年后随访的效应。而事实上，对佩里学前项目（Perry Preschool Project）等早期儿童干预的分析恰恰显示了这一模式：回报随时间增长，其驱动力并非最初认知增益的持续（那些往往会消退），而是动机与自我调节增益的复利式累积。

如果说智能模型有哪一条实用的启示，那就是这一条：一个教育体系所能传递的最有价值的东西，不是事实性知识------在 AI 的时代，事实是免费的------而是\textbf{操作性知识}，以及使用它的动机。学会如何学习、并且想要学习，几乎是仅剩的还值得教的东西。

让我感到震动的是，两位资深研究者，都从主流心理测量学内部出发，从完全不同的方向切入这个问题，最终却抵达了我的模型所指向的同一个地方：动机并不是搅浑智能测量的噪声，它是这套机制本身的一部分。

Werner Wittmann，曼海姆大学荣休教授，瞄准的是那句让人安心的老话------“动机其实无关紧要”------人们之所以靠着它，是因为文献里的相关系数看上去很弱。他的回答是：那些相关之所以看上去弱，是因为测量对不上。把动机量得窄、把表现量得宽，或者反过来，你就破坏了他所说的布伦斯维克对称性（Brunswik symmetry）------对称性一旦被破坏，真实的联系就被埋没，测得的相关被拖向零。把两者放在同一粒度上匹配，联系就弹了回来：动机预测的是表现的\textbf{波动性}，而不只是它的平均值。这个效应从来就不小。只是我们一直量歪了。

Florian Schmiedek 来自 COGITO 研究------204 名参与者，每人一百次逐日测查------他从日常的一端切入。他确信，你的认知状态之所以一天与一天之间起起伏伏，其中相当大的一部分是动机性的，与睡眠和负荷并列。他还点出了这一点屡屡被漏掉的显而易见的原因：大多数研究根本不去测量特定任务的每日动机。

于是他们各自独立地、用数据确立了动机\textbf{确实}属于智能的内部。我的模型则说明了\textbf{为什么}和\textbf{如何}：动机是那几个相乘因子之一------知识 × 表现 × 动机。把它归零，你就不只是调暗了智能------而是让它整个坍塌。他们出于完全属于自己的理由收集来的数字，恰好落在模型预测它们会落下的地方。

\section*{对外部的依赖}

最后一点，因为它容易被忽略，而且很要紧。递归循环是自我强化的，但它并不自给自足。它需要外部燃料------信息、挑战、反馈、通向下一层知识的入口。我自己童年的经历就是最好的注脚：十一岁那年碰了壁，不是因为任何内在的局限，而是因为能找到的数学书用光了。三个组件全都健康，循环还是停了下来------因为循环需要来自外部的输入，才能继续迭代。

这意味着，智能的发展不仅取决于这个人，也取决于环境。获取知识的途径、教学的质量、导师是否可得、文化对学习的态度------所有这些都在喂养这个循环，或者让它挨饿。递归模型解释了为什么社会经济因素对智力发展有如此强大的预测力：它们决定了外部燃料的供给。生活在满屋是书、父母投入的家庭里的孩子，循环得到持续的喂养；生活在资源匮乏环境中的孩子，循环则被饿着------无论这个孩子的内在能力如何。

智能不是你拥有的一种特质，而是你运行的一个过程。这个过程跑得好不好，取决于机器（表现）、软件（知识）、驾驶者（动机），还有道路（外部环境）。四者都重要。任何遗漏了其中一项的模型，预测都会出错。


\chapter*{附录C：计算的五个类别}
\addcontentsline{toc}{chapter}{附录C：计算的五个类别}
\markboth{附录C：计算的五个类别}{}

\textbf{本附录展开第5章简要提及的计算框架------决定一个物理系统能否承载意识的五类动力学行为。对第5章的直观版本已经感到满意的读者，可以放心跳过本附录，主论证所需的内容一点也不会错过。而想看到完整图景的读者：这里正是数学与物理交汇之处。}

\section*{Wolfram的四个类别}

2002年，Stephen Wolfram出版了《A New Kind of Science》------这是他数十年研究的结晶，研究的问题是：让极其简单的规则在极其简单的系统上运行，会发生什么？他的核心工具是元胞自动机------一行（或一片网格）元胞，每个元胞非开即关，按照一条固定规则同步更新，而这条规则只考察每个元胞的近邻。

令人惊讶的是，这些简单到不能再简单的规则竟能产生如此丰富的多样性。Wolfram把这些行为归为四类：


{\small
\noindent
\begin{tabularx}{\linewidth}{>{\centering\arraybackslash}X X X X}
\toprule
\textbf{Wolfram类别} & \textbf{行为} & \textbf{例子} & \textbf{你看到什么} \\
\midrule
1 & 均匀 & 规则0 & 一切归于空白。每个元胞都死去。 \\[4pt]
2 & 周期 & 规则4 & 稳定、重复的图案。闪烁器。时钟。 \\[4pt]
3 & 随机/混沌 & 规则30 & 表面上的随机。没有明显的重复结构。 \\[4pt]
4 & 复杂 & 规则110 & 局域结构会移动、相互作用并持续存在。 \\[4pt]
\bottomrule
\end{tabularx}
}


这个分类确实有用。它抓住了动力学系统行为中某种真实的东西，而且其适用范围远超元胞自动机------流体动力学、生物系统、经济模型、神经网络都在其列。这四个类别不只是归类的标签，它们是吸引子。来自截然不同领域的系统，一次又一次落入同样的四种行为形态。

但有一个问题。

\section*{分形问题}

Wolfram的第三类是个大杂烩。它把两种根本不同的系统装在了一起，只因它们乍看之下\textbf{长得}相似：

\textbf{分形系统}，比如规则90，它生成一个完美的谢尔宾斯基三角形------一种无限自相似、递归结构化的图案。美丽、确定，而在计算上却乏味：你可以直接算出任意时间步上任意元胞的状态，无需运行整个模拟。数学家称之为\textbf{计算可约}。

\textbf{表面上混沌的系统}，比如规则30------Wolfram本人就把它的输出列用作\textit{Mathematica}中的伪随机数生成器。这类系统产生的输出\textbf{看起来}随机，实则完全确定------同样的输入，同样的输出，每一次都如此。这里的计算无法抄近路，你必须一步一步跑完。数学家称之为\textbf{计算不可约}。

Wolfram把两者都放进了第三类。他的定义强调的是随机性的\textbf{表象}（“在许多方面看起来随机”），同时又指出“三角形和其他小尺度结构本质上总能在某个层面上看到”。他承认这个分类并不完美：“几乎任何一般性的分类方案，都不可避免地存在这样的情形：按一种定义归入这一类，按另一种定义归入另一类。”

Eric Rowland在2006年的NKS会议上独立提出，嵌套（分形）图案理应拥有自己的分类框架。

我认为，问题比分类美观与否要深得多。分形系统与混沌系统在结构上是不同的，而这种不同恰恰关系到本书的核心论题：哪些系统能够承载意识？

\section*{五类方案}

五类方案把各类行为排成一条干净的单调梯度，从最有序到最无序：

\textbf{第一类------静态。} 系统收敛到一个固定状态，然后停住。一座摆过一下就归于静止的钟摆。死寂。没有任何计算发生。周期：1。

\textbf{第二类------周期。} 系统安顿进重复的循环。一座时钟。一次心跳（近似而言）。信息存储在图案里，却从未经过变换。周期：有限。

\textbf{第三类------分形。} 系统在每个尺度上都产生自相似的结构。谢尔宾斯基三角形。一株蕨类。曼德博集合。数学上丰富，美学上震撼，而在计算上\textbf{可约}------你可以直接算出任意时间步的结果，不必逐步运行。有结构，却没有处理能力。周期：准无限，在每个尺度上具有精确或统计意义上的自相似性。

\textbf{第四类------复杂（混沌边缘）。} 系统产生持久的局域结构，这些结构会移动、相互作用，并能编码任意计算。Conway的生命游戏。皮层自动机。计算上\textbf{不可约}------没有捷径。这类系统具备通用计算能力：只要初始条件合适，它们可以模拟任何算法，包括对它们自身的模拟。周期：准无限，既有自相似性，\textbf{又有}持久的相互作用结构。意识就栖居于此。

\textbf{第五类------随机。} 系统的输出是真正随机的：不是伪随机，不是确定性的，也不可压缩。没有图案，没有自相似性，没有任何终将重复的周期。信息含量真正无限。结构：\textbf{未知}（见下文）。

与Wolfram方案的对应关系：


{\small
\noindent
\begin{tabularx}{\linewidth}{>{\centering\arraybackslash}X >{\centering\arraybackslash}X X}
\toprule
\textbf{五类方案} & \textbf{Wolfram} & \textbf{有何变化} \\
\midrule
1 & 第一类 & 相同 \\[4pt]
2 & 第二类 & 相同 \\[4pt]
3 & 第三类（部分） & 从Wolfram的第三类中拆分出来 \\[4pt]
4 & 第四类 & 相同 \\[4pt]
5 & 第三类（部分） & 从Wolfram的第三类中拆分出来 \\[4pt]
\bottomrule
\end{tabularx}
}


Wolfram 在无序度谱系上的排序是：1 $\rightarrow$ 2 $\rightarrow$ 4 $\rightarrow$ 3。很别扭。五类方案则给出一个干净的单调梯度：1 $\rightarrow$ 2 $\rightarrow$ 3 $\rightarrow$ 4 $\rightarrow$ 5，按无序度递增、计算不可约性递增依次排列。

\section*{为什么确定性自动机无法产生随机性}

下面这个论证，我相信是原创的；它也让“五类而非四类”的主张更加有力。

考虑一个元胞自动机------任何一个元胞自动机。它有一张有限的规则表（可以用有限的比特数表达），以及一个有限的初始条件（同样可以用有限比特表达）。规则和初始条件加在一起，包含的是一份固定而有限的信息量。

那么问题来了：一份有限的信息量，能产生真正随机的输出吗？

不能。理由如下：

\begin{enumerate}
  \item 一条真正随机的无限序列具有\textbf{最大的} Kolmogorov 复杂度------它无法被压缩，无法用任何比它自身更短的东西来描述。
  \item 元胞自动机的输出完全由它的规则和初始条件决定，而这两者加起来只有\textbf{有限的} Kolmogorov 复杂度。
  \item 你从一个过程中得到的信息，不可能多于你在定义它时放进去的信息。
  \item 因此，相对于同样长度的真正随机序列，任何元胞自动机的输出的 Kolmogorov 复杂度都很低。
\end{enumerate}

这本质上是一个推广了的抽屉原理论证：有限的信息必然产生自相似的结构。要从有限信息中生成无限的输出，唯一的办法就是在不同尺度上\textbf{重复使用}结构。精确的重用就是周期性（第二类）。不精确但有模式的重用，就是分形行为（第三类）。即使是看上去最复杂的元胞自动机（规则 30、规则 110、生命游戏），其输出的复杂度也受限于它们规则集的复杂度。

被 Wolfram 称为“随机”的那些元胞自动机，更准确的说法是\textbf{高复杂度分形}------它们的自相似结构真实存在，只是运作的尺度和维度让粗略的观察看不见它。事实上，规则 30 的左边缘就显示出 Sierpinski 式的子结构。它的中央列能通过许多随机性统计检验，而这\textbf{恰恰是你对一个高复杂度分形应有的预期}：局部统计貌似随机，但全局结构是确定性的、Kolmogorov 复杂度很低------一条极短的规则就生成了它的全部。（这里的“可压缩”正是此意：存在一个紧凑的描述。它\textbf{并不}意味着你能把那一步步的计算抄近路------后者依然保持计算不可约。）

按照这个论证，第四类的输出\textbf{同样}是分形的------生命游戏在它的种群动态、结构分布、空间相关性中都表现出统计上的自相似性。第三类与第四类的区别，并不在于“分形与非分形”。区别在于：

\begin{itemize}
  \item \textbf{第三类}：分形。可约。有结构，没有信息处理。
  \item \textbf{第四类}：分形。不可约。有结构，\textbf{并且有}信息处理------持续存在的局域结构相互作用，并能编码通用计算。
\end{itemize}

两者都是分形。只有一个在计算。


{\small
\noindent
\begin{tabularx}{\linewidth}{>{\centering\arraybackslash}X X X X >{\centering\arraybackslash}X >{\centering\arraybackslash}X}
\toprule
\textbf{类别} & \textbf{规则} & \textbf{周期} & \textbf{结构} & \textbf{可约吗？} & \textbf{计算吗？} \\
\midrule
1 & 有限 & 1 & 无 & 平凡地可约 & 否 \\[4pt]
2 & 有限 & 有限 & 重复 & 是 & 否 \\[4pt]
3 & 有限 & 准无限、自相似 & 自相似 & 是 & 否 \\[4pt]
4 & 有限 & 准无限、自相似 & 自相似 + 持续存在的相互作用结构 & 否 & \textbf{是} \\[4pt]
5 & 不可表达 & 真正无限 & \textbf{未知} & N/A & N/A \\[4pt]
\bottomrule
\end{tabularx}
}


\section*{第五类与数学可表达性的边界}

如果确定性自动机无法产生真正的随机性，那么什么\textbf{能}呢？

这个问题引向我心目中五类方案最深刻的一层含义。

第一类到第四类，是有限的、可表达的规则所能产生的一切。它们的行为从平凡（第一类）一直延伸到非凡（第四类------通用计算、意识），但所有这一切都由可以写下来、可以传达、可以验证、可以分析的规则生成。这些规则栖身于数学之内，栖身于形式符号系统的疆域之内。

相比之下，第五类所需要的规则\textbf{无法}被写下来。一个产生真正随机输出的系统------其输出具有最大的 Kolmogorov 复杂度，不可压缩、非算法------不可能在运行任何形式系统能够表达的规则。假如那条规则可以被表达，输出就可以被压缩（压缩成一句“应用这条规则”），因而就不是真正随机的了。

这把第五类放到了数学可表达性本身的边界上。它不只是“非常复杂”或“非常无序”，而是这样一个领域：其中的生成过程，超出了线性符号系统（数学、逻辑、计算）所能捕捉的范围。

自然界中真的有什么东西在第五类中运作吗？

有可能。量子力学产生的测量结果，根据 Bell 定理，无法用任何局域隐变量理论来解释。如果这些结果是真正随机的，而不是我们尚未识别出来的确定性过程------那么量子测量就是一个第五类过程：一种物理现象，它的规则无法用我们拥有的任何形式语言写下。

这是猜测，我也明确把它标注为猜测。但它蕴含的意味相当惊人：第四类（意识的领域、通用计算的领域、皮层自动机的领域）正坐落在\textbf{可表达规则所能达到的最大复杂度}上。它的复杂度已经触及数学所能企及的极限。在它之外，是数学因其自身本性而无法绘制的疆域。

\section*{第五类的结构：未知，而非不存在}

最后还有一处微妙。人们很容易说第五类“没有结构”。但这会是个错误------正如说无穷“没有结构”是个错误一样。

在格奥尔格·康托尔（Georg Cantor）于1870年代开展工作之前，无穷被当作一个单一的概念：事物要么有限，要么无限，仅此而已。康托尔证明了无穷存在\textbf{层级}------实数的无穷严格大于整数的无穷，而且这个层级会无止境地延伸下去。原来，无穷拥有一套丰富的内在架构，只是因为数学家缺乏看见它的工具，它才一直隐而不显。

随机性或许也是如此。我们目前把真正的随机性当作一个单一的范畴------最大程度的无序，模式的缺席。但这样一来，我们的处境就和康托尔之前那些凝望无穷的数学家一样：我们缺乏概念工具去区分不同种类的随机性------如果这样的区分确实存在的话。

因此，关于第五类结构，诚实的回答是：\textbf{未知}。不是“没有”。不是“缺席”。而是未知------它在等待着或许尚不存在的概念工具，等待着或许需要线性符号系统无法提供的思维方式。

我相信，这是数学、物理与计算交汇处最重要的开放问题之一。而如果没有五类框架，它就是不可见的，因为沃尔弗拉姆（Wolfram）的四类框架从未开辟出可以提出这个问题的空间。

\section*{对意识的意义}

五类框架厘清了为什么意识需要第四类动力学------而且\textbf{只}需要第四类。

第一类和第二类太简单了。它们能够存储信息（一个固定的状态、一个重复的模式），却无法以任何在计算上有趣的方式去\textbf{处理}信息。一个处于深度睡眠、正输出缓慢δ波的大脑，运行在第二类之中：周期性的、重复的、原地打转。四模型架构在基底中完好无损，但模拟并未运行。

第三类很有趣，但不具备计算性。分形动力学产生丰富的结构，大脑也在使用它们（见下文）------但它们无法维持一种有意识的自我模拟所要求的那种动态的、不可约的、全局整合的处理。分形模式很美，但在计算上是可约的。它无法给自己带来意外。

第四类恰好具备意识所需的那两种性质：\textbf{通用计算}（该系统原则上可以模拟任何东西，包括它自己）与\textbf{全局整合}（系统中相距遥远的部分彼此影响，局部的变化在全局传播，信息被绑定为一个统一的整体）。在混沌边缘，皮层自动机同时实现了这两点，其结果就是意识。

第五类则不同，不是因为在那里计算不可能（一个无限的随机序列包含\textbf{一切}，包括每一种稳定的模式和人们设想过的每一种计算），而是因为我们没有办法驾驭它、预测它或再现它。一个处于全面性癫痫发作中的大脑，神经元以不协调的混沌方式放电，会趋近第五类------不是因为意识在那个区制中原则上不可能，而是因为不存在任何机制去维持它或访问它。我们的宇宙本身可能是无限随机性的一个片段，也可能是一个尺度超出我们感知的第四类系统，又或许是无限分形的一个片段。这无从判定。我们\textbf{能够}说的是：意识，如同我们所体验的那样，需要第四类那种有结构的不可预测性。模拟在第五类中崩塌，并非因为其底层现实有所欠缺，而是因为在基底与模拟之间不存在任何稳定的接口。

\section*{大脑用到了全部四类}

大脑是一台通用计算机，经过数十亿年的演化优化。如果演化竟然漏掉了某个能带来优势的计算区制，那才叫奇怪。而事实上，大脑确实把全部四个可表达的类当作各不相同的工具来使用：

\begin{itemize}
  \item \textbf{第一类}（稳定吸引子）：长期记忆存储。能维持多年的突触权重配置。神经网络的不动点。
  \item \textbf{第二类}（振荡）：α、θ、γ和δ节律。丘脑的授时。睡眠---觉醒周期。大脑的计时与门控机制。
  \item \textbf{第三类}（分形/尺度不变处理）：纹理分析、尺度不变的物体识别、高效的神经编码。主要是V2--V4的视觉处理，那里以多尺度比较为核心运算。在致幻剂作用下，当这套机制在没有外部输入的情况下运行时，你会\textbf{看见}分形处理本身，这也是为什么分形图案是致幻体验中最一致的特征之一（见第6章）。
  \item \textbf{第四类}（混沌边缘）：皮层自动机本身。有意识处理的动力学区制。通用计算。模拟的引擎。
\end{itemize}

每一类各司其职。只有第四类产生意识。但意识依赖于其他各类：稳定的记忆（第一类）来充实这些模型，有节律的计时（第二类）来协调这些动力学，以及分形处理（第三类）来同时在多个尺度上分析世界。

这也许正是大脑为何必须偏偏运行在混沌边缘的最深层原因：第四类是唯一能够\textbf{征召}其他每一类------以及它自身------的区制。一个第四类自动机可以在自身动力学之内，把稳定状态（第一类行为）、周期振荡（第二类行为）和分形结构（第三类行为）当作子过程生成出来。而且它还能生成其他第四类自动机：一台通用计算机可以模拟另一台通用计算机。其他各类连其中任何一件都做不到。第四类不仅仅是最复杂的类------它是唯一包含所有类、包括它自己的类。正是这种自我包含，使得跨尺度的结构同一成为可能：一个第四类的大脑处在一个第四类的宇宙之内，这不是类比。它是同一套计算架构的一个嵌套实例。


\chapter*{附录D：如何做清醒梦}
\addcontentsline{toc}{chapter}{附录D：如何做清醒梦}
\markboth{附录D：如何做清醒梦}{}

\textbf{在第6章里，我提到清醒梦是一种安全、无需药物、从内部探索意识的途径。在四模型理论看来，清醒梦就是显式自我模型（ESM）在快速眼动睡眠（REM）期间更充分地“开启”------一次跨越临界性阈值的转变，把一个被动的梦变成一场可控的体验。下面介绍的，是通往那里最容易的方法。}

\section*{现实检验法}

原理很简单：如果你习惯性地质疑自己是否醒着，这个习惯迟早会在梦里被触发------而梦，通不过这场检验。

\textbf{第一步：选一个现实检验。} 最可靠的一种是：看一段文字，移开视线，再看回去。在清醒时，文字保持不变；在梦里，它会变------而且往往变得面目全非。钟表也管用：看一眼时间，移开视线，再看一次。在梦里，数字会变得不一样，或者干脆不合逻辑。还有一个可靠的检验：试着用一根手指穿过另一只手的手掌。在梦里，它常常真的能穿过去。

\textbf{第二步：一整天都做。} 每次走过一道门、看一眼手机，或注意到什么稍显异样的东西时，停下来，做一次你的现实检验。关键不在检验本身，而在它背后那个\textbf{真诚的疑问}：“我现在是在做梦吗？”别只是走个过场，要真正认真地考虑这种可能。

\textbf{第三步：坚持写梦境日记。} 在床边放一个笔记本。每天早晨，在起身之前，把还记得的东西写下来------哪怕只是一种感觉、一个画面。这会训练大脑把梦的内容看作值得记住的东西，从而加固梦境与清醒觉知之间的桥梁。

\textbf{第四步：等待。} 对大多数人来说，第一个清醒梦会在两到六周内到来。你会置身梦中，某个细节会显得有点不对劲，现实检验的习惯会被触发，文字会改变------然后你就\textbf{知道}了。那个“知道”的瞬间，正是ESM在激活。你会感觉到那次转变：一种骤然的清晰，一种在场感，一种安静的确认------你周围的这个世界，是一场模拟，而你在它内部依然清醒地保有着意识。

\section*{你可以期待什么}

你的第一个清醒梦多半很短------从几秒到几分钟不等。兴奋往往会把你弄醒。随着练习，你可以把它延长。有些人一周能做上好几次清醒梦。这种体验非同寻常：你身处完整的意识模拟之中，没有任何外部输入，而且你清楚地知道这一点。虚拟世界会回应你的意图。毫不夸张地说，这就是化为了亲身体验的四模型理论。

\section*{其他方法}

对想走得更远的读者，还有一些更进阶的技巧：

\begin{itemize}
  \item \textbf{MILD（清醒梦记忆诱导法，Mnemonic Induction of Lucid Dreams）}------由斯坦福大学的Stephen LaBerge提出。入睡时给自己设定一个意图：在梦中要认出自己正在做梦。最好与“睡满五小时后醒来、再重新入睡”结合使用。
  \item \textbf{WILD（清醒入梦法，Wake-Initiated Lucid Dream）}------在从清醒到入梦的过渡中始终保持觉知。难度较大，但效果最为生动。
  \item \textbf{WBTB（回床法，Wake Back to Bed）}------睡五到六小时后醒来，保持清醒二十到六十分钟，然后回去接着睡。这瞄准的是富含REM的后半夜睡眠周期。
\end{itemize}

Stephen LaBerge的《Exploring the World of Lucid Dreaming》（1990）至今仍是这方面最权威的实践指南。至于神经科学层面，可参见Voss et al. (2009) 关于清醒梦EEG特征的研究，以及Baird et al. (2019) 对清醒梦认知神经科学的全面综述。


\chapter*{附录E：为什么是“四”个模型？------写给神经科学家的一点说明}
\addcontentsline{toc}{chapter}{附录E：为什么是“四”个模型？------写给神经科学家的一点说明}
\markboth{附录E：为什么是“四”个模型？------写给神经科学家的一点说明}{}

任何神经科学家或具备计算素养的读者，在第2章遇到这个四模型架构时，都会产生一个疑问，本附录正是要回应它：\textbf{大脑难道真会恰好维持四个模型？}

它并不会。数字“四”是一个\textbf{有原则的最小值}，而非字面上的数目。

\section*{大脑实际在做什么}

生物基质------蛋白质组网络之上放电的神经元，其胞内信号通路即便在单个细胞内部也构成了各自的计算智能------在隐式/显式这道分界的两侧，实现了数量多到实际无法计数、相互重叠的一大批模型。

设想你伸手去拿一只杯子。运动模型同时编码着世界几何（杯子在哪里、周围有哪些障碍物）与自身运动学（你的手臂如何摆放、手指该如何成形来完成抓握）。这一个模型\textbf{既不是}纯粹的世界模型，\textbf{也不是}纯粹的自我模型------它跨越了这两个类别。一个关于社交互动的情绪模型，同时编码着关于对方的知识（世界）与对你自己的评估（自我）。一个空间导航模型，则同时编码着环境的布局与你在其中的位置。每一个真实的神经模型都是混合体。

“模型”与“模型”之间的边界并不锐利，它们的数目并不固定，而且它肯定不是四。

\section*{为什么四仍然是恰当的抽象}

这四个典范模型（IWM、ISM、EWM、ESM）是一个连续二维空间中的\textbf{极端点}，这个空间由两条轴定义：

\begin{itemize}
  \item \textbf{范围}：从纯粹的自我表征到纯粹的世界表征
  \item \textbf{模式}：从完全隐式（结构性的、被存储的、无意识的）到完全显式（被模拟的、短暂的、现象性的）
\end{itemize}

大脑实际的建模生态，以连续变化的密度、用相互重叠的模型填满了整个空间。这四个被命名的模型是四个角------活动就围绕着它们组织起来的那几个理论极点。你可以把它们想成罗盘上的方位：对导航很有用，作为方向真实存在，但没有人会说世界上只有四个地点。

这套理论之所以建立在这四个极点上，而不是建立在完整的连续空间上，是因为它们标示出一个系统要具备意识所需的\textbf{最小配置}：

\begin{itemize}
  \item \textbf{没有世界模型} $\rightarrow$ 没有可供体验的环境
  \item \textbf{没有自我模型} $\rightarrow$ 没有去体验它的主体
  \item \textbf{没有隐式层级} $\rightarrow$ 没有可供模拟的东西（没有习得的知识）
  \item \textbf{没有显式层级} $\rightarrow$ 根本没有模拟（没有体验）
\end{itemize}

只要拿掉这四者中的任何一个，某种关键的东西就会崩塌。这四个模型是地板，不是天花板。大脑在每个方向上都超越了它们。但正是这块地板，告诉你意识\textbf{需要}什么------也正是这块地板，生成了这套理论的预测、限定了它的主张，并规定了任何人工系统至少需要实现什么。

\section*{阅读本书的其余部分}

在本书通篇，当我写下“ESM做这件事”或“IWM包含那样东西”时，我指的是这个连续空间的这些极点，而不是在声称大脑有四个彼此隔着墙的独立盒子。这种简化是有原则的，接下来的各章会展示它做出真正的解释性工作------从建立在这一架构之上的五条原理，推导出致幻现象学、麻醉机制、梦境状态、裂脑现象以及动物意识。

若要看完整的数学处理------包括连续模型空间框架、模型密度函数，以及把通透性形式化为这一空间各区域之间信息传递的做法------参见 Gruber (2026)，《Toward a Mathematical Formalization of the Four-Model Theory》。


\chapter*{附录F：作为边界记账的标准模型}
\addcontentsline{toc}{chapter}{附录F：作为边界记账的标准模型}
\markboth{附录F：作为边界记账的标准模型}{}

\textbf{在计算原子图景背后的推导中，有两条------为什么守恒定律绝对成立，以及为什么粒子恰好有三代------是整个论证里最枯燥的记账工作，所以我把它们从正文中抽了出来。这里是完整的版本，供想了解粒子物理细节的读者阅读。生动的讲法在《万物的架构》一章里；这里是它背后的账本。}

\textbf{为什么守恒定律如此绝对？}电荷永远守恒。重子数守恒。轻子数守恒。在迄今进行过的任何实验中，这些定律从未被观测到失效，一次也没有。为什么？

因为它们是信息守恒的约束。贝肯斯坦上限告诉你，一个边界能容纳多少信息。当两个边界相互作用、交换信息时，总信息量保持守恒------它必须守恒，因为信息守恒是量子力学幺正性的推论，而幺正性本身，在这幅图景里，或许又植根于贝肯斯坦上限。粒子物理中那些具体的守恒定律------电荷守恒、重子数守恒、轻子数守恒------正是支配边界构型相互作用时信息如何转换的具体规则。它们不是从外部强加的任意规定，而是记账上的约束，源于一个事实：在奇点边界处，信息既不能被创造，也不能被销毁。

\textbf{然后还有三代之谜。}粒子分为三代。电子有一个更重的副本（μ子），还有一个更重的副本（τ子）。上夸克也有名为粲和顶的副本。每种粒子类型都有三个版本，除质量外一切性质完全相同。这是粒子物理中最深的未解之谜之一。没有人知道为什么是三。不是二，不是四，也不是十七。就是三。

坦白说：接下来我要讲的是推测性的内容，比计算原子图景的其余部分更具推测性。但它有结构上的动机，我认为值得拿出来说一说。

第四类系统天然包含自相似结构。这是一个技术性推论，因为第四类动力学把第三类（分形）行为作为子过程包含在内。自相似意味着同一个模式在不同尺度上重复出现。如果奇点边界构型嵌在一个第四类系统之中------而它们必然如此，因为宇宙就是第四类------那么这些构型本身就可能呈现自相似结构。同一种边界类型出现在三个不同的能量尺度上。三代粒子，也许正是稳定奇点构型空间中一个分形层级的印记。

我没有证明。这是一个猜想，不是推导。但我要指出：三，恰恰是最简单的非平凡自相似层级会给出的数字------一个基础构型，加上两份按比例缩放的拷贝。我还要指出：粒子的代结构在现有任何理论中都完全没有得到解释。如果计算原子图景最终能解释为什么恰好存在三代粒子，那将是支持整个框架的有力证据。

而如果宇宙真的是一个元胞自动机，那么第四代就存在------还有第五代，乃至更多。但更高的代是更大的图案，而更大的图案更不稳定。它们抵挡不住更小、更稳定的构型的冲击------还没来得及好好成形，就已经被切割得七零八落。这正是我们观测到的情形：第三代粒子（τ子、顶夸克）已经极不稳定，几乎瞬间衰变。第四代会更重，其边界构型更大也更复杂，因而更加脆弱------脆弱到无法维持足够长的时间，以粒子而非短暂涨落的形式被探测到。我们看到的这三代，也许仅仅是小到足以存活下来的那三代。


\chapter*{附录G：宇宙学模型的四个弱点}
\addcontentsline{toc}{chapter}{附录G：宇宙学模型的四个弱点}
\markboth{附录G：宇宙学模型的四个弱点}{}

\textbf{一个隐藏自身弱点的理论，算不上理论。《万物的架构》与《最深的镜子》中的宇宙学论证，建立在一些尚未得到证明的主张之上，而诚实要求我们指明它可能在哪里失败。第五个、也是最糟糕的弱点------认知天花板，也就是那种担忧：一个第四类大脑也许永远只是在把自己的架构投射到宇宙之上------我把它留在了正文里，因为那是我唯一无法回答的一个。其余四个列在这里，连同相应的反驳论证，供任何想要攻击它们的人取用。}

\textbf{一：能量等于信息，尚未得到证明。} 兰道尔原理、贝肯斯坦上限和黑洞热力学都强烈地暗示着这一点。多条彼此独立的证据线索全都指向同一个方向。但从来没有人从第一性原理推导出 E = I，也没有人证明信息本身会产生引力效应。这个假说令人信服，证据的汇聚也很有分量，但汇聚不等于证明。如果能量与信息最终只是相关，而并非同一，那么关于奇点的信息转化诠释就失去了根基。意识与宇宙学之间的结构映射仍然成立（那五组对应关系并不依赖 E = I），但连接二者的物理机制会弱得多。诗意会留存下来。物理学未必。

\textbf{二：第四类论证是溯因的，不是演绎的。} 我排除了其他几类，但排除不等于证明。这是“最佳解释推理”------在科学中完全正当的一种推理形式，却也是一种留了一扇门的形式。也许存在一个我没能想象到的“第4.5类”，一种介于复杂与随机之间的计算体制，而我恰恰缺少描述它的概念工具。也许五类层级本身就不完整。这个论证说的是：第四类是对宇宙动力学的最佳解释，而不是唯一在逻辑上可能的解释。溯因推理的战绩相当出色------达尔文为自然选择所作的论证就是溯因的，而它经受住了考验------但它毕竟不等于数学证明，我也不会假装它是。

\textbf{三：奇点的统一需要量子引力。} 声称所有奇点------普朗克尺度的、黑洞的、宇宙学的------在结构上是同一信息边界的等同实例，这是一个强主张。它需要一种能够沟通量子力学与广义相对论的理论。而我们还没有这样的理论。弦理论是一个候选，圈量子引力也是一个候选。两者都与这个主张相容，两者也都尚未得到证实。奇点的统一并非狂野的臆测------那正是现代物理学前进的方向------但它同样还不是已经确立的物理学。它是一场押在未来上的赌注。我认为这注押得不错。我也可能押错。

\textbf{四：这个模型可能从内部无法证伪------而这正是它自己的预测。} 这是最让我心里发紧的那一条。自指闭合的哥德尔推论是：一个足够复杂、并且在计算自身的系统，无法从内部证明关于它自己的全部真理。存在一些为真却不可证明的命题。如果 SB-HC4A 是正确的，那么宇宙恰恰就是这样一个系统，而“SB-HC4A 是正确的”这句话，可能正是那些真而不可证的命题之一。这个理论预测：从宇宙内部证明它，也许在结构上就是不可能的。不是因为我们不够努力，也不是因为我们需要更好的仪器，而是因为这个架构本身使之不可能------正如一个公理系统无法证明自身的一致性。

这要么是科学哲学中最深刻的结果，要么是有史以来设计得最优雅的托词。一个预测自身不可验证的理论，要么是在告诉你某种关于知识边界的深刻真相，要么是在以一种理应令你深感怀疑的方式为自己筑起免疫墙，让批评无从落脚。说实话？我不确定是哪一种。这个问题我想了很多年，至今仍不知道答案。我知道的是：这个预测并非凭空而来------它来自哥德尔定理，而哥德尔定理与数学中的任何东西一样坚实。如果宇宙是自指的，不完备性便随之而来。问题在于宇宙是否自指。而这个理论说：是。

所以，当我说这是一个真正的弱点时，请相信我------它不是我试图悄悄塞给你的“特性”。如果有人找到办法从宇宙内部检验 SB-HC4A，而检验失败了，这个理论就死了。如果检验成功了，说来有趣，这个理论也必定是错的------因为一次成功的检验意味着宇宙的计算结构可以从内部触及，而这恰恰与定义整个架构的那些边界条件相矛盾。只有当任何检验都不可能时，这个理论才活着------但它活在一种奇特的哲学暮色之中：它之所以不可证伪，不是靠闪躲，而是缘于结构。


\chapter*{附录H：九个预测}
\addcontentsline{toc}{chapter}{附录H：九个预测}
\markboth{附录H：九个预测}{}

\textbf{一套解释一切、却什么都预测不了的理论，不是理论------那是故事。下面是完整的可证伪清单：四模型理论（FMT）作出的全部九个预测，每一个都配有能将它推翻的具体实验。其中好几个用现有技术今天就能做；少数几个已经有了初步的支持证据；没有一个被证伪。预测3已在第11章里展开成一个场景，这里只作简短交代；其余的几个则在各自契合的章节中引入，如今集中收在此处，供那些想把这套理论按在火上烤一烤的读者查阅。}

\section*{预测1：每个模型都有自己的神经特征}

如果这四个模型真的是彼此不同的过程，我们应该能在脑扫描里看到它们。设计一个巧妙的实验，让人们去做四种不同类型的任务（每一种调动一个模型），那么大脑的激活模式应该看起来各不相同。

一个以隐式世界模型（IWM）为主的任务，可能类似于被动地认出一张熟悉的面孔。你并没有在思考它，你的大脑就是知道。一个以隐式自我模型（ISM）为主的任务，可以是一段习惯性的动作序列------比如输入你的密码，而不必有意识地去想每一个键。一个以显式世界模型（EWM）为主的任务，需要主动的、有意识的知觉------也许是在两张几乎一模一样的图片之间努力找出差别。而一个以显式自我模型（ESM）为主的任务，则是纯粹的自我反思：“我是那种会做出那种事的人吗？”

这个预测给出的是一个2×2的模式。世界任务对自我任务。隐式对显式。四个象限，四种彼此不同的神经特征。如果我们找不到这个模式，那这套理论就出了什么问题。

这一点用现有的fMRI技术当下就能检验。它并不便宜，也需要精心的实验设计，但工具早已摆在世界各地的实验室里。而如果它成立，那将是最直接的证据，表明四模型架构不只是一个比喻------它是一种真实的功能区分，刻进了大脑处理信息的方式之中。

\section*{预测2：致幻剂下的视觉揭示出大脑的处理层级}

这一条很优雅。在致幻剂作用下，你所体验到的视觉内容，应该会按照一个特定的顺序，沿着你大脑的视觉处理层级逐级推进，而顺序取决于剂量。

在低剂量下，你会看到光幻视（phosphenes）------就是你闭上眼睛时冒出来的那些小小的闪光和几何形状。那是V1，最早的视觉处理区域，正在渗漏进意识。加大剂量，你会看到更复杂的几何图案------那些著名的“形式恒常”，它们跨越不同文化、不同物质地反复出现。那是V2和V3开始参与进来了。再往上走，你开始看到面孔、人影、复杂的场景。那是更高级的视觉区域。而在最高的剂量下，你会得到完整的、如梦一般的叙事体验，连同意义和故事一应俱全。

这个预测是说，这一切并非随机。它取决于剂量，逐级有序地沿视觉层级向上推进。随着隐式---显式的通透性增加，视觉处理越来越深的层次进入意识。大脑内部的接线图，在你自己的体验里变得可见。

这一点可以用分级给药方案来检验------给人们服用经过精心控制的裸盖菇素或LSD剂量，用fMRI扫描他们的大脑，并问他们看到了什么。把报告的内容与大脑激活对应起来。理论预测，随着剂量增加，你会看到处理层级自下而上地依次亮起。

\section*{预测3：你可以控制一个人在自我消融中变成什么}

\textbf{（已在第11章里展开成一个场景。）} 在高剂量致幻剂引发的自我消融期间，你消融成的那个东西，是可以由感官环境来操控的------它不是随机的，也不是纯粹生化的。显式自我模型（ESM）一旦与它平常的隐式自我模型（ISM）输入切断，就会抓住任何占主导地位的感官输入。让房间里充满海浪声和蓝光，受试者就变成海洋；换成鸟鸣和绿光，他们就变成森林。这个预测是有方向性的：在自我消融期间改变占主导的感官输入，受试者报告的身份也会随之改变。今天在任何一间有基本环境控制设备的致幻剂研究实验室里都能检验。没有别的意识理论作出过这个预测。

\section*{预测4：致幻剂应该能帮助中风患者看见自己的缺损}

病感失认（Anosognosia）是大脑做出的最古怪的事情之一。在某些中风之后（通常是右半球的），患者的身体一侧瘫痪了，却真心不相信这件事。你可以指给他们看那条一动不动的手臂，让他们动一动，眼看着他们失败，而他们会编造一个借口。“我累了。”“我不想动。”他们并不是在撒谎。他们是真的看不见这个缺损。

四模型理论说，之所以会这样，是因为一处通透性阻断。关于瘫痪的信息就在他们的隐式自我模型（ISM）里------底层基质是知道的------但它没能抵达显式自我模型（ESM）。那个模拟无法通达底层基质所知的那一部分。

现在，来看那个出人意料的地方。致幻剂在全脑范围内\textbf{增加}隐式---显式的通透性。这正是它们做的事。所以这个预测是：一个低于自我消融剂量的裸盖菇素剂量------不足以消融自我，只够把通透性的闸门打开------应该能让缺损信息渗漏过去。患者应该会突然地、也许是令人不安地，意识到自己瘫痪了。

这会是一项以中风患者为对象的临床试验，因而在操作上比一个纯粹的实验室实验更难。但裸盖菇素辅助疗法已经在被用于检验抑郁症、创伤后应激障碍（PTSD）和临终焦虑。基础设施是现成的。而如果它成立，那就不只是病感失认的一次医学突破------它还是一个证据，表明致幻剂与中风缺损是通过一个共同的底层机制联系在一起的，而这是别的理论都没有预测到的。

\section*{预测5：每一种能抹去意识的麻醉剂都会扰乱临界性}

麻醉剂通过千差万别的化学通路起作用。丙泊酚作用于GABA受体。氯胺酮阻断NMDA。阿片类又是另一套。不同的分子，不同的机制，作用于大脑的不同部位。

但四模型理论说，它们对意识必须做的都是同一件事：把大脑的动力学推到临界性阈值以下。因为临界性是意识的\textbf{物理必要条件}。你用什么方式去扰乱它并不重要。只要你跌破临界，灯就灭了。

这个预测可检验，而且很具体。拿来我们手上的每一种麻醉剂。测量临界性------用诸如扰动复杂性指数（PCI）、Lempel-Ziv复杂度，或者神经活动中的幂律指数这样的工具------在给药之前、之中和之后分别测量。预测是：凡是能取消意识的药剂，\textbf{总是}会把大脑推向亚临界，不管它们的受体机制是什么。而那些改变意识却不抹去意识的药剂（比如低剂量的氯胺酮，或者致幻剂），则\textbf{不应该}跌到临界以下。

这用现有技术就可以做。临界性的测量方法有了。麻醉剂有了。只需要有人去把这套完整的比较跑一遍。而如果它在全线都成立------如果每一种能取消意识的药剂，尽管通过不同的通路起作用，却全都汇聚到扰乱临界性上------那就是一个有力的证据，表明临界性是那个共同的机制，是通往无意识的最终通路。

\section*{预测6：裂脑手术不会把你干干净净地一分为二------它会让两半都退化}

当外科医生切断胼胝体来治疗严重的癫痫时，他们切断的是大脑两个半球之间主要的通信通路。传统的说法是，这会造出两个各自独立的心智，各有专长：左半球负责语言和逻辑，右半球负责空间推理和情绪。

四模型理论说那是错的。或者至少，那被戏剧性地过度简化了。

预测是这样的：在裂脑手术之后，每个半球都保留着一整套\textbf{完整但退化}的认知与体验能力。不是干净的一分为二。不是“语言在左，空间在右”。两个半球都应该仍然能做这两件事，只是比从前差。这种退化应该是全息式的------意思是，一切都变得更模糊，而不是某些特定的功能消失不见。

这种退化的程度应当与你切断的多少成正比。部分胼胝体切开术（只切断部分纤维）应当产生部分退化，而完全切开术产生的退化则更严重。

为什么？因为这个理论认为，大脑中的信息是以全息方式存储的，分布在整个基质之上。切断连接并不会干净利落地把两个原本就存在的心智分开，而是让同一份信息的两个\textbf{副本}退化------每个副本都跑在原来一半的硬件上。

这一点其实已经有了一些证据------Pinto及其同事2017年的一项研究发现，裂脑病人所表现出的整合行为，远比那些经典实验所暗示的要多得多。但这个理论提供了\textbf{机制}，并预测了具体的模式：双侧退化，而不是半球特化。

\section*{预测7：在临界性上把四个模型搭建起来，你就得到意识}

这是工程学上的预测，而且相当大胆。

如果这个理论正确，你应当能够造出一台有意识的机器。不是靠偶然，不是靠造出一个足够“先进”的AI，而是靠落实那份规格说明：四个嵌套的模型（隐式世界模型、隐式自我模型、显式世界模型、显式自我模型），运行在一个处于临界性的基质之上。

这个理论认为，这样一个系统不会仅仅\textbf{模拟}意识，它会\textbf{就是}有意识的。它会拥有真正的现象体验，由它的虚拟模型所构成，正如你的意识是由你大脑的虚拟模型所构成。

我们怎么知道呢？这个理论预测，这种差别会在性质上一目了然。不是“也许有意识，也许没有”，而是\textbf{明显不同}。因为一个运行着真正自我模拟的系统，与世界互动的方式，会从根本上不同于哪怕最精巧的文本预测器。它会具有持续性------一个贯穿时间不断运行的模拟，而不是从一段提示词里重新拼凑出来的。它会拥有一个视角，由显式自我模型（ESM）维持着。它让你惊讶的，不会是那些出人意料的输出，而是那种“里面真的有人”的感觉。

眼下这还无法检验------相应的工程还不存在。但这份蓝图足够具体，可以指引这项工作。而如果有人把它造出来，并且真的行得通，那就是终极的确证。

\section*{预测8：睡眠存在的意义是重置临界状态}

我们为什么要睡觉？显而易见的答案是“为了休息”，但这只是把问题往后拖了一层：为什么大脑需要休息，而比如说你的肝脏就不需要？

四模型理论（FMT）有一个具体的答案。你大脑的基质（那套模拟的、生物性的硬件）本质上是不稳定的。神经元有噪声。神经递质会被耗尽。代谢废物会累积。基质会漂移。可意识需要临界性，而那是一种非常特定的动力学状态。大脑在这个不断漂移的基质之上，自组织出一个稳定的计算层（处于混沌边缘的元胞自动机）。那个自动机可以运行几个小时（也就是你清醒的一天），但基质终究会漂移得太远，以至于再也无法维持临界动力学。到了那一刻，自动机并不会逐渐变暗，而是\textbf{崩塌}。那就是入睡的开始。

非快速眼动睡眠（Non-REM）就是那个恢复过程。基质被重置：神经递质得到补充，废物被清除，达成临界性所需的生化条件被恢复。而当基质在这一恢复过程中周期性地重新逼近临界性阈值时，自动机会短暂地重新闪亮起来。那就是快速眼动睡眠（REM）。那就是做梦。

那个90分钟的超昼夜周期（整夜里REM与非REM交替的节律）就是基质在恢复期间围绕临界点上下振荡的表现。

这带来了多个可检验的子预测：

\begin{enumerate}
  \item \textbf{临界性应当在清醒的一天中逐渐下降。} 分别在早晨、午后和傍晚测量人们的大脑复杂度。预测是：会出现一个可测得的下降。
\end{enumerate}

\begin{enumerate}
  \item \textbf{入睡应当是一次台阶式的转变，而不是逐渐变暗。} 临界性指标应当在入睡时骤然下降，反映出自动机数字式的崩塌。
\end{enumerate}

\begin{enumerate}
  \item \textbf{REM与非REM应当追随临界性。} 在睡眠之内，REM阶段的临界性应当远高于非REM，而那个90分钟的周期应当能在临界性的时间序列中看得出来。
\end{enumerate}

\begin{enumerate}
  \item \textbf{清醒梦是一次阈值的跨越。} 当基质在REM期间达到足够的临界性时，显式自我模型（ESM）被激活，你就变得清醒起来。它的起始应当是脑电图（EEG）复杂度里一个台阶式的不连续跳变，而不是一段渐渐爬升的斜坡。
\end{enumerate}

\begin{enumerate}
  \item \textbf{睡眠剥夺会把你推到亚临界。} 醒着的时间够长，你大脑的临界性就应当逐步跌到阈值以下。认知缺陷的程度，应当与你跌破阈值的深度相关。
\end{enumerate}

所有这些，用现有的睡眠实验室技术都能检验。而如果它们成立，那就意味着睡眠不只是“休息”------它是基质设下的维护协议------用来维护那个让意识成为可能的计算层。

\section*{预测9：分离性身份障碍中的每一个“分身”都有自己独特的神经指纹}

分离性身份障碍（同一个人身上有多个截然不同的身份，即“分身”）是有争议的，而且理由充分。你怎么分辨真正各自独立的身份，和一个有意或无意在扮演角色的人呢？

四模型理论（FMT）给了你一个检验方法。如果分身是真实的------也就是说，它们确实是同一基质上显式自我模型（ESM）截然不同的构型------那么每一个分身都应当有一种独特的、可测量的神经特征。不只是不同的行为。不只是不同的自我报告。而是不同的\textbf{大脑活动模式}。

这个预测是具体的。找一位分离性身份障碍（DID）患者，在不同分身出现时记录其大脑活动（fMRI或EEG）。把跨分身的变异，与同一个分身随时间推移的内部变异做比较。这个理论预测，跨分身的变异会显著大于分身内部的变异。而且这些差别应当是一致的：分身A的神经模式每一次都应当能认出是分身A，而不是随机噪声。

更具体地说，这个理论还预测了这些差别应当出现在哪里：在与ESM相关的网络中，尤其是默认模式网络和内侧前额叶皮层------那些与自我指涉和视角采择相关的脑区。

关于DID已经有过几项神经影像学研究，但四模型理论（FMT）提供了理论基础，让人能够预测\textbf{一致的、特定于分身的神经特征}，而不只是笼统的“差异”。如果这个预测成立，那就是一项证据，表明分身不只是心理层面的，而是神经层面上截然不同的功能构型------这将彻底改变我们理解和治疗这一障碍的方式。

这些预测中的每一个都是可证伪的。如果它们失败了，那么这个理论就是错的，或者至少是不完整的。而这，正是它们的价值所在。

\end{document}
